\documentclass[paper=b5,fontsize=13pt,bookmarkpackage=false,twoside]{scrartcl}

\usepackage[T1]{fontenc}
\usepackage[utf8]{inputenc}
\usepackage{newpxtext}
\usepackage{amsmath,amssymb,mathtools}
\usepackage{newpxmath}
\usepackage{bm}
\usepackage{array,booktabs}
\usepackage{enumitem}
\usepackage{microtype}
\usepackage{graphicx}
\usepackage{xcolor}
\definecolor{LinkBlue}{HTML}{174A8B}
\definecolor{CiteViolet}{HTML}{7A3E9D}
\definecolor{UrlTeal}{HTML}{00796B}
% Reserve the page builder's maximum depth below the nominal text box so that
% deep page-ending formulae still leave at least 0.91 cm of visible margin.
\newlength{\visiblebodydepthreserve}
\setlength{\visiblebodydepthreserve}{\maxdepth}
\usepackage[
  b5paper,
  left=0.7cm,
  right=0.7cm,
  top=0.91cm,
  bottom=\dimexpr 0.91cm+\visiblebodydepthreserve\relax,
  includehead,
  headheight=0.61cm,
  headsep=0.12cm
]{geometry}
\usepackage{scrlayer-scrpage}
\clearpairofpagestyles
\ohead*{\pagemark}
\setkomafont{pageheadfoot}{\normalfont}
\setkomafont{pagenumber}{\normalfont}
\pagestyle{scrheadings}
\usepackage[
  unicode,
  colorlinks=true,
  hypertexnames=false,
  linkcolor=LinkBlue,
  citecolor=CiteViolet,
  urlcolor=UrlTeal,
  bookmarksopen=true,
  bookmarksnumbered=false,
  pdfdisplaydoctitle=true,
  pdfstartpage=1,
  pdfstartview=Fit,
  pdfpagemode=UseOutlines
]{hyperref}

\hypersetup{
  pdftitle={Allgemeine Eigenwerttheorie Hermitescher Funktionaloperatoren / General Eigenvalue Theory of Hermitian Functional Operators},
  pdfauthor={J. von Neumann},
  pdfsubject={Faithful English translation},
  pdfkeywords={von Neumann, Hermitian operators, eigenvalue theory, spectral theorem}
}

\linespread{1.1}
\raggedbottom
\setlength{\parindent}{1.5em}
\setlength{\parskip}{0pt}
\setlength{\emergencystretch}{2em}
\allowdisplaybreaks[4]
\urlstyle{same}

\setkomafont{disposition}{\normalfont\bfseries}
\setkomafont{section}{\normalfont\bfseries\Large}
\setkomafont{subsection}{\normalfont\bfseries\large}
\setkomafont{subsubsection}{\normalfont\bfseries\normalsize}
\setkomafont{pageheadfoot}{\normalfont\small}
\setkomafont{pagenumber}{\normalfont\small}

\NewDocumentEnvironment{namedstatement}{m m}
  {\par\addvspace{0.55\baselineskip}\phantomsection
   \noindent\textbf{#1%
   \if\relax\detokenize{#2}\relax\else\ #2\fi.}\
   \itshape\ignorespaces}
  {\par\addvspace{0.25\baselineskip}}

\NewDocumentEnvironment{proof}{}
  {\par\addvspace{0.25\baselineskip}
   \noindent\textit{Proof.}\ \ignorespaces}
  {\par\addvspace{0.25\baselineskip}}

\newcounter{translatornote}
\NewDocumentCommand{\translatornote}{m}{%
  \stepcounter{translatornote}%
  \begingroup
  \renewcommand{\thefootnote}{T\arabic{translatornote}}%
  \footnote{\textit{Translator's note:} #1}%
  \addtocounter{footnote}{-1}%
  \endgroup
}

\begin{document}
\pagenumbering{arabic}

\begin{center}
  \vspace*{0.5\baselineskip}
  {\LARGE\bfseries
    Allgemeine Eigenwerttheorie Hermitescher\\
    Funktionaloperatoren\par}
  \vspace{0.65\baselineskip}
  {\Large\bfseries
    General Eigenvalue Theory of Hermitian\\
    Functional Operators\par}
  \vspace{0.7\baselineskip}
  {\large J. von Neumann\par}
  \vspace{0.35\baselineskip}
  {\normalsize
   \href{https://doi.org/10.1007/BF01782338}
   {Mathematische Annalen 102, 49--131}\par}
\end{center}
\vspace{0.6\baselineskip}

\section*{Introduction}\label{sec:introduction}
\addcontentsline{toc}{section}{Introduction}

\subsection*{I.}\label{intro:I}
\addcontentsline{toc}{subsection}{I.}

If we consider all functions on a certain metric space $\Omega$ (for example, the real line, the complex plane, the surface of the unit sphere, the interval $0,1$, and so forth) that satisfy certain regularity conditions (for example, that are continuous and, apart from finitely many corners, continuously differentiable; that are twice continuously differentiable; that have a finite integral of the square of the absolute value over $\Omega$\footnote{\label{fn:source-1}The functions may take complex values.}; and so forth), and that may in addition be subject to certain boundary conditions, then, as is well known, a \emph{linear operator} means the following: a correspondence $R$ that assigns to every function $f$ in our class a function $Rf$ (which need no longer belong to this class), but does so in such a way that always
\[
 R(a_1f_1+\cdots+a_kf_k)
 =a_1Rf_1+\cdots+a_kRf_k
 \qquad
 (a_1,\ldots,a_k\ \text{complex constants}).
\]

\newpage
If a general notion of measure (for example, in the sense of Lebesgue) exists in $\Omega$, if $dv$ is the volume element in $\Omega$ (on the line: $dx$; in the plane: $dx\,dy$; on the surface of the unit sphere: $\sin\vartheta\,d\vartheta\,d\varphi$; and so forth), and if integration over this entire space is denoted by $\int_\Omega$, then a linear operator $R$ is called \emph{self-adjoint} or \emph{Hermitian} if, for all admissible functions $f,g$,
\begingroup
\postdisplaypenalty=10000
\[
 \int_\Omega f\,\overline{Rg}\,dv
 =
 \int_\Omega Rf\,\overline g\,dv
\]
\endgroup
\noindent
holds.\footnote{\label{fn:source-2}$\overline f$ is the function whose value at every point is the complex conjugate of the value of $f$. Our terminology differs somewhat from the customary one. A differential operator $R$, for example, is usually called self-adjoint when
\[
 f\,\overline{Rg}-Rf\,\overline g
\]
is the complete differential of an expression in $f,g$ and their derivatives, so that
\[
 \int_\Omega f\,\overline{Rg}\,dv-\int_\Omega Rf\,\overline g\,dv
\]
is an expression depending only on the boundary values of $f,g$ and their derivatives. Our self-adjointness (Hermitian character) follows from this only when that expression vanishes as a consequence of suitable boundary conditions. Thus a differential operator that is self-adjoint in the usual sense may, for us, become self-adjoint only on a sufficiently restricted domain of definition (cf.\ \hyperref[intro:VII]{Introduction VII}).}

Examples of Hermitian linear operators (which we shall abbreviate as H.O.) are the following. We now also display the arguments of the functions, and let $P$ denote a general point of $\Omega$:
\begin{align*}
 Rf(P)&=f(P),\\
 Rf(P)&=x\,f(P),\\[-0.2em]
 &\qquad\text{($x$ any coordinate of $P$ in $\Omega$)},\\
 Rf(P)&=\int_\Omega \varphi(P',P)f(P')\,dv',\\[-0.2em]
 &\qquad\text{(}\varphi(P',P)=\overline{\varphi(P,P')},
 \text{ i.e.\ the kernel is Hermitian}\text{)},\\
 Rf(P)&=Lf(P),\\[-0.2em]
 &\qquad\text{($L$ any self-adjoint differential or partial differential}\\[-0.2em]
 &\qquad\text{\quad expression; cf.\ note \ref{fn:source-2})}.
\end{align*}

As is well known, the last two types of H.O. determine an eigenvalue problem which, when the kernel $\varphi(P,P')$, respectively the coefficients of $L$, possess sufficient regularity (namely, when only a ``point spectrum'' is present), can be formulated as follows.\footnote{\label{fn:source-3}The eigenvalue problem for regular integral operators was solved by Hilbert (\emph{Göttinger Nachrichten}, 1904, pp.~49--91); his method was later considerably simplified by Courant, Hellinger, F.~Riesz, E.~Schmidt, Toeplitz, and others. In solving the eigenvalue problem, differential operators (at least self-adjoint operators of second order) can be replaced by integral operators (Hilbert, \emph{Göttinger Nachrichten}, 1904, pp.~213--259) having the same eigenfunctions but reciprocal eigenvalues.}

An eigenvalue is a number $\lambda$ for which there exists a function $f\ne0$ such that
\[
 Rf=\lambda f;
\]
$f$ is then an eigenfunction. In general there are infinitely many eigenvalues $\lambda_1,\lambda_2,\ldots$, to which eigenfunctions $\varphi_1,\varphi_2,\ldots$ can be assigned so that they form a complete system of orthonormal functions. In the following display, the mathematically intended order of the two cases is used.\translatornote{In the source, the two values in this displayed case distinction are interchanged.} That is,
\[
 \int_\Omega \varphi_m\,\overline{\varphi_n}\,dv
 =
 \begin{cases}
 1,&m=n,\\
 0,&m\ne n
 \end{cases}
 \qquad\text{(orthonormality),}
\]
and, for every pair of admissible functions $f,g$,
\[
 \int_\Omega f\,\overline g\,dv
 =
 \sum_{n=1}^{\infty}
 \left(\int_\Omega f\,\overline{\varphi_n}\,dv\right)
 \overline{\left(\int_\Omega g\,\overline{\varphi_n}\,dv\right)}
 \qquad\text{(completeness).}
\]
As for every complete orthonormal system, the expansion theorem with convergence in the mean holds:
\[
 \int_\Omega
 \left|f-\sum_{n=1}^{N}a_n\varphi_n\right|^2\,dv\longrightarrow0,
 \qquad
 a_n=\int_\Omega f\,\overline{\varphi_n}\,dv,
\]
as $N\to\infty$.

If $\varphi(P,P')$, respectively the coefficients of $L$, cease to behave quite regularly, complications arise that are characterized by the occurrence of a continuous spectrum. Nevertheless, the concepts of eigenvalue and eigenfunction can be retained through suitable generalizations (in particular, by weakening the regularity requirements on the eigenfunction, or even by considering so-called differential eigenfunctions\footnote{\label{fn:source-4}Hellinger, Göttingen inaugural dissertation, 1907, p.~7; Carleman, \emph{Sur les équations intégrales à noyau réel et symétrique}, Uppsala, 1923, p.~82.}). The completeness theorems likewise remain valid under suitable generalization (replacement of the sum by an integral, and so forth); but the entire construction becomes substantially more complicated and loses much of its original intuitive character.\footnote{\label{fn:source-5}The eigenvalue problem for integral equations with a ``bounded'' kernel was solved by Weyl (Göttingen inaugural dissertation, 1908) with the aid of Hilbert's theory of bounded bilinear forms (\emph{Göttinger Nachrichten}, 1906, pp.~157--209). Weyl also settled the eigenvalue problem for an extensive class of self-adjoint differential operators of second order whose coefficients may be singular (\emph{Mathematische Annalen} \textbf{68} (1910), pp.~220--269). A still more general class of integral operators was investigated by Carleman; cf.\ note \ref{fn:source-4}.}

\subsection*{II.}\label{intro:II}
\addcontentsline{toc}{subsection}{II.}

If we replace the continuous spaces $\Omega$ considered in \hyperref[intro:I]{I} by the ``discrete space'' $1,2,3,\ldots$ of positive integers, replacing functions $f(P)$ ($P$ ranging over $\Omega$) by sequences $a_n$ ($n=1,2,\ldots$) and integrals $\int_\Omega f\,dv$ by sums $\sum_{n=1}^{\infty}a_n$ in the corresponding fashion, then we arrive at notions entirely analogous to those in \hyperref[intro:I]{I}.

The simplest examples of H.O. (indeed, as can be shown after the concept is made more precise, the only ones) are linear transformations in infinitely many variables with a Hermitian matrix:
\[
 R(x_1,x_2,\ldots)=(y_1,y_2,\ldots),
 \qquad
 y_n=\sum_{m=1}^{\infty}\alpha_{mn}x_m
 \quad(n=1,2,\ldots),
 \qquad
 \alpha_{mn}=\overline{\alpha_{nm}}.
\]
(Here the role played by the ``regularity conditions'' for functions in $\Omega$ is played by conditions such as the requirement that all $\sum_{m=1}^{\infty}\alpha_{mn}x_m$ converge, and the like. We shall consider these matrices in greater detail in \hyperref[app:III]{Appendix III}.) These linear transformations are plainly analogues of the integral-kernel transformations in continuous spaces $\Omega$:
\[
 Rf(P)=\int_\Omega\varphi(P',P)f(P')\,dv'.
\]

For these transformations, or for the corresponding Hermitian forms
\[
 \sum_{m,n=1}^{\infty}\alpha_{mn}x_m\overline{y_n},
\]
the eigenvalue problem was likewise posed and solved by Hilbert, provided the matrix $\{\alpha_{mn}\}$ (respectively the Hermitian form $\sum_{m,n=1}^{\infty}\alpha_{mn}x_m\overline{y_n}$) satisfies certain conditions corresponding to the ``regularity'' of the kernel $\varphi(P,P')$ in integral-kernel transformations.

The same phenomena occur. Under high regularity (complete continuity), there is only a point spectrum, with the same properties as in \hyperref[intro:I]{I}, except that $\int_\Omega$ must be replaced by $\sum_{n=1}^{\infty}$. Thus $\lambda$ is an eigenvalue if there is a sequence $(x_1,x_2,\ldots)$ such that
\[
 R(x_1,x_2,\ldots)=\lambda(x_1,x_2,\ldots),
 \qquad\text{i.e.}\qquad
 \sum_{m=1}^{\infty}\alpha_{mn}x_m=\lambda x_n
 \quad(n=1,2,\ldots),
\]
without $(x_1,x_2,\ldots)=(0,0,\ldots)$, and with finiteness of $\sum_{n=1}^{\infty}|x_n|^2$ also required as a ``regularity condition''; $(x_1,x_2,\ldots)$ is then an eigensequence. If $\lambda_1,\lambda_2,\ldots$ are all the eigenvalues, eigensequences
\[
 (x_{11},x_{12},\ldots),\quad
 (x_{21},x_{22},\ldots),\quad\ldots
\]
can be assigned to them so that they form a complete orthonormal system:
\[
 \sum_{m=1}^{\infty}x_{pm}\overline{x_{qm}}
 =
 \begin{cases}
 1,&p=q,\\
 0,&p\ne q
 \end{cases}
 \qquad\text{(orthonormality),}
\]
and\footnote{\label{fn:source-6}This is, as is well known, equivalent to the proper analogue of the completeness relation (cf.\ \hyperref[intro:I]{I}),
\[
 \sum_{n=1}^{\infty}a_n\overline{b_n}
 =
 \sum_{p=1}^{\infty}
 \left(\sum_{n=1}^{\infty}a_n\overline{x_{pn}}\right)
 \overline{\left(\sum_{n=1}^{\infty}b_n\overline{x_{pn}}\right)}.
\]}
\[
 \sum_{m=1}^{\infty}x_{mp}\overline{x_{mq}}
 =
 \begin{cases}
 1,&p=q,\\
 0,&p\ne q
 \end{cases}
 \qquad\text{(completeness).}
\]

Under lower regularity (boundedness), a continuous spectrum can again occur, causing the same complications as those sketched in \hyperref[intro:I]{I} for continuous spaces. Since we shall later have to concern ourselves with these in considerable detail, we shall not discuss them further here.\footnote{\label{fn:source-7}The completely continuous and bounded forms (respectively operators) were discovered by Hilbert, who solved their eigenvalue problem (\emph{Göttinger Nachrichten}, 1906, pp.~157--209). The bounded bilinear forms were investigated further by Hellinger; cf.\ note \ref{fn:source-4}, and \emph{Journal für die reine und angewandte Mathematik} \textbf{136} (1909), pp.~210--273.} Some sort of ``regularity'' assumption has, moreover, always been made in the literature to date, both for matrices and for integral operators; the eigenvalue problem has never been discussed on the basis of Hermitian character alone.

\subsection*{III.}\label{intro:III}
\addcontentsline{toc}{subsection}{III.}

The analogous behavior of the continuous spaces $\Omega$ among themselves, and also with the discrete space $1,2,\ldots$ considered here, is known to rest on the following fact.

If $\varphi_1,\varphi_2,\ldots$ is a complete system of orthonormal functions in the space $\Omega$ under consideration, then the correspondence
\[
 f\longleftrightarrow(x_1,x_2,\ldots),
 \qquad
 x_n=\int_\Omega f\,\overline{\varphi_n}\,dv
 \quad(n=1,2,\ldots)
\]
assigns to every $f$ with finite integral $\int_\Omega|f|^2\,dv$ a sequence $(x_1,x_2,\ldots)$ with finite sum $\sum_{n=1}^{\infty}|x_n|^2$. Distinct $f$ correspond to distinct $(x_1,x_2,\ldots)$; and, by the theorem of Fischer and F.~Riesz,\footnote{\label{fn:source-8}For example, \emph{Göttinger Nachrichten}, 1907, pp.~116--122.} every $(x_1,x_2,\ldots)$ with finite $\sum_{n=1}^{\infty}|x_n|^2$ does in fact correspond to an $f$ with finite $\int_\Omega|f|^2\,dv$.

The correspondence is linear; that is,
\[
\begin{gathered}
 f\longleftrightarrow(x_1,x_2,\ldots)
 \quad\Longrightarrow\quad
 af\longleftrightarrow(ax_1,ax_2,\ldots),\\
 \left.
\begin{aligned}
 f&\longleftrightarrow(x_1,x_2,\ldots),\\
 g&\longleftrightarrow(y_1,y_2,\ldots)
\end{aligned}
\right\}
 \quad\Longrightarrow\quad
 f+g\longleftrightarrow(x_1+y_1,x_2+y_2,\ldots),
\end{gathered}
\]
and,\footnote{\label{fn:source-9}We shall prove these generally known facts, together with the Fischer--F.~Riesz theorem, in \hyperref[chap:I]{Chapter I} and \hyperref[app:I]{Appendix I}.}
\[
 \int_\Omega f\,\overline g\,dv
 =\sum_{n=1}^{\infty}x_n\overline{y_n}.
\]
Thus all the notions used in describing eigenvalues and eigenfunctions are invariant under this correspondence; they must therefore exhibit the same behavior in the continuous spaces considered in \hyperref[intro:I]{I} and in the discrete space of \hyperref[intro:II]{II}.

It is essentially by these considerations, as is well known, that Hilbert reduced the eigenvalue theory of integral operators in continuous spaces to that of H.O. for sequences $(x_1,x_2,\ldots)$ (i.e.\ functions on the discrete space $1,2,\ldots$).\footnote{\label{fn:source-10}This is historically inaccurate insofar as the Fischer--F.~Riesz theorem was discovered only later. The same considerations have played an important role in the new quantum theory; cf.\ Schrödinger, \emph{Annalen der Physik} \textbf{79}, no.~8 (1926), pp.~734--756.}

\subsection*{IV.}\label{intro:IV}
\addcontentsline{toc}{subsection}{IV.}

The subject of this paper is the construction of a general eigenvalue theory for \emph{all} H.O. We shall proceed abstractly, arranging matters from the outset so that our results apply uniformly to all the continuous spaces named in \hyperref[intro:I]{I}, as well as to the discrete space in \hyperref[intro:II]{II} (that is, to the operators on the functions in these spaces).

To this end, in \hyperref[chap:I]{Chapter I} we shall give a general characterization (by five properties A through E; see there) that applies to each of the following function spaces:

\begin{quote}
All complex functions $f$, measurable in the Lebesgue sense, on a space $\Omega$ (the number line, the complex plane, the surface of the unit sphere, the interval $0,1$, and so forth; cf.\ \hyperref[app:I]{Appendix I}), with volume element $dv$, for which
\[
 \int_\Omega|f|^2\,dv
\]
is finite. All sequences $(x_1,x_2,\ldots)$ of complex numbers for which
\[
 \sum_{n=1}^{\infty}|x_n|^2
\]
is finite.
\end{quote}

In the first case, it should be observed that functions $f,g$ for which $f(P)=g(P)$ everywhere except on a Lebesgue null set are not to be regarded as distinct.

The elements of these function spaces admit all vector operations: they can be added and multiplied by complex numbers. Moreover, for two such elements the ``inner product''
\[
 \int_\Omega f\,\overline g\,dv,
 \qquad\text{respectively}\qquad
 \sum_{n=1}^{\infty}x_n\overline{y_n},
\]
is defined. We denote the inner product by $(f,g)$ (unless something special is said, we shall also regard sequences $(x_1,x_2,\ldots),(y_1,y_2,\ldots),\ldots$ as functions, as in \hyperref[intro:II]{II}, and denote them by $f,g,\ldots$). With its aid, the Hermitian character of the operators $R$ can be defined in every case by
\[
 (f,Rg)=(Rf,g).
\]
(In \hyperref[intro:I]{I} this was the definition, and it is equivalent to the one in \hyperref[intro:II]{II}.) As with vectors, the inner product permits the definition of an absolute value
\[
 |f|=\sqrt{(f,f)}
 \qquad
 \left(
 \sqrt{\int_\Omega|f|^2\,dv}
 \ \text{or}\
 \sqrt{\sum_{n=1}^{\infty}|x_n|^2}
 \right),
\]
and of a distance $|f-g|$ (thus the analogue of ordinary Euclidean distance). Our space is consequently ``metrized'' and ``topologized'' by it; that is, topological expressions such as closedness, continuity, and so forth become meaningful in the space.

It may not be superfluous to devote a word to the consistent restriction to functions $f$ with finite $\int_\Omega|f|^2\,dv$. In eigenvalue problems for sequences, it has been customary since the fundamental investigations of Hilbert (cf.\ note \ref{fn:source-7}) and E.~Schmidt\footnote{\label{fn:source-11}For example, \emph{Rendiconti del Circolo Matematico di Palermo} \textbf{25} (1908), pp.~57--73.} to require the finiteness of $\sum_{n=1}^{\infty}|x_n|^2$; for functions $f$ on continuous spaces, $\int_\Omega|f|^2\,dv$ is in any event finite for the eigenfunctions of the point spectrum. In the continuous spectrum, admittedly, neither $\sum_{n=1}^{\infty}|x_n|^2$ nor $\int_\Omega|f|^2\,dv$ is finite (if eigensequences or eigenfunctions still exist at all and recourse need not be had to ``differential solutions''). We shall nevertheless show, by generalizing Hilbert's methods, that it is especially advantageous to treat the eigenfunctions in the continuous spectrum as improper objects; that is, to formulate the eigenvalue problem without mentioning them explicitly.

When the discrete space $1,2,\ldots$ is taken as its basis, our function space is the so-called complex Hilbert space (i.e.\ the infinite-dimensional Euclidean space): the set of all sequences $(x_1,x_2,\ldots)$ with finite $\sum_{n=1}^{\infty}|x_n|^2$. Since, as was observed in \hyperref[intro:III]{III}, all our manifolds of functions are isomorphic (that is, can be mapped one-to-one onto one another in such a way that the operations $f+g$, $a f$, and $(f,g)$ pass into their analogues), we shall call all of them Hilbert spaces. They are all to be regarded as special realizations of the general abstract Hilbert space, which is characterized by the ``intrinsic'' properties A through E alone (and, up to isomorphism, completely; cf.\ \hyperref[chap:I]{Chapter I}).

We denote the abstract Hilbert space by $\mathfrak H$.

\subsection*{V.}\label{intro:V}
\addcontentsline{toc}{subsection}{V.}

In the abstract Hilbert space, by an H.O. we understand, as stated in \hyperref[intro:I]{I}, a function $R$ (defined for certain $f$ in $\mathfrak H$ and taking values in $\mathfrak H$) that possesses the following properties:
\begin{align}
 R(a_1f_1+\cdots+a_kf_k)
 &=a_1Rf_1+\cdots+a_kRf_k
 \quad
 (a_1,\ldots,a_k\ \text{complex constants}),
 \tag{$\alpha$}\label{eq:intro-V-alpha}\\
 (f,Rg)&=(Rf,g),
 \tag{$\beta$}\label{eq:intro-V-beta}
\end{align}
insofar as both sides have meaning. For the time being, moreover, we assume that the domain of $R$ is a linear manifold; that is, whenever it contains $f_1,\ldots,f_k$, it also contains $a_1f_1+\cdots+a_kf_k$. (The final definition of an H.O. will in fact be somewhat different, but this does not concern us at first.)

We call an H.O. $R$ \emph{bounded} if $|f|\le1$ implies $|Rf|\le C$ for some fixed constant $C$ independent of $f$. Since generally $|af|=|a|\,|f|$, this is equivalent to
\[
 |Rf|\le C|f|,
 \qquad\text{or also}\qquad
 |Rf-Rg|\le C|f-g|;
\]
thus it expresses a kind of continuity of $R$.\footnote{\label{fn:source-12}If we regard $\mathfrak H$ as a sequence space (Hilbert space proper), then the wording of our definition of boundedness differs from Hilbert's: in Hilbert's terminology we require the boundedness of $R$ folded with itself. Nevertheless, the two formulations are equivalent; we discuss this more closely in \hyperref[chap:II]{Chapter II}, \hyperref[thm:12]{Theorem 12}. It should also be noted that the entire topology of $\mathfrak H$ is always to be understood in the sense of so-called ``strong convergence''; cf.\ Weyl, dissertation, p.~9.}

In Hilbert's theory of bounded Hermitian forms (i.e.\ H.O.), one may always assume that $R$ is defined for every $f$ in $\mathfrak H$: if this is not the case, it can be extended to all of $\mathfrak H$ by virtue of its continuity. But if we wish to establish a general theory of H.O. that also goes beyond the bounded ones, we may not require $R$ to be defined everywhere in $\mathfrak H$. (Indeed, by a theorem of Toeplitz, every H.O. meaningful everywhere must be bounded; cf.\ \hyperref[thm:48]{Theorem 48}, $\alpha$), $\gamma$), and \hyperref[fn:61]{note 61}.) For example, let $\Omega$ be the interval $0,1$, and let $R$ be the operator
\[
 Rf=i\frac{df}{dx}
\]
for all differentiable functions vanishing at $0$ and $1$. Then $R$ is an H.O., but plainly it cannot be defined for nondifferentiable functions. Likewise, for a highly singular integral kernel $\varphi(P,P')$, the expression
\[
 \int_\Omega\varphi(P',P)f(P')\,dv'
\]
will not converge for every $f$ with finite $\int_\Omega|f|^2\,dv$. Another example, of a different but equally characteristic kind, is this: let $\Omega$ be the number line and $Rf(x)=x f(x)$. This $R$ too is an H.O.; but since
\[
 \int_{-\infty}^{\infty}|f(x)|^2\,dx
\]
may be finite while
\[
 \int_{-\infty}^{\infty}x^2|f(x)|^2\,dx
\]
is infinite, this $R$ also does not always have meaning.

The restriction that \eqref{eq:intro-V-alpha} and \eqref{eq:intro-V-beta} are to hold only insofar as both sides have meaning is therefore quite essential. As a minimum, however, we shall require that
\begin{equation}
 \text{the domain of $R$ be everywhere dense in $\mathfrak H$.}
 \tag{$\gamma$}\label{eq:intro-V-gamma}
\end{equation}
What this means can be seen without difficulty from, for example, $i\,d/dx$ and multiplication by $x$: the former is always meaningful for all polynomials vanishing at $0$ and $1$, and the latter for all functions that vanish outside an arbitrary but finite interval; both sets of functions are everywhere dense.

With this arrangement, however, we run the danger of requiring too little: our H.O. could already be meaningless at points at which they could still be defined naturally, merely because the domain was narrowed arbitrarily. (For example, for $i\,d/dx$ one could, despite the everywhere-denseness of the domain, have required analyticity rather than mere differentiability.) To avoid this difficulty, we introduce the concept of a maximal H.O.

An H.O. $S$ is called an \emph{extension} of the H.O. $R$ if, wherever $Rf$ has meaning, $Sf$ also has meaning and agrees with $Rf$. If $S\ne R$ (that is, if $S$ has meaning at more points than $R$), then $S$ is a \emph{proper extension}. An H.O. $R$ having no further proper extensions is called \emph{maximal}. (Abbreviations: ext., max.)

Since, as we shall see in \hyperref[thm:35]{Theorem 35}, every H.O. can be extended to a max.\ H.O., we must investigate max.\ H.O. first and foremost. To do this, however, it is necessary to give a general formulation of the eigenvalue problem, in the abstract Hilbert space, that takes account of point and continuous spectra on an equal footing. This will be done in the next paragraph, following Hilbert's formulation of the problem.

\subsection*{VI.}\label{intro:VI}
\addcontentsline{toc}{subsection}{VI.}

Consider first the $k$-dimensional Euclidean space. If an H.O. is given there, that is, a correspondence
\begin{align*}
 R(x_1,\ldots,x_k)&=(y_1,\ldots,y_k),\\
 y_n&=\sum_{m=1}^{k}\alpha_{mn}x_m
 \quad(n=1,\ldots,k),\\
 \alpha_{mn}&=\overline{\alpha_{nm}}.
\end{align*}
(there are of course no convergence questions here), then, as is well known, $R$ can always be put into diagonal form by introducing new Cartesian coordinates, that is, by a unitary transformation. Thus there is a unitary transformation matrix $\{u_{mn}\}$ such that
\[
 \sum_{p=1}^{k}u_{mp}\overline{u_{np}}
 =
 \begin{cases}
 1,&m=n,\\
 0,&m\ne n,
 \end{cases}
 \qquad
 \sum_{p=1}^{k}u_{pm}\overline{u_{pn}}
 =
 \begin{cases}
 1,&m=n,\\
 0,&m\ne n,
 \end{cases}
\]
and, after the transformation
\[
 \xi_n=\sum_{m=1}^{k}u_{mn}x_m,
 \qquad
 \eta_n=\sum_{m=1}^{k}u_{mn}y_m
 \quad(n=1,\ldots,k),
\]
$R$ is described simply by
\[
 \eta_n=\lambda_n\xi_n
 \quad(n=1,\ldots,k),
\]
where $\lambda_1,\ldots,\lambda_k$ are real constants, the eigenvalues of $R$.\footnote{\label{fn:source-13}As one sees, $\lambda_1,\ldots,\lambda_k$ are eigenvalues, and
\[
 (\overline{u_{11}},\ldots,\overline{u_{k1}}),\ldots,
 (\overline{u_{1k}},\ldots,\overline{u_{kk}})
\]
are corresponding eigenvectors; cf.\ \hyperref[intro:II]{II}.} That is,
\[
 \alpha_{mn}
 =\sum_{p=1}^{k}\lambda_pu_{mp}\overline{u_{np}}
 \qquad(m,n=1,\ldots,k).
\]

Here $\lambda_1,\ldots,\lambda_k$, apart from their order, are uniquely determined by $R$, but the $u_{mn}$ are not. As is immediately seen, they may be replaced by $\vartheta_nu_{mn}$, where $\vartheta_1,\ldots,\vartheta_k$ are arbitrary numbers of absolute value $1$; and when several of the $\lambda_n$ are equal, even the corresponding rows $(u_{1n},\ldots,u_{kn})$ may be transformed unitarily among themselves in an arbitrary manner.

An ambiguous formulation of the eigenvalue problem such as this one (transformability to diagonal form) is ill suited for passage from $k$-dimensional Euclidean space to infinite-dimensional, that is, Hilbert space. The correct formulation of the finite-dimensional problem, which is easy to generalize, is obtained by stating it uniquely after Hilbert. This is done as follows.

Set
\[
 e_{mn}(\lambda)
 =\sum_{\lambda_p\le\lambda}u_{mp}\overline{u_{np}}
 \qquad(\lambda\ \text{an arbitrary real number});
\]
then the matrices $\{e_{mn}(\lambda)\}$ are H.O. $E(\lambda)$ with the readily verified property
\begin{equation}
 E(\lambda)E(\mu)=E(\mu)E(\lambda)=E(\lambda)
 \qquad(\lambda\le\mu).
 \tag{$\alpha$}\label{eq:intro-VI-alpha}
\end{equation}
(In particular, $E(\lambda)^2=E(\lambda)$.)\footnote{\label{fn:source-14}The meaning of $AB$ for operators $A,B$ that are meaningful everywhere is clear.} Let $l_1,\ldots,l_h$ ($h\le k$) be the mutually distinct values among $\lambda_1,\ldots,\lambda_k$, arranged in increasing order. Then $E(\lambda)$, as a function of $\lambda$, is constant on the intervals
\[
 \lambda<l_1,\qquad
 l_1\le\lambda<l_2,\qquad\ldots,\qquad
 l_{h-1}\le\lambda<l_h,\qquad
 l_h\le\lambda,
\]
and has jumps from the left at $l_1,\ldots,l_h$. In the first of these intervals it is plainly $0$; in the last, by the unitarity of $\{u_{mn}\}$, it is the identity matrix, hence the identity operator $1$. The last formula for $\alpha_{mn}$ immediately gives, with $l_0$ arbitrary but $l_0<l_1$,\footnote{\label{fn:source-15}The final expression is a Stieltjes integral.}
\[
 \alpha_{mn}
 =\sum_{p=1}^{h}l_p\bigl(e_{mn}(l_p)-e_{mn}(l_{p-1})\bigr)
 =\int_{-\infty}^{\infty}\lambda\,de_{mn}(\lambda).
\]
It follows at once, if, as in \hyperref[intro:IV]{IV}, $(f,g)=\sum_{n=1}^{k}x_n\overline{y_n}$ when $f=(x_1,\ldots,x_k)$ and $g=(y_1,\ldots,y_k)$, that
\begin{equation}
 (f,Rg)
 =\int_{-\infty}^{\infty}\lambda\,d\bigl(f,E(\lambda)g\bigr).
 \tag{$\beta$}\label{eq:intro-VI-beta}
\end{equation}

It is easy to see that \eqref{eq:intro-VI-alpha} and \eqref{eq:intro-VI-beta}, together with $E(\lambda)=0$, respectively $1$, for small, respectively large, $\lambda$, and with right-continuity of $E(\lambda)$ as a function of $\lambda$, determine the operators $E(\lambda)$ completely. Since the diagonal form of $R$ can easily be recovered from the $E(\lambda)$ (respectively the $e_{mn}(\lambda)$), this is the unique formulation of the eigenvalue problem.

Let us now return to the abstract Hilbert space $\mathfrak H$. Conditions \eqref{eq:intro-VI-alpha} and \eqref{eq:intro-VI-beta} may be carried over verbatim, provided only that the $E(\lambda)$ are defined precisely: they are to be H.O. meaningful everywhere.\footnote{\label{fn:source-16}In the $k$-dimensional space $E(\lambda)^2=E(\lambda)$; hence, as is easily verified, $|E(\lambda)f|\le|f|$. One may therefore expect it to be bounded and thus take it to be meaningful everywhere; cf.\ also \hyperref[chap:III]{Chapter III}, \hyperref[thm:14]{Theorems 14}, \hyperref[thm:15]{15}, and \hyperref[thm:16]{16}.} We formulate the additional requirements concerning the dependence of $E(\lambda)$ on $\lambda$ as follows:
\begin{equation}
\begin{gathered}
 E(\lambda)\longrightarrow0\quad(\lambda\to-\infty),
 \qquad
 E(\lambda)\longrightarrow1\quad(\lambda\to+\infty),\\
 E(\lambda)\longrightarrow E(\lambda_0)
 \quad\text{as }\lambda\to\lambda_0\text{ with }\lambda\ge\lambda_0.
\end{gathered}
\tag{$\gamma$}\label{eq:intro-VI-gamma}
\end{equation}
\footnote{\label{fn:source-17}For operators $A_n,B$ meaningful everywhere, $A_n\to B$ means that $A_nf\to Bf$ for every $f$. The identity operator $1$ is defined by $1f=f$.}
(That precisely this is the correct generalization will be shown by its success.)

A family $E(\lambda)$ satisfying \eqref{eq:intro-VI-alpha} and \eqref{eq:intro-VI-gamma} is called a \emph{resolution of the identity} (abbreviated: r.i.); if $R$ is related to it by \eqref{eq:intro-VI-beta}, $R$ is the corresponding operator. Convergence questions do of course arise, but we shall see (cf.\ \hyperref[thm:36]{Theorem 36}) that they are best resolved by strengthening \eqref{eq:intro-VI-beta} as follows: $Rg$ is to have meaning if and only if
\begin{equation}
 \int_{-\infty}^{\infty}\lambda^2\,d\,|E(\lambda)g|^2
 \tag{$\beta'$}\label{eq:intro-VI-beta-prime}
\end{equation}
is finite; then, as we shall show, all the integrals in \eqref{eq:intro-VI-beta} are absolutely convergent, and \eqref{eq:intro-VI-beta} is to hold. (It follows from \eqref{eq:intro-VI-alpha} that $|E(\lambda)f|^2$, as a function of $\lambda$, never decreases; cf.\ \hyperref[thm:17]{Theorems 17} and \hyperref[thm:18]{18}. The integrand in \eqref{eq:intro-VI-beta-prime} is therefore nonnegative; hence the integral in \eqref{eq:intro-VI-beta-prime} is either absolutely convergent or diverges properly, namely to $+\infty$.) It is easy to see that an r.i. (that is, one satisfying \eqref{eq:intro-VI-alpha} and \eqref{eq:intro-VI-gamma}) uniquely determines the corresponding operator, if one exists: the domain is given by \eqref{eq:intro-VI-beta-prime}, and all $(f,Rg)$ by \eqref{eq:intro-VI-beta}, hence also $Rg$ itself.

The eigenvalue problem thus consists in deciding: Is there an associated H.O. for every r.i.? Is it maximal? Must different operators belong to different r.i.? Which H.O. belong to r.i.?\footnote{\label{fn:source-18}Hilbert proved that every bounded H.O. has one and only one r.i.; moreover, precisely those r.i. occur for which $E(\lambda)$ becomes constant, hence $0$ or $1$, for all sufficiently small or sufficiently large $\lambda$. Condition \eqref{eq:intro-VI-beta-prime} then drops out, since it is always satisfied. Citation: note \ref{fn:source-7}. In our formulation the distinction between point and continuous spectrum is evidently not yet completed (in \eqref{eq:intro-VI-gamma}, $\lambda\ge\lambda_0$ is required); this is immaterial for our purposes.} We shall see that the first three questions are to be answered affirmatively (cf.\ \hyperref[thm:36]{Theorem 36}). For the fourth we must therefore ask further: the H.O. associated with r.i. form a subclass of the class of all max.\ H.O.; what is this subclass?

\subsection*{VII.}\label{intro:VII}
\addcontentsline{toc}{subsection}{VII.}

We shall see (\hyperref[thm:36]{Theorems 36} and \hyperref[thm:38]{38}) that this subclass does not comprise all max.\ H.O. There is one exceptional operator $\overline R$,\footnote{\label{fn:source-19}\hyperref[app:IV]{Appendix IV} is devoted, among other things, to this operator.} and all other exceptions can be built from it and from the H.O. having an r.i. by certain trivial operations, not to be discussed here (cf.\ \hyperref[thm:38]{Theorem 38}). For $\overline R$ we shall see that its properties in fact exclude every reasonable formulation of an eigenvalue problem, but that it is by no means ``pathological'': it is easy to describe (cf.\ note \ref{fn:source-19}).

Nevertheless, max.\ H.O. that are bounded, definite, or real\footnote{\label{fn:source-20}Let $\varphi_1,\varphi_2,\ldots$ be a complete orthonormal system in $\mathfrak H$ (cf.\ \hyperref[def:3]{Definition 3}); then every $f$ can be written uniquely in the form $\sum_{n=1}^{\infty}x_n\varphi_n$. It is called real if all the $x_n$ are real. $\mathfrak Rf$ is obtained from $f$ by replacing every $x_n$ by $\Re x_n$ (the real part). $R$ is real if, whenever $Rf$ has meaning, $R\mathfrak Rf$ also has meaning and equals $\mathfrak R(Rf)$. (Thus an arbitrary choice of a complete orthonormal system enters into the concept of ``reality.'')} must possess an r.i.; thus certain simple properties are incompatible with the absence of an r.i. (cf.\ \hyperref[thm:40]{Theorem 40}). We may also interpret the occurrence of this $\overline R$ as showing that the Hermitian character of infinite-dimensional operators does not prevent the exceptional cases of elementary-divisor theory from occurring, in contrast to the finite-dimensional case or, in infinite dimensions, to real, bounded, or definite operators; but only a single exceptional type can occur.

For the moment, this may suffice concerning max.\ H.O. (In particular, we shall not discuss further how, and to what extent, knowledge of the r.i. -- that is, of the ``Hilbert spectral form'' -- solves the eigenvalue problem as it was formulated in \hyperref[intro:I]{I} and \hyperref[intro:II]{II}. In this connection we refer to the investigations cited in notes \ref{fn:source-4} and \ref{fn:source-5}.) It remains, however, to discuss those H.O. whose maximality is not assured, or is absent. H.O. are usually given in precisely this way: for example, as self-adjoint differential operators defined for all $n$-times differentiable functions ($n$ being their order) that satisfy certain boundary conditions. We mentioned in \hyperref[intro:V]{V} that every H.O. can be extended to a max.\ H.O.; it was not said, however, that this can be done in only one way.

We shall also obtain simple results on this point, including the following: if an H.O. is not itself max., or is not max.\ after a certain trivial extension, it can be made max.\ in a continuum of different ways (cf.\ the concluding remark of \hyperref[chap:VIII]{Chapter VIII}). Bounded H.O. always fall in the first case.

A simple example may illustrate how this nonunique extendibility arises. Let $\Omega$ be the interval $0,1$, and let $R$ be the operator $i\,d/dx$, defined for all differentiable functions. Since
\begin{align*}
 (f,Rg)-(Rf,g)
 &=-i\int_0^1\bigl\{f(x)\overline{g'(x)}+f'(x)\overline{g(x)}\bigr\}\,dx\\
 &=-i\bigl\{f(1)\overline{g(1)}-f(0)\overline{g(0)}\bigr\},
\end{align*}
$R$ is not Hermitian; but the operators $R_0$ and $R_\vartheta$, where $\vartheta$ is a number of absolute value $1$, are Hermitian. They arise from $R$ by the following restrictions of its domain:
\[
 f(0)=f(1)=0,
 \qquad\text{respectively}\qquad
 f(1)=\vartheta f(0).
\]
$R_\vartheta$ has maximal extensions, for example $S_\vartheta$. All the $S_\vartheta$ are also maximal extensions of $R_0$. They are all distinct: for if $S_{\vartheta_1}=S_{\vartheta_2}=S$, then $S$ would have meaning for an $f$ and a $g$ satisfying
\[
 f(0)=g(0)=1,\qquad f(1)=\vartheta_1,\qquad g(1)=\vartheta_2,
\]
and consequently
\[
 (f,Sg)-(Sf,g)
 =-i\bigl(\vartheta_1\overline{\vartheta_2}-1\bigr)\ne0,
\]
so that $S$ would not be an H.O.\translatornote{The source prints \(+i\) in the integration-by-parts formula and omits the factor \(i\) in this last display.  With the inner-product convention of this paper and \(R=i\,d/dx\), the boundary form has the factor \(-i\); both signs and the omitted factor have been corrected.} Thus the multiplicity of the maximal extensions here stems from the multiplicity of the possible boundary conditions. We shall see that, even in the most general case, the situation displays a certain analogy with this one (cf.\ \hyperref[thm:26]{Theorems 26} and \hyperref[thm:27]{27}).\footnote{\label{fn:source-21}Carleman had already found signs of this nonuniqueness for singular integral operators, but in that case he could not solve the eigenvalue problem completely (the analogue of the decisive condition \eqref{eq:intro-VI-alpha} from \hyperref[intro:VI]{VI} was lacking). Cf.\ note \ref{fn:source-4}.}

\subsection*{VIII.}\label{intro:VIII}
\addcontentsline{toc}{subsection}{VIII.}

The connection between the H.O. of $\mathfrak H$ and matrices is, as is well known, the following. Let $R$ be an H.O. and let $\varphi_1,\varphi_2,\ldots$ be a complete orthonormal system (cf.\ \hyperref[def:3]{Definition 3}); then the matrix $\{\alpha_{mn}\}$, where
\[
 \alpha_{mn}=(\varphi_m,R\varphi_n),
\]
is assigned to $R$. For bounded H.O. this assignment is entirely clear; for unbounded ones it involves certain complications. It will occupy us in \hyperref[app:III]{Appendix III}. The unitary transformation theory of unbounded matrices leads to a particularly peculiar pathology, but this will be investigated elsewhere.\footnote{\label{fn:source-22}Cf.\ a forthcoming paper by the author in \emph{Journal für die reine und angewandte Mathematik}, ``Zur Theorie der unbeschränkten Matrizen.''}

Nor shall we discuss here the relation of these results to the new quantum theory (the so-called ``transformation theory''). In any event, they show that the general theory of H.O. by no means everywhere exhibits the behavior generally presumed and anticipated in ``transformation theory'' on the basis of analogy with bounded, or even finite-dimensional, operators.

\subsection*{IX.}\label{intro:IX}
\addcontentsline{toc}{subsection}{IX.}

It remains to say something about the method and arrangement of this paper.

The essential device is to investigate a certain operator $U$, which could symbolically be denoted by
\[
 U=\frac{R+i1}{R-i1},
\]
in place of $R$. For if $R$ is an H.O., then $U$ preserves length (that is, $|Uf|=|f|$), and is therefore continuous and hence easier to handle (cf.\ \hyperref[thm:22]{Theorems 22}, \hyperref[thm:23]{23}, and \hyperref[thm:24]{24}). This ``Cayley transformation''
\[
 U=\frac{R+i1}{R-i1}
\]
is the core of our method.

All other methods fail: the elegant procedures of Hellinger and F.~Riesz, because they require iteration of the operator $R$, which is straightforward only for H.O. meaningful everywhere (and hence bounded), but is initially questionable for arbitrary H.O.; all maximum--minimum methods, as well as E.~Schmidt's method, because they are restricted from the outset to cases without continuous spectrum. Only Hilbert's original method, approximation by finite-dimensional section operators, permits certain results to be obtained, but only with very great difficulty and only a fraction of what we can achieve.\footnote{\label{fn:source-23}Using this method the author was able to settle the case of real H.O. (cf.\ note \ref{fn:source-20}), where the exceptional case indicated in \hyperref[intro:VII]{VII} had not yet appeared. The process of extending H.O., however, could not be surveyed as it will be, for example, in \hyperref[chap:VIII]{Chapter VIII}; and the method was so nonconstructive that, for instance, the well-ordering theorem had to be invoked (cf.\ the announcement in \emph{Jahresbericht der Deutschen Mathematiker-Vereinigung} \textbf{37}, nos.~1--4 (1928), pp.~11--15). Nor could nonreal H.O. be approached in this way. The author does not wish to miss this opportunity to thank E.~Schmidt for the interest he has shown in these investigations, especially in the extension of the results to nonreal H.O., and to note that conversations with him repeatedly furnished essential suggestions. In particular, the concept of ``hypermaximality'' of H.O. (cf.\ \hyperref[def:9]{Definition 9}), as well as the recognition that this property is necessary and sufficient for the existence of an r.i., is due to E.~Schmidt.}

The arrangement of the subject is as follows. \hyperref[chap:I]{Chapters I} through \hyperref[chap:IV]{IV} are preparatory: \hyperref[chap:I]{I}, structure of the space $\mathfrak H$; \hyperref[chap:II]{II}, general facts about operators in $\mathfrak H$; \hyperref[chap:III]{III}, linear manifolds and their projection operators; \hyperref[chap:IV]{IV}, reducibility of operators. In \hyperref[chap:V]{Chapters V} through \hyperref[chap:VIII]{VIII}, the principal work of solving the eigenvalue problem is carried out: \hyperref[chap:V]{V}, Cayley transformation; \hyperref[chap:VI]{VI}, extension elements; \hyperref[chap:VII]{VII}, deficiency indices; \hyperref[chap:VIII]{VIII}, the extension process. In \hyperref[chap:IX]{Chapters IX} and \hyperref[chap:X]{X}, the exact normal forms for max.\ H.O. are then established: \hyperref[chap:IX]{IX}, eigenvalue representation; \hyperref[chap:X]{X}, operators without an eigenvalue representation. Finally, \hyperref[chap:XI]{Chapters XI} and \hyperref[chap:XII]{XII} contain several general theorems about unbounded operators: \hyperref[chap:XI]{XI}, semiboundedness; \hyperref[chap:XII]{XII}, pathology of unboundedness.\footnote{\label{fn:source-24}The latter will be sharpened considerably in the author's paper announced in note \ref{fn:source-22}.} There are also several appendices. \hyperref[app:I]{Appendix I} proves that the function spaces of \hyperref[intro:I]{Introduction I} really do satisfy conditions A through E of \hyperref[chap:I]{Chapter I}, and hence that our theory is applicable to them.\footnote{\label{fn:source-25}Together with several considerations from \hyperref[chap:I]{Chapter I}, this of course yields a proof of the Fischer--F.~Riesz theorem.} In \hyperref[app:II]{Appendix II} we prove a theorem on the solution of the eigenvalue problem for unitary operators (more generally, for ``normal'' operators). This is a question about bounded operators that can be solved by F.~Riesz's method, but that has never yet been settled within the framework of Hilbert's theory.\footnote{\label{fn:source-26}It essentially amounts to the fact that two commuting bounded H.O. admit a common eigenvalue form (or spectral form) in Hilbert's sense. Cf.\ Hellinger--Toeplitz, encyclopedia article II, C.~13, p.~1562 ff., where the question is solved for completely continuous H.O.} We need this for \hyperref[chap:VIII]{Chapter VIII}; but since methodologically it belongs to Hilbert's theory, we prefer to treat it separately in this way. In \hyperref[app:III]{Appendix III} we examine more closely the connection between operators and matrices (cf.\ \hyperref[intro:VIII]{Introduction VIII} and note \ref{fn:source-23}); in \hyperref[app:IV]{Appendix IV} the operator $\overline R$ (cf.\ \hyperref[intro:VII]{Introduction VII}) is investigated more precisely.

Apart from this, the paper is a logically self-contained whole that can be read without any prior knowledge of the literature on ``infinitely many variables.'' (Accordingly, the contents of \hyperref[chap:I]{Chapters I} and \hyperref[chap:III]{III}, and of \hyperref[app:I]{Appendix I}, apart from their arrangement, are not new.)

\section*{I. The Abstract Hilbert Space}\label{chap:I}
\addcontentsline{toc}{section}{I. The Abstract Hilbert Space}

We characterize the abstract complex Hilbert space by five properties A through E. It proves convenient to append to some of them several simple theorems that already follow from the particular property alone.

\subsection*{A. $\mathfrak H$ is a linear space.}\label{axiom:A}

That is: in $\mathfrak H$ there are an addition $f+g$, a subtraction $f-g$, a scalar multiplication $af$ ($f,g$ in $\mathfrak H$, $a$ a complex number), and an element $0$; and these obey the rules of ordinary vector algebra.

\begin{namedstatement}{Definition}{1}\label{def:1}
A subset $\mathfrak M$ of $\mathfrak H$ is called a \emph{linear manifold} (abbreviated: lin.\ mfd.) if, together with $f_1,\ldots,f_n$, it also contains $a_1f_1+\cdots+a_nf_n$. If $\mathfrak M$ is any subset of $\mathfrak H$, then the set of all
\[
 a_1f_1+\cdots+a_nf_n
 \qquad
 (f_1,\ldots,f_n\text{ in }\mathfrak M;\ a_1,\ldots,a_n\text{ constants})
\]
is the smallest linear manifold containing $\mathfrak M$; it is called the linear manifold \emph{spanned by} $\mathfrak M$.
\end{namedstatement}

\begin{namedstatement}{Definition}{2}\label{def:2}
The $n$ elements $f_1,\ldots,f_n$ are called \emph{linearly independent} (abbreviated: lin.\ indep.) if
\[
 a_1f_1+\cdots+a_nf_n=0
 \quad\Longrightarrow\quad
 a_1=\cdots=a_n=0.
\]
The $n$ linear manifolds $\mathfrak M_1,\ldots,\mathfrak M_n$ are called linearly independent if
\[
 f_1+\cdots+f_n=0,
 \qquad
 f_1\in\mathfrak M_1,\ldots,f_n\in\mathfrak M_n,
\]
implies $f_1=\cdots=f_n=0$.
\end{namedstatement}

\subsection*{B. There is in $\mathfrak H$ an inner product, analogous to that of vector calculus, which induces a metric.}\label{axiom:B}

That is: there is a function $(f,g)$ ($f,g$ in $\mathfrak H$, and $(f,g)$ a complex number) having the following properties:
\[
\begin{array}{ll}
1.\ (af,g)=a(f,g),&
2.\ (f_1+f_2,g)=(f_1,g)+(f_2,g),\\[0.3em]
3.\ (f,g)=\overline{(g,f)},&
4.\ (f,f)\ge0,\quad\text{and it equals $0$ only when $f=0$.}
\end{array}
\]
From 1 and 2, by 3, it follows that
\[
 1'.\ (f,ag)=\overline a(f,g),
 \qquad
 2'.\ (f,g_1+g_2)=(f,g_1)+(f,g_2).
\]
The formula
\[
 |f|=\sqrt{(f,f)}
\]
defines an ``absolute value,'' and $|f-g|$ defines the metric (cf.\ \hyperref[thm:2]{Theorem 2} and note \ref{fn:source-27}).

\begin{namedstatement}{Theorem}{1}\label{thm:1}
One always has
\[
 |(f,g)|\le |f|\,|g|.
\]
\end{namedstatement}

\begin{proof}
First, by 2, $2'$, 3, and 4,
\begin{align*}
 \Re(f,g)
 &=
 \left(\frac{f+g}{2},\frac{f+g}{2}\right)
 -
 \left(\frac{f-g}{2},\frac{f-g}{2}\right)\\
 &\le
 \left(\frac{f+g}{2},\frac{f+g}{2}\right)
 +
 \left(\frac{f-g}{2},\frac{f-g}{2}\right)\\
 &=\frac12\bigl((f,f)+(g,g)\bigr)
 =\frac{|f|^2+|g|^2}{2}.
\end{align*}
We replace $f,g$ by $af,a^{-1}g$ ($a>0$): the left-hand side remains invariant, while the right-hand side has minimum $|f|\,|g|$. We then replace $f$ by $\vartheta f$ ($|\vartheta|=1$): the right-hand side remains invariant, while the left-hand side has maximum $|(f,g)|$. Hence
\[
 |(f,g)|\le |f|\,|g|.
\]
\end{proof}

\begin{namedstatement}{Theorem}{2}\label{thm:2}
One always has
\[
 |af|=|a|\,|f|,
 \qquad
 |f+g|\le|f|+|g|.
\]
\end{namedstatement}

\begin{proof}
The first assertion follows immediately from 1 and $1'$, the second from \hyperref[thm:1]{Theorem 1}:
\[
 (f+g,f+g)=(f,f)+(g,g)+2\Re(f,g),
\]
\[
 |f+g|^2
 =|f|^2+|g|^2+2\Re(f,g)
 \le |f|^2+|g|^2+2|f|\,|g|,
\]
and therefore $|f+g|\le|f|+|g|$.
\end{proof}

\begin{namedstatement}{Definition}{3}\label{def:3}
Two elements $f,g$ are \emph{orthogonal} (abbreviated: orth.) if $(f,g)=0$; two linear manifolds are orthogonal if every element of one is orthogonal to every element of the other. A set $\mathfrak M$ is called \emph{orthonormal} (abbreviated: orthonorm.) if, for $f,g$ in $\mathfrak M$,
\[
 (f,g)=
 \begin{cases}
 1,&f=g,\\
 0,&f\ne g.
 \end{cases}
\]
It is called \emph{complete} (abbreviated: complete orthonorm.) if it is not a proper subset of another orthonormal set. This plainly says that
\[
 (f,g)=0
 \quad
 (f\text{ fixed and }g\text{ arbitrary in }\mathfrak M)
\]
excludes $|f|=1$, hence also excludes $|f|>0$, and therefore implies $f=0$. (As is well known, for several elements or linear manifolds, pairwise orthogonality implies linear independence.)
\end{namedstatement}

\begin{namedstatement}{Definition}{4}\label{def:4}
By \hyperref[thm:2]{Theorem 2}, $|f-g|$ is a distance in $\mathfrak H$,\footnote{\label{fn:source-27}A ``distance'' $\overline{f,g}$ is, as is well known, required to satisfy the following axioms:
\[
\begin{array}{ll}
1.&\overline{f,g}\ge0,\quad\text{and it is $0$ only when $f=g$};\\
2.&\overline{f,g}=\overline{g,f};\\
3.&\overline{f,h}\le\overline{f,g}+\overline{g,h}.
\end{array}
\]
In linear spaces one adds
\[
 4.\quad
 \overline{f+h,g+h}=\overline{f,g},
 \qquad
 \overline{af,ag}=|a|\,\overline{f,g}.
\]
By \hyperref[thm:2]{Theorem 2}, all these properties are possessed by $\overline{f,g}=|f-g|$. Note that, by its bilinearity (1 and $1'$ in \hyperref[axiom:B]{B}) and \hyperref[thm:1]{Theorem 1}, $(f,g)$ is continuous.} and hence $\mathfrak H$ is topologized. A \emph{closed} (abbreviated: cl.) linear manifold is therefore one that contains its accumulation points. If to the linear manifold spanned by $\mathfrak M$ we adjoin all of its accumulation points, we obtain the smallest closed linear manifold containing $\mathfrak M$: the closed linear manifold spanned by $\mathfrak M$.
\end{namedstatement}

\subsection*{C. In the metric $|f-g|$, $\mathfrak H$ is separable.}\label{axiom:C}

That is: a certain countable set is everywhere dense in $\mathfrak H$.

\subsection*{D. $\mathfrak H$ possesses arbitrarily many (finitely many) linearly independent elements.\texorpdfstring{\protect\footnotemark}{}}\label{axiom:D}
\footnotetext{\label{fn:source-28}This excludes the possibility that $\mathfrak H$ be a finite-dimensional Euclidean space. Conditions \hyperref[axiom:A]{A} and \hyperref[axiom:B]{B} were, moreover, modeled on an axiomatic system for vector space given by H.~Weyl for finitely many dimensions (\emph{Raum, Zeit, Materie}, 5th ed., Berlin, 1923, \S\S~2--4).}

\subsection*{E. $\mathfrak H$ is complete.}\label{axiom:E}

That is: if a sequence $f_1,f_2,\ldots$ in $\mathfrak H$ satisfies the Cauchy convergence condition (for every $\varepsilon>0$ there is an $N=N(\varepsilon)$ such that $m,n\ge N$ implies $|f_m-f_n|\le\varepsilon$), then it is convergent (there exists an $f$ in $\mathfrak H$ such that for every $\varepsilon>0$ there is an $N=N(\varepsilon)$ for which $n\ge N$ implies $|f_n-f|\le\varepsilon$).

In \hyperref[app:I]{Appendix I} we shall show that all the function spaces named in \hyperref[intro:I]{Introduction I} and \hyperref[intro:II]{II} have properties A through E, and hence are special cases of the abstract Hilbert space $\mathfrak H$. Here, however, we proceed to prove that $\mathfrak H$ is in fact completely characterized by properties A through E: that $\mathfrak H$ is isomorphic to ordinary Hilbert space. For this it is necessary to prove several generally known theorems about orthonormal and complete orthonormal systems.

\begin{namedstatement}{Theorem}{3}\label{thm:3}
Every orthonormal set $\mathfrak M$ is at most countable; every complete orthonormal set $\mathfrak M$ is countably infinite.
\end{namedstatement}

\begin{namedstatement}{Remark}{}\label{rem:after-thm3}
Henceforth we therefore write all orthonormal and complete orthonormal sets as sequences $\varphi_1,\varphi_2,\ldots$ (which in the former case may terminate).
\end{namedstatement}

\begin{proof}
Let $\mathfrak M$ be orthonormal. For any two $f,g$ in $\mathfrak M$ with $f\ne g$,
\[
 |f-g|^2=|f|^2+|g|^2-2\Re(f,g)=2,
 \qquad
 |f-g|=\sqrt2.
\]
By the separability of $\mathfrak H$ (condition \hyperref[axiom:C]{C}), $\mathfrak M$ is therefore at most countable.

It remains to show that a finite orthonormal system $\mathfrak M$ is not complete. Let its elements be $\varphi_1,\ldots,\varphi_n$. Among the elements $a_1\varphi_1+\cdots+a_n\varphi_n$ there cannot be $n+1$ that are linearly independent; hence, by condition \hyperref[axiom:D]{D}, $\mathfrak H$ must contain an element $f$ different from all of them. Then
\[
 f^*=f-a_1\varphi_1-\cdots-a_n\varphi_n
\]
is certainly nonzero and is orthogonal to all $\varphi_1,\ldots,\varphi_n$ if we set
\[
 a_k=(f,\varphi_k)
 \qquad(k=1,\ldots,n).
\]
By the concluding observation in \hyperref[def:3]{Definition 3}, $\mathfrak M$ is therefore incomplete.
\end{proof}

\begin{namedstatement}{Theorem}{4}\label{thm:4}
Let $\varphi_1,\varphi_2,\ldots$ be an orthonormal system. Then the series
\[
 \sum_{n=1}^{\infty}(f,\varphi_n)\,
 \overline{(g,\varphi_n)}
\]
is absolutely convergent for all $f,g$ in $\mathfrak H$; in particular, when $f=g$,
\[
 \sum_{n=1}^{\infty}|(f,\varphi_n)|^2\le|f|^2.
\]
\end{namedstatement}

\begin{proof}
For $a_n=(f,\varphi_n)$, one has, as is well known,
\begin{align*}
 \left|f-\sum_{n=1}^{N}a_n\varphi_n\right|^2
 &=
 |f|^2
 -2\Re\sum_{n=1}^{N}(f,a_n\varphi_n)
 +\sum_{m,n=1}^{N}(a_m\varphi_m,a_n\varphi_n)\\
 &=
 |f|^2
 -2\Re\sum_{n=1}^{N}\overline{a_n}(f,\varphi_n)
 +\sum_{m,n=1}^{N}a_m\overline{a_n}(\varphi_m,\varphi_n)\\
 &=
 |f|^2-2\sum_{n=1}^{N}|a_n|^2+\sum_{n=1}^{N}|a_n|^2\\
 &=
 |f|^2-\sum_{n=1}^{N}|a_n|^2.
\end{align*}
Since the left-hand side is always nonnegative,
\[
 \sum_{n=1}^{N}|a_n|^2\le|f|^2,
 \qquad
 \sum_{n=1}^{N}|(f,\varphi_n)|^2\le|f|^2.
\]
Consequently $\sum_{n=1}^{\infty}|(f,\varphi_n)|^2$ converges, absolutely, and is at most $|f|^2$. Since
\[
 \left|(f,\varphi_n)\overline{(g,\varphi_n)}\right|
 \le
 \frac12|(f,\varphi_n)|^2
 +
 \frac12|(g,\varphi_n)|^2,
\]
the series $\sum_{n=1}^{\infty}(f,\varphi_n)\overline{(g,\varphi_n)}$ also converges absolutely. This proves everything.
\end{proof}

\begin{namedstatement}{Theorem}{5}\label{thm:5}
Let $\varphi_1,\varphi_2,\ldots$ be an orthonormal system. The series
\[
 \sum_{n=1}^{\infty}a_n\varphi_n
\]
converges if and only if
\[
 \sum_{n=1}^{\infty}|a_n|^2
\]
converges.
\end{namedstatement}

\begin{proof}
By condition \hyperref[axiom:E]{E}, it suffices to show that the Cauchy convergence condition has the same form for the two series. This follows, for $m\ge n$, from
\begin{align*}
 \left|
 \sum_{p=1}^{m}a_p\varphi_p
 -
 \sum_{p=1}^{n}a_p\varphi_p
 \right|^2
 &=
 \left|\sum_{p=n+1}^{m}a_p\varphi_p\right|^2\\
 &=
 \sum_{p,q=n+1}^{m}(a_p\varphi_p,a_q\varphi_q)\\
 &=
 \sum_{p,q=n+1}^{m}a_p\overline{a_q}(\varphi_p,\varphi_q)\\
 &=
 \sum_{p=n+1}^{m}|a_p|^2
 =
 \sum_{p=1}^{m}|a_p|^2-\sum_{p=1}^{n}|a_p|^2.
\end{align*}
\end{proof}

\begin{namedstatement}{Corollary}{}\label{cor:thm5}
If
\[
 f=\sum_{n=1}^{\infty}a_n\varphi_n
\]
(in the case of convergence), then
\[
 (f,\varphi_n)=a_n.
\]
\end{namedstatement}

\begin{proof}
For $N\ge n$,
\[
 \left(\sum_{p=1}^{N}a_p\varphi_p,\varphi_n\right)
 =
 \sum_{p=1}^{N}a_p(\varphi_p,\varphi_n)
 =a_n.
\]
By continuity of the inner product, we may let $N\to\infty$, obtaining $(f,\varphi_n)=a_n$.\translatornote{The source prints $a_m$ in this final clause; $a_n$ is clearly intended.}
\end{proof}

\begin{namedstatement}{Theorem}{6}\label{thm:6}
Let $\varphi_1,\varphi_2,\ldots$ be an orthonormal system. For every $f$, the series
\[
 f'=\sum_{n=1}^{\infty}a_n\varphi_n,
 \qquad
 a_n=(f,\varphi_n)
 \quad(n=1,2,\ldots),
\]
converges, and $f-f'$ is orthogonal to all $\varphi_1,\varphi_2,\ldots$.
\end{namedstatement}

\begin{proof}
The series converges by \hyperref[thm:4]{Theorems 4} and \hyperref[thm:5]{5}; and by the corollary to \hyperref[thm:5]{Theorem 5},
\[
 (f',\varphi_n)=a_n=(f,\varphi_n),
 \qquad
 (f-f',\varphi_n)=0.
\]
This proves everything.
\end{proof}

\begin{namedstatement}{Theorem}{7}\label{thm:7}
For an orthonormal system $\varphi_1,\varphi_2,\ldots$ to be complete, each of the following three conditions is necessary and sufficient:

\begin{enumerate}
\renewcommand{\labelenumi}{$\alpha)$}
\item\label{cond:thm7-alpha}
$\varphi_1,\varphi_2,\ldots$ span the closed linear manifold $\mathfrak H$.

\renewcommand{\labelenumi}{$\beta)$}
\item\label{cond:thm7-beta}
For every $f$,
\[
 f=\sum_{n=1}^{\infty}a_n\varphi_n,
 \qquad
 a_n=(f,\varphi_n)
 \quad(n=1,2,\ldots).
\]

\renewcommand{\labelenumi}{$\gamma)$}
\item\label{cond:thm7-gamma}
For every $f,g$,
\[
 (f,g)
 =
 \sum_{n=1}^{\infty}
 (f,\varphi_n)\overline{(g,\varphi_n)}.
\]
\end{enumerate}
\end{namedstatement}

\begin{proof}
First, completeness implies $\beta)$: for the $f-f'$ of \hyperref[thm:6]{Theorem 6} must vanish. Second, $\beta)$ implies $\alpha)$, because $f$ is an accumulation point of
\[
 \sum_{n=1}^{N}a_n\varphi_n,
\]
and hence of the linear manifold spanned by $\varphi_1,\varphi_2,\ldots$; it is therefore an element of the corresponding closed linear manifold. Third, $\beta)$ implies $\gamma)$, since
\[
 \sum_{n=1}^{N}a_n\varphi_n\longrightarrow f,
 \qquad
 \left|\sum_{n=1}^{N}a_n\varphi_n\right|^2\longrightarrow|f|^2
 \quad(N\to\infty),
\]
while
\[
 \left|\sum_{n=1}^{N}a_n\varphi_n\right|^2
 =
 \sum_{m,n=1}^{N}a_m\overline{a_n}(\varphi_m,\varphi_n)
 =
 \sum_{n=1}^{N}|a_n|^2.
\]
Thus
\[
 \sum_{n=1}^{\infty}|a_n|^2
 =
 \sum_{n=1}^{\infty}|(f,\varphi_n)|^2
 =
 |f|^2.
\]
If here we replace $f$ by $(f+g)/2$ and $(f-g)/2$ and subtract, we obtain the real part of $\gamma)$; if we replace $f,g$ by $if,g$, we obtain its imaginary part.

Fourth, $\alpha)$ implies completeness. Otherwise there would be an $f\ne0$ orthogonal to all $\varphi_1,\varphi_2,\ldots$. The entire closed linear manifold spanned by $\varphi_1,\varphi_2,\ldots$, that is, $\mathfrak H$, would then be orthogonal to $f$, and hence so would $f$ itself:
\[
 |f|^2=(f,f)=0,\qquad f=0,
\]
contrary to the assumption. Fifth, $\gamma)$ also implies completeness; for the same $f$, condition $\gamma)$, with $f=g$, would give
\[
 |f|^2=(f,f)=\sum_{n=1}^{\infty}0^2=0,
 \qquad f=0.
\]

Thus we have the logical scheme
\[
\begin{gathered}
\text{completeness}\Longrightarrow\beta),\\
\beta)\Longrightarrow\alpha)\Longrightarrow\text{completeness},
\qquad
\beta)\Longrightarrow\gamma)\Longrightarrow\text{completeness}.
\end{gathered}
\]
Consequently these four assertions are indeed equivalent.
\end{proof}

\begin{namedstatement}{Theorem}{8}\label{thm:8}
For every sequence $f_1,f_2,\ldots$ of arbitrary elements of $\mathfrak H$, there is an orthonormal system $\varphi_1,\varphi_2,\ldots$ spanning the same linear manifold, and therefore also the same closed linear manifold. Either of the two sequences may already terminate after finitely many terms.
\end{namedstatement}

\begin{namedstatement}{Remark}{}\label{rem:after-thm8}
Choose $f_1,f_2,\ldots$ everywhere dense, as condition \hyperref[axiom:C]{C} permits. It then plainly spans the closed linear manifold $\mathfrak H$. The same is true of the orthonormal system $\varphi_1,\varphi_2,\ldots$ constructed by \hyperref[thm:8]{Theorem 8}; hence, by condition $\alpha)$ of \hyperref[thm:7]{Theorem 7}, it is complete. This proves the existence of complete orthonormal systems.
\end{namedstatement}

\begin{proof}
First, we may replace $f_1,f_2,\ldots$ by a subsequence consisting entirely of linearly independent elements and spanning the same linear manifold. Indeed, let $g_1$ be the first $f_n\ne0$, let $g_2$ be the first $f_n$ not of the form $a_1g_1$ with $a_1$ arbitrary, let $g_3$ be the first $f_n$ not of the form $a_1g_1+a_2g_2$ with $a_1,a_2$ arbitrary, and so on. (If there is no $f_n$ outside the set of all $a_1g_1+\cdots+a_pg_p$, we stop with $g_p$.) The sequence $g_1,g_2,\ldots$ plainly has the property just required.

We now apply E.~Schmidt's ``orthogonalization procedure'' to $g_1,g_2,\ldots$:
\[
\begin{array}{lll}
\gamma_1=g_1,
&&
\displaystyle\varphi_1=\frac{1}{|\gamma_1|}\gamma_1,
\\[0.9em]
\gamma_2=g_2-(g_2,\varphi_1)\varphi_1,
&&
\displaystyle\varphi_2=\frac{1}{|\gamma_2|}\gamma_2,
\\[0.9em]
\gamma_3=g_3-(g_3,\varphi_1)\varphi_1-(g_3,\varphi_2)\varphi_2,
&&
\displaystyle\varphi_3=\frac{1}{|\gamma_3|}\gamma_3,
\\[-0.1em]
\multicolumn{3}{c}{\dotfill}
\end{array}
\]
The sequence $\varphi_1,\varphi_2,\ldots$ plainly spans the same linear manifold as $g_1,g_2,\ldots$, and hence as $f_1,f_2,\ldots$; it is also immediately seen to form an orthonormal system.
\end{proof}

On the basis of these facts, we can now show that $\mathfrak H$ can be mapped one-to-one onto ordinary Hilbert space, that is, the set of all sequences of complex numbers $(x_1,x_2,\ldots)$ with finite $\sum_{n=1}^{\infty}|x_n|^2$, in such a way that, with $\longleftrightarrow$ denoting the correspondence,
\[
 f\longleftrightarrow(x_1,x_2,\ldots)
 \quad\Longrightarrow\quad
 af\longleftrightarrow(ax_1,ax_2,\ldots),
\]
\[
 \left.
\begin{aligned}
 f&\longleftrightarrow(x_1,x_2,\ldots),\\
 g&\longleftrightarrow(y_1,y_2,\ldots)
\end{aligned}
\right\}
\quad\Longrightarrow\quad
 f+g\longleftrightarrow(x_1+y_1,x_2+y_2,\ldots),
\]
and, moreover,
\[
 (f,g)=\sum_{n=1}^{\infty}x_n\overline{y_n}.
\]

Indeed, let $\varphi_1,\varphi_2,\ldots$ be an arbitrary but fixed complete orthonormal system, and assign to $f$ the sequence $(x_1,x_2,\ldots)$ with
\[
 x_n=(f,\varphi_n)
 \qquad(n=1,2,\ldots).
\]
That $(x_1,x_2,\ldots)$ belongs to Hilbert space follows from \hyperref[thm:4]{Theorem 4}; that an $f$ exists for every such $(x_1,x_2,\ldots)$ follows from \hyperref[thm:5]{Theorem 5} and its corollary. One-to-one correspondence follows from completeness: if $(f,\varphi_n)=(g,\varphi_n)$ for every $n$, then $f-g$ is orthogonal to every $\varphi_n$, and hence is $0$. Of the three concluding assertions, the first two are clear, while the third follows from condition $\gamma)$ of \hyperref[thm:7]{Theorem 7}.

Thus it has been shown that $\mathfrak H$ is isomorphic to the ordinary Hilbert space, and hence that any two systems $\mathfrak H$ satisfying A through E are isomorphic to one another; that is, A through E form a ``categorical'' system of axioms: they determine all the properties of $\mathfrak H$.

\section*{II. Operators in Hilbert Space}\label{chap:II}
\addcontentsline{toc}{section}{II. Operators in Hilbert Space}

An operator is a function defined on a subset of $\mathfrak H$ (possibly on all of $\mathfrak H$) and taking values in $\mathfrak H$. We shall first define several especially important kinds of operators.

\begin{namedstatement}{Definition}{5}\label{def:5}
An operator $R$ is \emph{linear} (abbreviated: lin.) if, whenever $Rf_1,\ldots,Rf_n$ are defined, $R(a_1f_1+\cdots+a_nf_n)$ is also defined, and in fact
\[
R(a_1f_1+\cdots+a_nf_n)=a_1Rf_1+\cdots+a_nRf_n.
\]
An operator $R$ is \emph{closed} (abbreviated: cl.) if, for every sequence $f_1,f_2,\ldots$ for which all $Rf_n$ are defined, the convergences
\[
f_n\longrightarrow f,\qquad Rf_n\longrightarrow f^*
\]
imply that $Rf$ is defined and equals $f^*$.\footnote{\label{fn:29}This is evidently a continuity-like property; in every ``compact'' topological space it is in fact equivalent to continuity. In Hilbert space, however, it means much less: for Hermitian operators it is essentially always present; cf. \hyperref[thm:9]{Theorem 9}.}
\end{namedstatement}

\begin{namedstatement}{Definition}{6}\label{def:6}
An operator $R$ is \emph{Hermitian} (abbreviated: a H.O.) if, first, its domain spans $\mathfrak H$ as a closed linear manifold,\footnote{\label{fn:30}To this extent we weaken condition $\gamma)$ of \hyperref[intro:V]{Introduction V} (with a view to the application to matrices; cf. \hyperref[app:III]{Appendix III}): there, density of the domain was required.} and, second, whenever $Rf$ and $Rg$ are defined,
\[
(f,Rg)=(Rf,g).
\]
An operator $R$ is \emph{length-preserving} if it is closed and linear and, whenever $Rf$ is defined,
\[
|f|=|Rf|.
\]
\end{namedstatement}

We begin by proving two theorems about Hermitian and length-preserving operators. We also recall that an operator $S$ is called an \emph{extension} (abbreviated: ext.) of an operator $R$ if, wherever $Rf$ is defined, $Sf$ is also defined and equals $Rf$; it is called a \emph{proper extension} (abbreviated: prop. ext.) if $R\ne S$, that is, if $S$ is defined at more points than $R$.

\begin{namedstatement}{Theorem}{9}\label{thm:9}
Let $R$ be a H.O. There exists a linear H.O. $\widehat R$ that is an extension of $R$ and of which every H.O. with these same properties is in turn an extension; and there exists a closed linear H.O. $\widetilde R$ that is an extension of $R$ (and hence also of $\widehat R$) and of which every H.O. with these same properties is in turn an extension.
\end{namedstatement}

\begin{proof}
If we could define $\widehat R$ as follows: $\widehat Rf$ is defined whenever
\[
f=a_1f_1+\cdots+a_nf_n
\]
and $Rf_1,\ldots,Rf_n$ are defined, and its value is then
\[
a_1Rf_1+\cdots+a_nRf_n;
\]
and define $\widetilde R$ as follows: $\widetilde Rf$ is defined whenever there exists a sequence $g_1,g_2,\ldots$ for which all $\widehat Rg_n$ are defined and
\[
g_n\longrightarrow f,\qquad \widehat Rg_n\longrightarrow f^*,
\]
and its value is then $f^*$---then these $\widehat R$ and $\widetilde R$ would plainly have all the desired properties. The question is whether these prescriptions are consistent.

They certainly are if
\[
a_1f_1+\cdots+a_mf_m=b_1g_1+\cdots+b_ng_n,
\]
where $Rf_1,\ldots,Rf_m,Rg_1,\ldots,Rg_n$ are defined, implies
\[
a_1Rf_1+\cdots+a_mRf_m=b_1Rg_1+\cdots+b_nRg_n,
\]
and if
\[
f_n\longrightarrow f,\qquad g_n\longrightarrow f,
\]
where all $\widehat Rf_n,\widehat Rg_n$ are defined, together with
\[
\widehat Rf_n\longrightarrow f^*,\qquad \widehat Rg_n\longrightarrow g^*,
\]
implies $f^*=g^*$. It plainly suffices to prove the following two assertions:
\[
a_1f_1+\cdots+a_nf_n=0
\quad(Rf_1,\ldots,Rf_n\text{ defined})
\]
implies
\[
a_1Rf_1+\cdots+a_nRf_n=0,
\]
and
\[
f_n\longrightarrow0,\qquad \widehat Rf_n\longrightarrow f^*
\quad(\widehat Rf_n\text{ all defined})
\]
implies $f^*=0$.

Suppose first that the premise of the first assertion holds. Whenever $Rg$ is defined,
\begin{align*}
\left(g,a_1Rf_1+\cdots+a_nRf_n\right)
 &=\sum_{p=1}^n(g,a_pRf_p)\\
 &=\sum_{p=1}^n(Rg,a_pf_p)\\
 &=\left(Rg,a_1f_1+\cdots+a_nf_n\right)=0.
\end{align*}
Thus all such $g$ are orthogonal to $a_1Rf_1+\cdots+a_nRf_n$, and so is their closed linear span, $\mathfrak H$. Hence $a_1Rf_1+\cdots+a_nRf_n$ is orthogonal to itself and therefore is $0$.

Now consider the premise of the second assertion. Whenever $\widehat Rg$ is defined, since $\widehat R$ is a H.O.,
\begin{align*}
(g,f^*)
 &=\left(g,\lim_{n\to\infty}\widehat Rf_n\right)
  =\lim_{n\to\infty}(g,\widehat Rf_n)\\
 &=\lim_{n\to\infty}(\widehat Rg,f_n)
  =\left(\widehat Rg,\lim_{n\to\infty}f_n\right)=0.
\end{align*}
Thus all such $g$ are orthogonal to $f^*$, and so is their closed linear span, $\mathfrak H$; this again implies $f^*=0$.
\end{proof}

\begin{namedstatement}{Theorem}{10}\label{thm:10}
Let $U$ be a length-preserving operator. Then its domain and its range are both closed linear manifolds, and $U$ maps them continuously and one-to-one onto one another. Wherever $U$ is defined,
\[
(f,g)=(Uf,Ug).
\]
\end{namedstatement}

\begin{namedstatement}{Remark}{}\label{rem:10}
If both the domain and the range equal $\mathfrak H$, then, in the customary terminology, $U$ is called \emph{unitary}.
\end{namedstatement}

\begin{proof}
Since $U$ is linear, its domain and range are linear manifolds. Moreover,
\[
|Uf-Ug|=|U(f-g)|=|f-g|.
\]
Thus $Uf=Ug$ implies $f=g$, so the mapping is one-to-one; the continuity of both the mapping and its inverse follows as well. Together with the closedness of $U$, this implies that its domain and range are both closed sets.

In $|f|^2=|Uf|^2$, replace $f$ successively by $(f+g)/2$ and $(f-g)/2$ and subtract. This gives
\[
\operatorname{Re}(f,g)=\operatorname{Re}(Uf,Ug).
\]
Replacing $f,g$ by $if,g$ gives the same result for the imaginary parts. Therefore
\[
(f,g)=(Uf,Ug).
\]
This proves everything.
\end{proof}

Several properties of $\widetilde R$ (for $R$ a H.O.) are evident:
\[
\widetilde{\widetilde R}=\widetilde R;
\]
if $S$ is an extension of $R$, then $\widetilde S$ is an extension of $\widetilde R$. We now define maximality for Hermitian operators.

\begin{namedstatement}{Definition}{7}\label{def:7}
A H.O. $R$ is called \emph{maximal} (abbreviated: max.) if it has no proper extension that is also a H.O.
\end{namedstatement}

Since $\widetilde R$ is an extension of $R$, a maximal $R$ must satisfy $R=\widetilde R$. Furthermore, let $R$ be a H.O., let $S$ be a maximal H.O., and let $S$ be an extension of $R$; then $S=\widetilde S$ is also an extension of $\widetilde R$. Thus every maximal H.O. is closed and linear, and a H.O. $R$ has exactly the same maximal extensions as the closed linear H.O. $\widetilde R$. In later chapters we shall investigate the extension of arbitrary Hermitian operators to maximal ones (cf. also \hyperref[intro:VI]{Introduction VI} and \hyperref[intro:VII]{VII}); what has just been said shows that throughout this investigation we may restrict ourselves to closed linear operators.

We define further:

\begin{namedstatement}{Definition}{8}\label{def:8}
Let $R$ be a H.O. The vector $f$ is an \emph{extension element} of $R$, and $f^*$ is \emph{assigned to $f$ by $R$}, if
\[
(f,Rg)=(f^*,g)
\]
for every $g$ for which $Rg$ is defined. The pair $f,f^*$ belongs to the $+$-, $0$-, or $-$-class according as
\[
\operatorname{Im}(f^*,f)>0,\qquad =0,\qquad <0.
\]
\end{namedstatement}

\begin{namedstatement}{Theorem}{11}\label{thm:11}
The vector $f^*$ is uniquely determined by $f$ (cf. \hyperref[def:8]{Definition 8}). If $Rf$ is defined, then $f$ is an extension element of the $0$-class, and $f^*=Rf$. The operator $R$ is maximal if and only if there are no further extension elements of the $0$-class.
\end{namedstatement}

\begin{proof}
If both $f^*$ and $f^{**}$ were assigned to $f$, then
\[
(f^*,g)=(f^{**},g),\qquad (f^*-f^{**},g)=0
\]
for every $g$ for which $Rg$ is defined, and therefore also for the entire closed linear span of such $g$, namely $\mathfrak H$. It follows once again that $f^*-f^{**}=0$. The second assertion is clear. For the third, the sufficiency of the condition is clear, while its necessity follows because any extension element of the $0$-class could be used to form a proper extension of $R$.
\end{proof}

The characterization of maximality given by \hyperref[thm:11]{Theorem 11} suggests the following concept.

\begin{namedstatement}{Definition}{9}\label{def:9}
A H.O. $R$ is called \emph{hypermaximal} (abbreviated: hypermax.) if it has no extension elements other than those for which $R$ is defined. (That is, if $R$ is maximal and the $+$- and $-$-classes are empty.)
\end{namedstatement}

The importance of this concept\footnote{\label{fn:31}It is due to E. Schmidt; cf. \hyperref[fn:source-23]{note 23}. When one restricts oneself to real Hilbert space, the distinction between maximal and hypermaximal does not arise.} will appear later; cf., for example, \hyperref[thm:36]{Theorem 36}.

We prove one more theorem concerning the continuity of linear operators.

\begin{namedstatement}{Theorem}{12}\label{thm:12}
For a linear operator $R$, each of the following assertions is equivalent to continuity:
\[
\begin{array}{ll}
\alpha)& |Rf|\le C|f|,\\[2mm]
\beta)& |(f,Rg)|\le C|f|\,|g|.
\end{array}
\]
If $R$ is a H.O., the following is equivalent as well:
\[
\gamma)\qquad |(f,Rf)|\le C|f|^2.\footnote{\label{fn:32}For a H.O. every $(f,Rf)$ must be real, since
\[
(f,Rf)=(Rf,f)=\overline{(f,Rf)}.
\]}
\]
\end{namedstatement}

\begin{namedstatement}{Remark}{}\label{rem:12}
Following Hilbert, for linear Hermitian operators we also say \emph{bounded} instead of continuous. We do not investigate continuity for nonlinear operators. (Hilbert's definition is $\beta)$.)
\end{namedstatement}

\begin{proof}
By linearity, $\alpha)$ gives
\[
|Rf-Rg|=|R(f-g)|\le C|f-g|,
\]
and hence continuity. In $\beta)$ we need only put $f=Rg$ to obtain
\[
|Rg|^2\le C|Rg|\,|g|,\qquad |Rg|\le C|g|,
\]
which is $\alpha)$. If $R$ is continuous, then, because $R0=0$ by linearity, there is an $\varepsilon>0$ such that $|f|\le\varepsilon$ implies $|Rf|\le1$. Hence, again by linearity,
\[
|Rf|\le\frac1\varepsilon |f|,
\]
and therefore
\[
|(f,Rg)|\le |f|\,|Rg|\le\frac1\varepsilon |f|\,|g|;
\]
thus $\beta)$ holds. Consequently $\alpha)$ and $\beta)$ really are equivalent to continuity.

Plainly $\gamma)$ says less than $\beta)$, so for a H.O. it is enough to derive $\beta)$ from $\gamma)$. This is done as in \hyperref[thm:1]{Theorem 1}. Substitute $(f+g)/2$ and $(f-g)/2$ for $f$ and subtract; the result is
\[
\operatorname{Re}(f,Rg)\le
C\left(\frac12|f|^2+\frac12|g|^2\right).
\]
As in \hyperref[thm:1]{Theorem 1}, replace $f,g$ by $af,a^{-1}g$ with $a>0$ and take the minimum of the right-hand side; then replace $f,g$ by $\vartheta f,g$ with $|\vartheta|=1$ and take the maximum of the left-hand side. We obtain
\[
|(f,Rg)|\le C|f|\,|g|,
\]
that is, $\beta)$. Here, however, $Rf$ and $Rg$ were assumed to be defined. We can dispense with the definedness of $Rf$: $Rf$ does not occur in $\beta)$, and the inequality holds for every $f$ for which $Rf$ is defined, hence on an everywhere-dense set (because $R$ is a linear H.O.); since both sides are continuous, it holds for every $f$ whatsoever.
\end{proof}

For a linear H.O. the condition for boundedness is
\[
-C|f|^2\le(f,Rf)\le C|f|^2
\]
(\hyperref[thm:12]{Theorem 12}, $\gamma)$). We split this condition in the usual way:

\begin{namedstatement}{Definition}{10}\label{def:10}
A linear H.O. $R$ is called \emph{semibounded above}, respectively \emph{semibounded below}, if always
\[
(f,Rf)\le C|f|^2,
\qquad\text{respectively}\qquad
(f,Rf)\ge -C|f|^2.
\]
(Together, the two conditions are equivalent to boundedness by \hyperref[thm:12]{Theorem 12}, $\gamma)$.) If, in particular,
\[
(f,Rf)\ge0,
\]
then $R$ is called \emph{definite}.
\end{namedstatement}

\section*{III. Linear Manifolds and Projection Operators}\label{chap:III}
\addcontentsline{toc}{section}{III. Linear Manifolds and Projection Operators}

\begin{namedstatement}{Definition}{11}\label{def:11}
Let $\mathfrak M,\mathfrak N$ be two linear manifolds, not necessarily closed. Then $\mathfrak M+\mathfrak N$ and $\mathfrak M\times\mathfrak N$, respectively the linear manifold spanned by their union and their intersection, are again linear manifolds. If, in particular, $\mathfrak M,\mathfrak N$ are closed linear manifolds, then $\mathfrak M\times\mathfrak N$ plainly is one as well. (Whether $\mathfrak M+\mathfrak N$ is also closed is uncertain; for orthogonal $\mathfrak M,\mathfrak N$ it follows from \hyperref[thm:17]{Theorem 17}.) In what follows, $\mathfrak M,\mathfrak N$ are to be assumed closed linear manifolds. If $\mathfrak M$ is a subset of $\mathfrak N$, then their difference $\mathfrak N-\mathfrak M$, that is, the set of all elements of $\mathfrak N$ orthogonal to every element of $\mathfrak M$, is also a closed linear manifold. The manifold $\mathfrak H-\mathfrak M$ is also called the \emph{complement} of $\mathfrak M$.
\end{namedstatement}

\begin{namedstatement}{Theorem}{13}\label{thm:13}
Let $\mathfrak M$ be a closed linear manifold. Every $f$ can be decomposed in one and only one way as
\[
f=g+h,\qquad g\in\mathfrak M,\quad h\in\mathfrak H-\mathfrak M.
\]
\end{namedstatement}

\begin{namedstatement}{Remark}{}\label{rem:13}
The vector $g$ is then called the \emph{projection} of $f$ into $\mathfrak M$. The assignment of $g$ to $f$ is plainly an operator defined everywhere; we call it the \emph{projection operator} of $\mathfrak M$ and denote it by $P_{\mathfrak M}$ (where $\mathfrak M$ is a closed linear manifold).
\end{namedstatement}

\begin{proof}
Uniqueness is clear. From
\[
g_1+h_1=g_2+h_2
\]
with $g_1,g_2\in\mathfrak M$ and $h_1,h_2\in\mathfrak H-\mathfrak M$, it follows that
\[
g_1-g_2=h_2-h_1.
\]
This vector belongs both to $\mathfrak M$ and to $\mathfrak H-\mathfrak M$, and is therefore orthogonal to itself and hence $0$; thus $g_1=g_2$ and $h_1=h_2$. It remains to prove existence.

Since $\mathfrak H$ is separable, there is a sequence $f_1,f_2,\ldots$ dense in $\mathfrak M$,\footnote{\label{fn:33}Every subset of a separable topological space is itself separable; see, for example, F. Hausdorff, \emph{Einführung in die Mengenlehre}, Leipzig--Berlin, 1914.} and this sequence therefore spans the closed linear manifold $\mathfrak M$. By \hyperref[thm:8]{Theorem 8}, there is a normalized orthogonal system $\varphi_1,\varphi_2,\ldots$ that does the same. By \hyperref[thm:6]{Theorem 6},
\[
g=\sum_{n=1}^{\infty}(f,\varphi_n)\varphi_n
\]
converges; together with $\varphi_1,\varphi_2,\ldots$, its value belongs to $\mathfrak M$. The vector $h=f-g$ is orthogonal to every $\varphi_1,\varphi_2,\ldots$, and therefore also to their closed linear span $\mathfrak M$; hence $h\in\mathfrak H-\mathfrak M$. This proves everything.
\end{proof}

\begin{namedstatement}{Theorem}{14}\label{thm:14}
The projection operator $P_{\mathfrak M}$ is an everywhere-defined maximal H.O., and
\[
P_{\mathfrak M}^{\,2}=P_{\mathfrak M}.
\]
\end{namedstatement}

\begin{proof}
We know that $P_{\mathfrak M}$ is defined everywhere. It is a H.O. if we can prove
\[
(f,P_{\mathfrak M}g)=(P_{\mathfrak M}f,g).
\]
Writing
\[
f=h+k,\qquad g=h'+k',
\]
where $h,h'\in\mathfrak M$ and $k,k'\in\mathfrak H-\mathfrak M$, this is the assertion $(f,h')=(h,g)$. But $(k,h')=0$ and $(h,k')=0$, so
\[
(f,h')=(h+k,h')=(h,h')=(h,h'+k')=(h,g).
\]
If $f\in\mathfrak M$, the decomposition of \hyperref[thm:13]{Theorem 13} is $f=f+0$, and hence $P_{\mathfrak M}f=f$. Since $P_{\mathfrak M}f\in\mathfrak M$, it follows that
\[
P_{\mathfrak M}(P_{\mathfrak M}f)=P_{\mathfrak M}f,
\]
that is, $P_{\mathfrak M}^{\,2}=P_{\mathfrak M}$. Since $P_{\mathfrak M}$ is defined everywhere, it must be maximal.
\end{proof}

The maximal Hermitian operators $E$ having the property $E^2=E$ are, however, of independent interest to us (they correspond to Hilbert's ``single forms''); we therefore first investigate them in their own right.

\begin{namedstatement}{Theorem}{15}\label{thm:15}
For every maximal H.O. $E$ with the property
\[
E^2=E,\footnote{\label{fn:34}That is, if $Ef$ is defined, then $E(Ef)$ is also defined and equals $Ef$.}
\]
there is a closed linear manifold $\mathfrak M$ such that
\[
P_{\mathfrak M}=E.
\]
The manifold $\mathfrak M$ is the range of $E$.
\end{namedstatement}

\begin{proof}
Since $\mathfrak M$ is plainly the range of $P_{\mathfrak M}$, it suffices to show that $P_{\mathfrak M}=E$ when $\mathfrak M$ is the range of $E$.

Because $E^2=E$, as is easily seen, $\mathfrak M$ is the set of all $f$ satisfying $Ef=f$; thus it is a closed linear manifold (since $E$ is maximal, it is also linear and closed). Whenever $Ef$ is defined, $f-Ef$ is orthogonal to every element of $\mathfrak M$, that is, to every $Eg$:
\[
(Eg,f-Ef)=(g,Ef-E^2f)=0.
\]
Hence $f-Ef\in\mathfrak H-\mathfrak M$. It follows that $P_{\mathfrak M}f=Ef$, so $P_{\mathfrak M}$ is an extension of $E$. Since $E$ is maximal, therefore $P_{\mathfrak M}=E$.
\end{proof}

\begin{namedstatement}{Definition}{12}\label{def:12}
An operator satisfying the conditions of \hyperref[thm:15]{Theorem 15} is called a \emph{projection operator} (abbreviated: P.O.; by \hyperref[thm:15]{Theorem 15} the P.O.'s are identical with the $P_{\mathfrak M}$ and correspond one-to-one to the closed linear manifolds $\mathfrak M$). A closed linear manifold $\mathfrak M$ and its P.O. $P_{\mathfrak M}$ are called \emph{corresponding}. It is immediately seen that the P.O.'s corresponding to the closed linear manifolds $(0)$\footnote{\label{fn:35}The set $(0)$ is the set whose sole element is the zero vector of $\mathfrak H$. We do not regard the empty set as a linear manifold.} and $\mathfrak H$ are respectively $0$ and $1$,\footnote{\label{fn:36}The symbols $0$ and $1$ denote two everywhere-defined operators, defined by $0f=0$ and $1f=f$.} and furthermore
\[
P_{\mathfrak H-\mathfrak M}=1-P_{\mathfrak M}.
\]
(Thus $0$ and $1$ are P.O.'s, and if $E$ is one, then so is $1-E$.)
\end{namedstatement}

\begin{namedstatement}{Theorem}{16}\label{thm:16}
A P.O. $E$ is always bounded and definite; in fact,
\[
(f,Ef)=|Ef|^2,\qquad |Ef|\le|f|.
\]
\end{namedstatement}

\begin{proof}
We have
\[
(f,Ef)=(f,E^2f)=(Ef,Ef)=|Ef|^2,
\]
and hence $(f,Ef)\ge0$. Apply this to the P.O. $1-E$:
\[
(f,(1-E)f)\ge0,\qquad (f,Ef)\le(f,f)=|f|^2.
\]
Thus $|Ef|^2\le|f|^2$, and hence $|Ef|\le|f|$. This proves the displayed formulas, which in turn imply boundedness and definiteness.
\end{proof}

\begin{namedstatement}{Theorem}{17}\label{thm:17}
Let $\mathfrak M,\mathfrak N$ be two closed linear manifolds, and let $E,F$ be their corresponding P.O.'s. Then $EF$ is a P.O. if and only if
\[
EF=FE,
\]
that is, if $E$ and $F$ commute; in that case we also call $\mathfrak M,\mathfrak N$ \emph{commuting}, and $EF$ is the P.O. of $\mathfrak M\times\mathfrak N$.

The operator $E+F$ is a P.O. if and only if
\[
EF=0
\]
(or, equivalently, if $FE=0$); this means that all of $\mathfrak M$ is orthogonal to all of $\mathfrak N$. We then also call $E,F$ \emph{orthogonal}, and $E+F$ is the P.O. of $\mathfrak M+\mathfrak N$.

The operator $E-F$ is a P.O. if and only if
\[
EF=F
\]
(or, equivalently, if $FE=F$); this means that $\mathfrak N$ is a subset of $\mathfrak M$. We then also call $F$ a \emph{part} of $E$, and $E-F$ is the P.O. of $\mathfrak M-\mathfrak N$.\footnote{\label{fn:37}Since these operators are defined everywhere, we may form $EF$ and $E\pm F$ without hesitation; otherwise their domains would have to be examined more precisely.}
\end{namedstatement}

\begin{proof}
The operator $EF$ is defined everywhere, and hence is a maximal H.O. if it satisfies the symmetry condition
\[
(EFf,g)=(f,EFg).
\]
Now
\[
(f,EFg)=(Ef,Fg)=(FEf,g),
\]
and therefore
\[
(EFf,g)-(f,EFg)=((EF-FE)f,g).
\]
For this to vanish for every $g$, one must have $(EF-FE)f=0$ for every $f$, and hence $EF-FE=0$. Thus our condition is already necessary and sufficient for $EF$ to be a maximal H.O. Since it implies
\[
(EF)^2=EFEF=EEFF=E^2F^2=EF,
\]
it is also necessary and sufficient for $EF$ to be a P.O.

For $E+F$, the H.O. property is clear, as is maximality (the operator is defined everywhere); it therefore suffices to examine when $(E+F)^2=E+F$. Since
\[
(E+F)^2=E^2+F^2+EF+FE=(E+F)+(EF+FE),
\]
this condition says that $EF+FE=0$. But this implies $EF=0$: multiply on the left by $E$ to obtain
\[
EF+EFE=0,
\]
then multiply this equality on the right by $E$ to obtain
\[
2EFE=0;
\]
together these give $EF=0$.\footnote{\label{fn:38}This device is due to E. Schmidt.} Conversely, if $EF=0$, then $EF$ is a P.O., and hence $FE=EF=0$ and $EF+FE=0$. Thus $EF=0$ is indeed characteristic; interchanging $E$ and $F$ shows that $FE=0$ is equivalent.

The operator $E-F$ is a P.O. if and only if
\[
1-(E-F)=(1-E)+F
\]
is one. Since $1-E$ and $F$ are P.O.'s, this says
\[
(1-E)F=0,\qquad F-EF=0,
\]
or, equivalently,
\[
F(1-E)=0,\qquad F-FE=0.
\]

It remains to prove the assertions concerning $\mathfrak M,\mathfrak N,\mathfrak M\times\mathfrak N$, and $\mathfrak M\pm\mathfrak N$. First suppose $EF=FE$. Every $EFf=FEf$ belongs to both $\mathfrak M$ and $\mathfrak N$, hence to $\mathfrak M\times\mathfrak N$. Conversely, if $f\in\mathfrak M\times\mathfrak N$, then $f\in\mathfrak M$ and $f\in\mathfrak N$, so
\[
EFf=Ef=f.
\]
Thus $\mathfrak M\times\mathfrak N$ is the range of $EF$, and $EF$ is its P.O.

Second suppose $EF=0$ (and hence $FE=0$). If $f\in\mathfrak M$ and $g\in\mathfrak N$, then
\[
(f,g)=(Ef,Fg)=(f,EFg)=0;
\]
thus $\mathfrak M$ is orthogonal to $\mathfrak N$, and the converse is equally clear. For every $f$, $Ef\in\mathfrak M$ and $Ff\in\mathfrak N$.
Hence $(E+F)f\in\mathfrak M+\mathfrak N$. If $f\in\mathfrak M+\mathfrak N$, then $f=g+h$ with $g\in\mathfrak M$ and $h\in\mathfrak N$, and therefore
\begin{align*}
(E+F)f
 &=(E+F)(g+h)\\
 &=Eg+Fg+Eh+Fh\\
 &=g+0+0+h=f.
\end{align*}
Thus $\mathfrak M+\mathfrak N$ is the range of $E+F$, and is therefore a closed linear manifold, while $E+F$ is its P.O. by \hyperref[thm:15]{Theorem 15}.

Third suppose $EF=F$ (and hence $FE=F$). This says that $EFf=Ff$ for every $f$, or that $Eg=g$ for every $g\in\mathfrak N$. Since this is the defining equation of $\mathfrak M$ (cf. the proof of \hyperref[thm:15]{Theorem 15}), it says that $\mathfrak N$ is a subset of $\mathfrak M$. Since
\[
E-F=E-EF=E(1-F),
\]
$E-F$ is the P.O. of
\[
\mathfrak M\times(\mathfrak H-\mathfrak N)=\mathfrak M-\mathfrak N.
\]
All our assertions are now proved.
\end{proof}

\begin{namedstatement}{Theorem}{18}\label{thm:18}
The operator $E$ is a part of $F$ if and only if
\[
|Ef|\le|Ff|
\]
for every $f$. The sum $E_1+\cdots+E_p$ is a P.O. (where $E_1,\ldots,E_p$ are P.O.'s) if and only if
\[
E_mE_n=0\qquad(m\ne n;\;m,n=1,\ldots,p).
\]
\end{namedstatement}

\begin{proof}
If $E$ is a part of $F$, then $E=EF$, and hence
\[
|Ef|=|EFf|\le|Ff|.
\]
Conversely, suppose this inequality holds for every $f$. If $f=Ef$, then
\[
|Ff|\ge|Ef|=|f|,\qquad
(f,Ff)=|Ff|^2\ge|f|^2=(f,f),
\]
and consequently
\[
|(1-F)f|^2=(f,(1-F)f)\le0,\qquad (1-F)f=0.
\]
Thus $f=Ff$. Therefore, if $E,F$ are the P.O.'s of $\mathfrak M,\mathfrak N$, then $\mathfrak M$ is a part of $\mathfrak N$, and hence $E$ is a part of $F$.

The sufficiency of $E_mE_n=0$ for $m\ne n$ for the P.O. character of $E_1+\cdots+E_p$ follows by repeated application of \hyperref[thm:17]{Theorem 17}. Its necessity is seen as follows. For $m\ne n$,
\begin{align*}
|E_mf|^2+|E_nf|^2
 &\le |E_1f|^2+\cdots+|E_pf|^2\\
 &=(f,E_1f)+\cdots+(f,E_pf)\\
 &=\bigl(f,(E_1+\cdots+E_p)f\bigr)\le|f|^2.
\end{align*}
Thus, if $E_nf=f$, then $|E_mf|^2\le0$, so $E_mf=0$. Taking $f=E_ng$, we obtain $E_mE_ng=0$ for every $g$, that is, $E_mE_n=0$, as asserted.
\end{proof}

We conclude with a theorem on the convergence of projection operators.

\begin{namedstatement}{Theorem}{19}\label{thm:19}
Let $E_1,E_2,\ldots$ be a sequence of P.O.'s, and suppose either that $E_m$ is a part of $E_n$ for every $m<n$, or that $E_n$ is a part of $E_m$ for every $m<n$. Then there is a P.O. $E$ such that, for every $f$,
\[
E_nf\longrightarrow Ef.
\]
\end{namedstatement}

\begin{proof}
If we replace $E_1,E_2,\ldots$ by $1-E_1,1-E_2,\ldots$, the first case becomes the second,\footnote{\label{fn:39}For if $E$ is a part of $F$, then $1-F$ is a part of $1-E$, for example because
\[
(1-E)(1-F)=1-E-F+EF=1-F.
\]} so it suffices to consider the first case.

As $n$ increases, $|E_nf|^2$ increases by \hyperref[thm:18]{Theorem 18}, while it remains at most $|f|^2$ by \hyperref[thm:16]{Theorem 16}. Hence, for every $\varepsilon>0$, there is an $N=N(\varepsilon)$ such that, for all $m,n\ge N$,
\[
\bigl||E_nf|^2-|E_mf|^2\bigr|\le\varepsilon.
\]
If also $m<n$, then $E_n-E_m$ is a P.O., and therefore
\begin{align*}
|E_nf-E_mf|^2
 &=(f,(E_n-E_m)f)\\
 &=(f,E_nf)-(f,E_mf)\\
 &=|E_nf|^2-|E_mf|^2\le\varepsilon.
\end{align*}
Thus the sequence $E_1f,E_2f,\ldots$ has a limit; denote it by $Ef$.

The operator $E$ is defined everywhere. Since
\[
(f,Eg)=\lim_{n\to\infty}(f,E_ng)
      =\lim_{n\to\infty}(E_nf,g)
      =(Ef,g),
\]
it is a maximal H.O. We need only prove that $E^2=E$. We have
\begin{align*}
(f,E^2g)
 &=(Ef,Eg)
  =\lim_{n\to\infty}(E_nf,E_ng)\\
 &=\lim_{n\to\infty}(f,E_n^2g)
  =\lim_{n\to\infty}(f,E_ng)
  =(f,Eg).
\end{align*}
Thus indeed $E^2=E$.
\end{proof}

\section*{IV. Reducibility of Operators}\label{chap:IV}
\addcontentsline{toc}{section}{IV. Reducibility of Operators}

For finite-dimensional matrices, reducibility is customarily understood to mean the possibility of applying a unitary transformation to the matrix so as to bring it into a form in which---if, say, it is $N=M_1+M_2$ dimensional---only entries lying simultaneously in the first $M_1$ rows and the first $M_1$ columns, or in the last $M_2$ rows and the last $M_2$ columns, may be nonzero. If the matrix is regarded as a linear transformation, this means that the space can be decomposed into two mutually orthogonal linear manifolds of dimensions $M_1$ and $M_2$, each of which is carried by the transformation into a part of itself.

For operators in $\mathfrak H$ we formulate the analogue as follows, restricting ourselves to closed linear operators.

\begin{namedstatement}{Definition}{13}\label{def:13}
A closed linear operator $R$ is said to be \emph{reduced by} the closed linear manifold $\mathfrak M$ if, whenever $Rf$ is defined, $RP_{\mathfrak M}f$ is also defined and
\[
RP_{\mathfrak M}f=P_{\mathfrak M}Rf.\footnote{\label{fn:40}This may also be understood as the commutativity of $R$ with $P_{\mathfrak M}$.}
\]
\end{namedstatement}

Every $R$ is plainly reduced by $(0)$ and by $\mathfrak H$; and if $\mathfrak M$ reduces $R$, then so does $\mathfrak H-\mathfrak M$, since
\[
P_{\mathfrak H-\mathfrak M}=1-P_{\mathfrak M}.
\]
The operator $R$ is called \emph{irreducible} if it is reduced only by $(0)$ and $\mathfrak H$.

If $R$ is reduced by $\mathfrak M$, then for every $f\in\mathfrak M$, $Rf$, whenever defined, plainly also belongs to $\mathfrak M$, since
\[
Rf=RP_{\mathfrak M}f=P_{\mathfrak M}Rf.
\]
Thus $R$ also gives rise to an operator in $\mathfrak M$ (that is, a function whose domain and range are subsets of $\mathfrak M$). We call this operator the \emph{part of $R$ lying in $\mathfrak M$}.

We now prove two theorems on the reduction of operators.

\begin{namedstatement}{Theorem}{20}\label{thm:20}
Let $R$ be a closed linear operator and $\mathfrak M$ a closed linear manifold. In order that $\mathfrak M$ reduce $R$, it is necessary that whenever $Rf$ is defined, $RP_{\mathfrak M}f$ also be defined, and that whenever $f\in\mathfrak M$, $Rf$, if defined, also belong to $\mathfrak M$. If $R$ is a H.O., these conditions are also sufficient.\footnote{\label{fn:41}In finite-dimensional spaces these conditions are also sufficient for unitary operators; in $\mathfrak H$ they are not, but that is immaterial here.}
\end{namedstatement}

\begin{proof}
We have already established necessity. It remains to show that the conditions are sufficient for a Hermitian $R$, that is, to prove
\[
RP_{\mathfrak M}f=P_{\mathfrak M}Rf.
\]
By assumption $RP_{\mathfrak M}f\in\mathfrak M$, so it remains to show that
\[
Rf-RP_{\mathfrak M}f=R(f-P_{\mathfrak M}f)
\]
belongs to $\mathfrak H-\mathfrak M$. Since $f-P_{\mathfrak M}f\in\mathfrak H-\mathfrak M$, it is enough to show that if $g\in\mathfrak H-\mathfrak M$, then $Rg$, whenever defined, also belongs to $\mathfrak H-\mathfrak M$.

Let $Rf$ be defined. Then $RP_{\mathfrak M}f\in\mathfrak M$, and hence
\[
(f,P_{\mathfrak M}Rg)
=(P_{\mathfrak M}f,Rg)
=(RP_{\mathfrak M}f,g)=0.
\]
The set of such $f$ is everywhere dense; hence $P_{\mathfrak M}Rg$ is orthogonal to everything and is therefore $0$. Thus $Rg\in\mathfrak H-\mathfrak M$.
\end{proof}

\begin{namedstatement}{Theorem}{21}\label{thm:21}
Let $\mathfrak M_1,\mathfrak M_2,\ldots$ be a possibly finite sequence of pairwise orthogonal closed linear manifolds which together span the closed linear manifold $\mathfrak H$. Let $R$ be a closed linear operator reduced by each of them, and let the parts of $R$ lying in $\mathfrak M_1,\mathfrak M_2,\ldots$ be $R_1,R_2,\ldots$, respectively.

Then $R$ is determined by $R_1,R_2,\ldots$ as follows: $Rf$ is defined if and only if
\[
f=f_1+f_2+\cdots,\qquad f_p\in\mathfrak M_p,
\]
all $R_1f_1,R_2f_2,\ldots$ are defined, and both series
\[
f_1+f_2+\cdots,\qquad R_1f_1+R_2f_2+\cdots
\]
converge (if the sequence $\mathfrak M_1,\mathfrak M_2,\ldots$ terminates, the convergence assumption of course disappears). Its value is then
\[
Rf=R_1f_1+R_2f_2+\cdots.
\]
\end{namedstatement}

\begin{proof}
That $R$ is defined at the points just described and takes the indicated values follows immediately from the definitions of $R_1,R_2,\ldots$, together with the fact that $R$ is closed and linear and is reduced by $\mathfrak M_1,\mathfrak M_2,\ldots$. We still have to prove that, if $Rf$ is defined, then $f$ has the stated form.

Put
\[
f_p=P_{\mathfrak M_p}f\qquad(p=1,2,\ldots).
\]
Then $f_p\in\mathfrak M_p$, and the $P_{\mathfrak M_p}$, like the $\mathfrak M_p$, are pairwise orthogonal. Moreover,
\[
R_pf_p=Rf_p=RP_{\mathfrak M_p}f=P_{\mathfrak M_p}Rf.
\]
It follows from the orthogonality of the $P_{\mathfrak M_p}$ that
\[
P_{\mathfrak M_1},\quad
P_{\mathfrak M_1}+P_{\mathfrak M_2},\quad
P_{\mathfrak M_1}+P_{\mathfrak M_2}+P_{\mathfrak M_3},\quad\ldots
\]
form an increasing sequence of P.O.'s, which by \hyperref[thm:19]{Theorem 19} has a P.O. $E$ as its limit. By continuity, every $\mathfrak M_p$ is a part of the closed linear manifold $\mathfrak M$ corresponding to $E$. By the assumption on $\mathfrak M_1,\mathfrak M_2,\ldots$, therefore $\mathfrak M=\mathfrak H$ and $E=1$.

Thus the series
\[
f_1+f_2+\cdots,\qquad R_1f_1+R_2f_2+\cdots
\]
converge, with respective sums
\[
1f=f,\qquad 1Rf=Rf.
\]
This proves everything.
\end{proof}

\section*{V. The Cayley Transform}\label{chap:V}
\addcontentsline{toc}{section}{V. The Cayley Transform}

After these preparations we can take up our actual task, the closer investigation of the various kinds of Hermitian operators. The basis of the following considerations is a device that transfers the Cayley transform of algebra to our setting.

Let $R$ be a closed linear H.O., and let $\mathfrak A$ be its domain. The set $\mathfrak A$ is a linear manifold and is everywhere dense (and is therefore closed only when $\mathfrak A=\mathfrak H$, that is, when $R$ is defined everywhere, which is by no means assumed). Consider the two operators $R\pm i1$, which are also defined on $\mathfrak A$. A direct calculation gives
\begin{align*}
((R\pm i1)f,(R\pm i1)g)
 &=(Rf,Rg)\pm(Rf,ig)\pm(if,Rg)+(f,g)\\
 &=(Rf,Rg)+(f,g).
\end{align*}
In particular, for $f=g$,
\[
|(R\pm i1)f|=|(R-i1)f|
=\sqrt{|Rf|^2+|f|^2}.
\]
Thus $f\ne0$ implies both $(R+i1)f\ne0$ and $(R-i1)f\ne0$, and by linearity $f\ne g$ implies
\[
(R+i1)f\ne(R+i1)g,\qquad
(R-i1)f\ne(R-i1)g.
\]
Suppose $R+i1$ and $R-i1$ map $\mathfrak A$ onto $\mathfrak E$ and $\mathfrak F$, respectively. These mappings among $\mathfrak A,\mathfrak E,\mathfrak F$ are therefore all one-to-one. In particular, denote the mapping from $\mathfrak E$ onto $\mathfrak F$ by $U$. Since $R\pm i1$, like $R$, are closed and linear, the same is true of $U$; and the preceding equality gives
\[
|Uf|=|f|.
\]
Thus $U$ is length-preserving by \hyperref[def:6]{Definition 6}. Since $\mathfrak E,\mathfrak F$ are its domain and range, both are closed linear manifolds by \hyperref[thm:10]{Theorem 10}.

\begin{namedstatement}{Theorem}{22}\label{thm:22}
Let $R$ be a closed linear H.O. with domain $\mathfrak A$, and suppose the operators $R+i1$ and $R-i1$ map $\mathfrak A$ onto $\mathfrak E$ and $\mathfrak F$, respectively. These mappings among $\mathfrak A,\mathfrak E,\mathfrak F$ are one-to-one; in particular, the mapping from $\mathfrak E$ onto $\mathfrak F$ is a length-preserving operator $U$. Both $\mathfrak E$ and $\mathfrak F$ are closed linear manifolds.
\end{namedstatement}

\begin{proof}
This was just proved.
\end{proof}

\begin{namedstatement}{Definition}{14}\label{def:14}
The operators $R,U$ of \hyperref[thm:22]{Theorem 22} are called the \emph{Cayley transforms} of one another. Here $R$ is always a closed linear H.O., and $U$ is always length-preserving.\footnote{\label{fn:42}For the moment we know only that every closed linear H.O. $R$ has exactly one Cayley transform $U$, which is length-preserving. The scope of the converse will follow from \hyperref[thm:24]{Theorem 24}.}
\end{namedstatement}

\begin{namedstatement}{Theorem}{23}\label{thm:23}
A length-preserving operator $U$ is the Cayley transform of the closed linear H.O. $R$ if and only if
\[
\begin{array}{ll}
\alpha)&\mathfrak A\text{ is the set of all }\varphi-U\varphi,\quad\varphi\in\mathfrak E,\\[1mm]
\beta)&R(\varphi-U\varphi)=i(\varphi+U\varphi)\quad\text{for every }\varphi\in\mathfrak E.
\end{array}
\]
\end{namedstatement}

\begin{proof}
Let $U$ be the Cayley transform of $R$. If $f\in\mathfrak A$, then
\[
(R+i1)f=\varphi\in\mathfrak E,\qquad U\varphi=(R-i1)f,
\]
and hence
\[
\varphi-U\varphi=2if,\qquad
f=\frac{\varphi}{2i}-U\frac{\varphi}{2i}.
\]
Conversely, if $\varphi\in\mathfrak E$, then $(R+i1)f=\varphi$ for some $f\in\mathfrak A$, and therefore
\[
(R-i1)f=U\varphi,\qquad
\varphi-U\varphi=2if.
\]
Thus $\varphi-U\varphi\in\mathfrak A$. Finally,
\[
\varphi+U\varphi=2Rf,
\]
and so
\[
R(\varphi-U\varphi)=2iRf=i(\varphi+U\varphi).
\]
Thus $\alpha)$ and $\beta)$ hold.

Conversely, suppose $U$ satisfies $\alpha)$ and $\beta)$. If $f\in\mathfrak A$, then
\[
f=\varphi-U\varphi,\qquad \varphi\in\mathfrak E,
\]
by $\alpha)$, and
\[
Rf=i(\varphi+U\varphi)
\]
by $\beta)$. Hence
\[
(R+i1)f=2i\varphi,\qquad (R-i1)f=2iU\varphi.
\]
Thus $(R+i1)f\in\mathfrak E$ and
\[
U(R+i1)f=(R-i1)f.
\]
On the other hand, starting with any $\varphi\in\mathfrak E$, we have
\[
f=\varphi-U\varphi\in\mathfrak A,\qquad
Rf=i(\varphi+U\varphi),
\]
and hence
\[
2i\varphi=(R+i1)f,\qquad
\varphi=(R+i1)\frac{f}{2i}.
\]
Thus $\mathfrak E$ is precisely the set of all $(R+i1)f$ with $f\in\mathfrak A$, and
\[
U(R+i1)f=(R-i1)f.
\]
Therefore $U$ really is the Cayley transform of $R$.
\end{proof}

\begin{namedstatement}{Theorem}{24}\label{thm:24}
For a length-preserving operator $U$ there exists a closed linear H.O. $R$ of which it is the Cayley transform if and only if the set of all
\[
\varphi-U\varphi,\qquad\varphi\in\mathfrak E,
\]
where $\mathfrak E$ is the domain of $U$, is everywhere dense. In that case $R$ is uniquely determined by $U$.
\end{namedstatement}

\begin{proof}
Necessity follows from $\alpha)$ of \hyperref[thm:23]{Theorem 23}, since $\mathfrak A$ is everywhere dense. Sufficiency and the unique determination of $R$ by $U$ likewise follow from $\alpha)$ and $\beta)$ of \hyperref[thm:23]{Theorem 23}, once we know that $\beta)$ does not involve an inconsistency in the definition of $R$, and that the $R$ so defined is a closed linear H.O.

The first point follows from the fact that $\varphi\ne\psi$ implies
\[
\varphi-U\varphi\ne\psi-U\psi,
\]
or equivalently that $\varphi\ne0$ implies $\varphi-U\varphi\ne0$. Indeed, if $\varphi=U\varphi$, then, for every $\chi\in\mathfrak E$,
\[
(\varphi,\chi)=(U\varphi,U\chi)=(\varphi,U\chi),
\qquad
(\varphi,\chi-U\chi)=0.
\]
But the vectors $\chi-U\chi$ are everywhere dense, so $\varphi$ is orthogonal to an everywhere-dense set, hence to all of $\mathfrak H$, and therefore $\varphi=0$.

For the second point, the closedness and linearity of $R$ follow from those of $U$, while the density of the domain of $R$, namely the set of all $\varphi-U\varphi$, was assumed. It remains only to prove
\[
(f,Rg)=(Rf,g),
\]
that is, that interchanging $f$ and $g$ in $(f,Rg)$ changes it into its complex conjugate. Put
\[
f=\varphi-U\varphi,\qquad g=\psi-U\psi.
\]
Then
\begin{align*}
(f,Rg)
 &=(\varphi-U\varphi,i(\psi+U\psi))\\
 &=-i\{(\varphi,\psi)-(U\varphi,\psi)
      +(\varphi,U\psi)-(U\varphi,U\psi)\}\\
 &=i\left((U\varphi,\psi)-\overline{(U\psi,\varphi)}\right)\\
 &=i(U\varphi,\psi)+\overline{i(U\psi,\varphi)},
\end{align*}
which makes the assertion evident.
\end{proof}

\begin{namedstatement}{Theorem}{25}\label{thm:25}
Let $R,S$ be closed linear Hermitian operators, and let $U,V$ be their Cayley transforms. Then $R$ is an extension, respectively a proper extension, of $S$ if and only if $U$ is an extension, respectively a proper extension, of $V$. A closed linear manifold $\mathfrak M$ reduces $R$ if and only if it reduces $U$.
\end{namedstatement}

\begin{proof}
Each assertion follows immediately from $R$ to $U$ by \hyperref[thm:22]{Theorem 22}, and from $U$ to $R$ by \hyperref[thm:23]{Theorem 23}.
\end{proof}

In the preceding theorems we reduced the problem of deciding the extension relations of $R$ to the much simpler corresponding problem for $U$. The Cayley transforms of all $R$ are precisely all the $U$ of \hyperref[thm:5]{Theorem 5}; here one must note that if $R$ is to be an extension of $S$, so that $U$ is an extension of a $V$ satisfying \hyperref[thm:5]{Theorem 5}, then $U$ need only be length-preserving.\footnote{\label{fn:43}For the vectors $\varphi-U\varphi$ include the vectors $\varphi-V\varphi$, and hence, like the latter, are everywhere dense.} Thus, instead of considering the initially rather opaque extensions of $R$, we need only consider the easily described length-preserving extensions of $V$, which is itself length-preserving. To orient ourselves in questions of this kind, we prove two simple, geometrically intuitive theorems about length-preserving mappings and closed linear manifolds.

First observe that by the \emph{dimension} of a closed linear manifold $\mathfrak M$ we mean the supremum of all finite $n$ for which $\mathfrak M$ contains $n$ linearly independent elements; it is therefore one of the numbers
\[
0,1,2,\ldots,\infty.
\]
(A substantially more general definition of the dimension of sets in $\mathfrak H$ will be given in the next chapter, \hyperref[def:16]{Definition 16}.)

\begin{namedstatement}{Theorem}{26}\label{thm:26}
Let $\mathfrak M$ be a closed linear manifold, and let $\varphi_1,\ldots,\varphi_n$ be a normalized orthogonal system spanning it (such a system certainly exists; cf. the proof of \hyperref[thm:13]{Theorem 13}; here $n=0,1,2,\ldots,\infty$, and in the last case the sequence $\varphi_1,\varphi_2,\ldots$ does not terminate). Then $\mathfrak M$ is $n$-dimensional.

Two closed linear manifolds $\mathfrak M,\mathfrak N$ can be respectively the domain and the range of a length-preserving operator $U$ if and only if they have the same dimension. Let $\varphi_1,\ldots,\varphi_n$ be a normalized orthogonal system spanning $\mathfrak M$. Then the normalized orthogonal systems $\psi_1,\ldots,\psi_n$ spanning $\mathfrak N$ correspond one-to-one to such operators $U$: to each such system there corresponds exactly one $U$ satisfying
\[
U\varphi_1=\psi_1,\ldots,U\varphi_n=\psi_n.
\]
\end{namedstatement}

\begin{proof}
Let $\varphi_1,\ldots,\varphi_n$ be a normalized orthogonal system spanning the closed linear manifold $\mathfrak M$. Since $\varphi_1,\ldots,\varphi_n$ are linearly independent, $\mathfrak M$ has at least $n$ dimensions. Thus, when $n=\infty$, we are done. For finite $n$, we must show that $n+1$ elements of $\mathfrak M$ can never be linearly independent. Indeed, they must have the form
\[
f_\mu=a_{\mu1}\varphi_1+\cdots+a_{\mu n}\varphi_n
\qquad(\mu=1,\ldots,n+1),
\]
and there must be a linear relation among the $n+1$ vectors in $n$ dimensions
\[
(a_{\mu1},\ldots,a_{\mu n})
\qquad(\mu=1,\ldots,n+1),
\]
and hence also among the $f_\mu$.

We turn to the second and third assertions. Let $\mathfrak M$ be the domain and $\mathfrak N$ the range of the length-preserving operator $U$. If $\varphi_1,\ldots,\varphi_n$ is a normalized orthogonal system spanning $\mathfrak M$, then, because $U$ preserves length,
\[
\psi_1=U\varphi_1,\ldots,\psi_n=U\varphi_n
\]
is a normalized orthogonal system spanning $\mathfrak N$. Thus $\mathfrak M$ and $\mathfrak N$ both have $n$ dimensions.

Conversely, let $\mathfrak M,\mathfrak N$ be two closed linear manifolds of the same dimension, say $n$. Let
\[
\varphi_1,\ldots,\varphi_p
\qquad\text{and}\qquad
\psi_1,\ldots,\psi_q
\]
be normalized orthogonal systems spanning $\mathfrak M$ and $\mathfrak N$, respectively. We must have $p=q=n$. The series
\[
\sum_{r=1}^{n}x_r\varphi_r,\qquad
\sum_{r=1}^{n}x_r\psi_r
\]
converge under the same circumstances: always when $n$ is finite, and when $n=\infty$ if $\sum_{r=1}^{\infty}|x_r|^2$ is finite, by \hyperref[thm:5]{Theorem 5}---and only then. They therefore exhaust $\mathfrak M$ and $\mathfrak N$, respectively. (The inner product of $\sum_{r=1}^{n}x_r\varphi_r$ with $\varphi_s$ is $x_s$, so it determines $x_1,\ldots,x_n$ uniquely.) We can therefore define an operator $U$ by
\[
U\left(\sum_{r=1}^{n}x_r\varphi_r\right)
 =\sum_{r=1}^{n}x_r\psi_r.
\]
It has domain $\mathfrak M$ and range $\mathfrak N$, and is plainly linear. If
\[
f=\sum_{r=1}^{n}x_r\varphi_r,
\]
then
\[
|f|^2=|Uf|^2=\sum_{r=1}^{n}|x_r|^2,
\qquad |f|=|Uf|.
\]
Thus the operator is continuous, and since its domain is closed, it is itself closed. Hence $U$ is length-preserving.

Conversely, length preservation, and hence in particular closed linearity, together with
\[
U\varphi_\mu=\psi_\mu\qquad(\mu=1,\ldots,n),
\]
implies
\[
U\left(\sum_{r=1}^{n}x_r\varphi_r\right)
=\sum_{r=1}^{n}x_r\psi_r.
\]
Our $U$ is therefore uniquely determined by these properties. This proves all the assertions of the theorem.
\end{proof}

\begin{namedstatement}{Theorem}{27}\label{thm:27}
Let $U$ be a length-preserving operator, and let $\mathfrak E,\mathfrak F$ be its domain and range (thus two closed linear manifolds). All length-preserving extensions of $U$ are obtained as follows.

Choose two closed linear manifolds $\mathfrak M,\mathfrak N$ of the same dimension, with
\[
\mathfrak M\subseteq\mathfrak H-\mathfrak E,\qquad
\mathfrak N\subseteq\mathfrak H-\mathfrak F,
\]
and a length-preserving operator $U'$ having $\mathfrak M$ as domain and $\mathfrak N$ as range. These determine exactly one length-preserving extension $V$ of $U$, characterized by the following conditions: its domain and range are, respectively,
\[
\mathfrak E+\mathfrak M,\qquad
\mathfrak F+\mathfrak N;
\]
on $\mathfrak E$ one has $Vf=Uf$, and on $\mathfrak M$ one has $Vf=U'f$.

By choosing $\mathfrak M,\mathfrak N,U'$ in all possible ways, one obtains all length-preserving extensions $V$ of $U$, each exactly once.
\end{namedstatement}

\begin{proof}
It is clear that every length-preserving extension $V$ of $U$ arises in this way. Let $\mathfrak K,\mathfrak L$ be the domain and range of $V$; it suffices to put
\[
\mathfrak M=\mathfrak K-\mathfrak E,\qquad
\mathfrak N=\mathfrak L-\mathfrak F,
\]
and to let $U'$ be the restriction of $V$ to $\mathfrak M$. Then all the conditions are satisfied. It is also clear that our conditions determine at most one $V$: if both $V'$ and $V''$ satisfy them, then $V'f=V''f$ throughout $\mathfrak E$ and throughout $\mathfrak M$, and hence, by linearity, throughout their common domain $\mathfrak E+\mathfrak M$; thus $V'=V''$. It remains only to show that every system $\mathfrak M,\mathfrak N,U'$ really determines at least one length-preserving $V$.

Since $\mathfrak M$ and $\mathfrak N$ are parts of $\mathfrak H-\mathfrak E$ and $\mathfrak H-\mathfrak F$, respectively, they are orthogonal to $\mathfrak E$ and $\mathfrak F$, respectively. Hence $\mathfrak M$ and $\mathfrak E$ are linearly independent, and the decomposition
\[
f=g+h,\qquad g\in\mathfrak E,\quad h\in\mathfrak M,
\]
is unique in $\mathfrak E+\mathfrak M$. Define
\[
V(g+h)=Ug+U'h.
\]
Since $g,h$ are orthogonal, as are $Ug,U'h$, and since $U,U'$ preserve length,
\begin{align*}
|V(g+h)|^2
 &=|Ug+U'h|^2\\
 &=|Ug|^2+|U'h|^2\\
 &=|g|^2+|h|^2=|g+h|^2.
\end{align*}
The operator $V$ is defined on the closed linear manifold $\mathfrak E+\mathfrak M$, is plainly linear, and by the preceding equality is continuous; it is therefore closed, and the same equality shows that it is length-preserving. Moreover, as required, it extends both $U$ and $U'$.
\end{proof}

By \hyperref[thm:26]{Theorems 26} and \hyperref[thm:27]{27} we survey all length-preserving extensions of a length-preserving $U$, and hence all closed linear Hermitian extensions of a closed linear H.O. $R$; all of them can be constructed explicitly. For example, $R$ is maximal, that is, it has no proper extension, exactly when $U$ has none---thus when it is impossible in \hyperref[thm:27]{Theorem 27}\translatornote{The source cites Theorem 26 in this sentence; the choice of \(\mathfrak M,\mathfrak N\) is made in Theorem 27.} to choose $\mathfrak M$ and $\mathfrak N$ different from $(0)$ (cf. \hyperref[fn:35]{note 35}). This is plainly the case only if
\[
\mathfrak H-\mathfrak E=(0)
\qquad\text{or}\qquad
\mathfrak H-\mathfrak F=(0),
\]
that is, if $\mathfrak E=\mathfrak H$ or $\mathfrak F=\mathfrak H$. These questions will be investigated more closely in the next three chapters.

\section*{VI. Extension Elements}\label{chap:VI}
\addcontentsline{toc}{section}{VI. Extension Elements}

We now wish to characterize more precisely the closed linear manifolds $\mathfrak E,\mathfrak F$, which already came into play in the preceding chapter, in relation to the concepts of \hyperref[def:8]{Definition 8}.

For the purposes of this chapter it is convenient to use the terminology of the residue-class calculus of algebra: if $\mathfrak M$ is a linear manifold, not necessarily closed, then
\[
f=g\pmod{\mathfrak M}
\]
means that $f-g\in\mathfrak M$. The symbols $R,U,\mathfrak A,\mathfrak E,\mathfrak F$ shall retain their previous meanings, as, for example, in \hyperref[thm:22]{Theorem 22}.

\begin{namedstatement}{Theorem}{28}\label{thm:28}
The manifolds $\mathfrak H-\mathfrak E$ and $\mathfrak H-\mathfrak F$ are, respectively, the sets of all extension elements $f$ to which $R$ assigns $if$ and $-if$.
\end{namedstatement}

\begin{proof}
By \hyperref[thm:23]{Theorem 23}, the assertion that $f^*$ is assigned to $f$ means
\[
(f^*,\varphi-U\varphi)
=(f,i(\varphi+U\varphi))
=(-if,\varphi+U\varphi).
\]
Thus $f^*=if$, respectively $f^*=-if$, means that $f$ (and with it $f^*=\pm if$) is orthogonal to every $\varphi\in\mathfrak E$, respectively to every $U\varphi$, that is, to all of $\mathfrak F$. Equivalently,
\[
f\in\mathfrak H-\mathfrak E,
\qquad\text{respectively}\qquad
f\in\mathfrak H-\mathfrak F.
\]
\end{proof}

Consequently all of $\mathfrak H-\mathfrak E$, respectively all of $\mathfrak H-\mathfrak F$, apart from $0$, belongs to the $+$-class, respectively the $-$-class, of extension elements.

\begin{namedstatement}{Theorem}{29}\label{thm:29}
The set of all extension elements of $R$ is
\[
\mathfrak A+(\mathfrak H-\mathfrak E)+(\mathfrak H-\mathfrak F).
\]
\end{namedstatement}

\begin{proof}
The extension elements plainly form a linear manifold. Since this manifold must, as we already know, contain
\[
\mathfrak A,\qquad \mathfrak H-\mathfrak E,\qquad
\mathfrak H-\mathfrak F,
\]
their sum is a part of it; it remains to show that it is no larger.

Let $f$ be an extension element. By the preceding calculation this means
\[
(f^*,\varphi-U\varphi)=(-if,\varphi+U\varphi),
\]
or
\[
(f^*+if,\varphi)=(f^*-if,U\varphi)
\]
for every $\varphi\in\mathfrak E$, with $f^*$ fixed. Let $U^{-1}$ be the inverse of $U$, with domain $\mathfrak F$ and range $\mathfrak E$. Then, wherever these operators are defined, that is, on $\mathfrak E$ or $\mathfrak F$ as appropriate,
\[
P_{\mathfrak F}U=U,\qquad U^{-1}U=1,
\qquad
P_{\mathfrak E}U^{-1}=U^{-1},\qquad UU^{-1}=1.
\]
Now
\begin{align*}
(f^*-if,U\varphi)
 &=(f^*-if,P_{\mathfrak F}U\varphi)\\
 &=(P_{\mathfrak F}f^*-iP_{\mathfrak F}f,U\varphi)\\
 &=(UU^{-1}P_{\mathfrak F}f^*
      -iUU^{-1}P_{\mathfrak F}f,U\varphi)\\
 &=(U^{-1}P_{\mathfrak F}f^*
      -iU^{-1}P_{\mathfrak F}f,\varphi).
\end{align*}
Therefore
\[
\left(
\{f^*+if\}
-\{U^{-1}P_{\mathfrak F}f^*
   -iU^{-1}P_{\mathfrak F}f\},
\varphi\right)=0.
\]
Since this holds for every $\varphi\in\mathfrak E$,
\[
P_{\mathfrak E}\left(
\{f^*+if\}
-\{U^{-1}P_{\mathfrak F}f^*
   -iU^{-1}P_{\mathfrak F}f\}
\right)=0,
\]
and hence
\[
\{P_{\mathfrak E}f^*+iP_{\mathfrak E}f\}
-\{U^{-1}P_{\mathfrak F}f^*
   -iU^{-1}P_{\mathfrak F}f\}=0.
\]
Equivalently,
\[
-U^{-1}P_{\mathfrak F}f^*+P_{\mathfrak E}f^*
=-i\bigl(U^{-1}P_{\mathfrak F}f+P_{\mathfrak E}f\bigr),
\]
or, after applying $-U$,
\[
P_{\mathfrak F}f^*-UP_{\mathfrak E}f^*
=i\bigl(P_{\mathfrak F}f+UP_{\mathfrak E}f\bigr).
\]

In general,
\[
g-P_{\mathfrak M}g=P_{\mathfrak H-\mathfrak M}g
=0\pmod{\mathfrak H-\mathfrak M}.
\]
It follows that
\[
P_{\mathfrak E}g
=P_{\mathfrak F}g
\pmod{(\mathfrak H-\mathfrak E)+(\mathfrak H-\mathfrak F)}.
\]
Hence the preceding equality gives
\[
P_{\mathfrak E}f^*-UP_{\mathfrak E}f^*
=i\bigl(P_{\mathfrak E}f+UP_{\mathfrak E}f\bigr)
\pmod{(\mathfrak H-\mathfrak E)+(\mathfrak H-\mathfrak F)}.
\]
The left-hand side belongs, by its very form, to $\mathfrak A$; therefore the right-hand side belongs to
\[
\mathfrak A+(\mathfrak H-\mathfrak E)+(\mathfrak H-\mathfrak F).
\]
Likewise,
\[
P_{\mathfrak E}f-UP_{\mathfrak E}f
\]
belongs by its very form to $\mathfrak A$. These two facts imply that $P_{\mathfrak E}f$ belongs to the same sum. Finally,
\[
f=P_{\mathfrak E}f\pmod{\mathfrak H-\mathfrak E},
\]
so the same is true of $f$. This proves everything.
\end{proof}

\begin{namedstatement}{Theorem}{30}\label{thm:30}
The linear manifolds
\[
\mathfrak A,\qquad \mathfrak H-\mathfrak E,\qquad
\mathfrak H-\mathfrak F
\]
are linearly independent.
\end{namedstatement}

\begin{proof}
We must infer from
\[
f+\varphi+\psi=0,
\qquad
f\in\mathfrak A,\quad
\varphi\in\mathfrak H-\mathfrak E,\quad
\psi\in\mathfrak H-\mathfrak F,
\]
that $f=\varphi=\psi=0$. Equivalently, it suffices to prove that if
\[
\varphi\in\mathfrak H-\mathfrak E,\qquad
\psi\in\mathfrak H-\mathfrak F,\qquad
\varphi+\psi\in\mathfrak A,
\]
then $\varphi=\psi=0$. Assume these latter premises.

By what has already been said, $\varphi,\psi,\varphi+\psi$ are extension elements to which $R$ assigns, respectively,
\[
i\varphi,\qquad -i\psi,\qquad R(\varphi+\psi),
\]
the last of which is defined. But $i(\varphi-\psi)$ is plainly also assigned to $\varphi+\psi$, and hence
\[
i(\varphi-\psi)=R(\varphi+\psi).
\]
Since $i\varphi$ and $-i\psi$ are assigned to $\varphi$ and $\psi$, respectively, we must have
\begin{align*}
(i\varphi,\varphi+\psi)
 &=(\varphi,R(\varphi+\psi))
  =(\varphi,i(\varphi-\psi)),\\
2i(\varphi,\varphi)&=0,\qquad \varphi=0,
\end{align*}
and
\begin{align*}
(-i\psi,\varphi+\psi)
 &=(\psi,R(\varphi+\psi))
  =(\psi,i(\varphi-\psi)),\\
-2i(\psi,\psi)&=0,\qquad \psi=0.
\end{align*}
\end{proof}

Thus every extension element $f$ can be decomposed in one and only one way as
\[
f=f^0+f^++f^-,
\qquad
f^0\in\mathfrak A,\quad
f^+\in\mathfrak H-\mathfrak E,\quad
f^-\in\mathfrak H-\mathfrak F;
\]
conversely, every such $f$ is an extension element. We call $f^0,f^+,f^-$, respectively, the $0$-, $+$-, and $-$-components of $f$. Since $Rf^0,if^+,-if^-$ are assigned to $f^0,f^+,f^-$, respectively, $R$ assigns to
\[
f=f^0+f^++f^-
\]
the vector
\[
f^*=Rf^0+if^+-if^-.
\]
This decomposition is important because it permits a more precise characterization of the $0$-, $+$-, and $-$-classes.

\begin{namedstatement}{Theorem}{31}\label{thm:31}
An extension element $f$ belongs to the $0$-, $+$-, or $-$-class according as
\[
|f^+|=|f^-|,\qquad |f^+|>|f^-|,\qquad |f^+|<|f^-|.
\]
\end{namedstatement}

\begin{proof}
Since
\[
f=f^0+f^++f^-,
\qquad
f^*=Rf^0+if^+-if^-,
\]
and since $if^+$ and $-if^-$ are assigned to $f^+$ and $f^-$, respectively, we have
\begin{align*}
(f^*,f)
={}&(Rf^0+if^+-if^-,f^0+f^++f^-)\\
={}&(Rf^0,f^0)
 +\{(Rf^0,f^+)+(if^+,f^0)\}\\
&+\{(Rf^0,f^-)+(-if^-,f^0)\}\\
&+\{(if^+,f^-)+(-if^-,f^+)\}\\
&+\{(if^+,f^+)+(-if^-,f^-)\}\\
={}&(f^0,Rf^0)
 +\{(f^0,if^+)+(if^+,f^0)\}\\
&+\{(f^0,-if^-)+(-if^-,f^0)\}\\
&+\{(f^+,-if^-)+(-if^-,f^+)\}\\
&+i\{|f^+|^2-|f^-|^2\}.
\end{align*}
The first term and the first three brace-enclosed expressions are plainly real, whereas the last brace-enclosed contribution is purely imaginary. It follows that
\[
\operatorname{Im}(f^*,f)=|f^+|^2-|f^-|^2,
\]
which proves the assertion.
\end{proof}

\section*{VII. The Deficiency Indices}\label{chap:VII}
\addcontentsline{toc}{section}{VII. The Deficiency Indices}

\begin{namedstatement}{Definition}{15}\label{def:15}
Let $R,U,\mathfrak A,\mathfrak E,\mathfrak F$ have their previous meanings, as in \hyperref[thm:22]{Theorem 22}. Let $m,n$ be the dimensions of
\[
\mathfrak H-\mathfrak E,\qquad
\mathfrak H-\mathfrak F,
\]
respectively. We call $m,n$ the \emph{deficiency indices} of $R$. Thus
\[
m,n=0,1,2,\ldots,\infty.
\]
\end{namedstatement}

These numbers will play a decisive role in questions of maximality and hypermaximality. There is a close relation between them and the three classes of extension elements, which we shall now establish. First we generalize the notion of dimension in $\mathfrak H$.

\begin{namedstatement}{Definition}{16}\label{def:16}
Let $\mathfrak M$ be a linear manifold, not necessarily closed, and let $\mathfrak U$ be a subset of $\mathfrak H$, which need not even be a linear manifold. By the \emph{dimension of $\mathfrak U$ modulo $\mathfrak M$} we mean the supremum of all finite $n$ for which there exist $n$ linearly independent elements whose entire linear manifold,\footnote{\label{fn:44}Since $n$ is finite, this linear manifold is automatically closed.} possibly with the exception of $0$,\footnote{\label{fn:45}This exception will prove convenient below.} lies in $\mathfrak U$ and outside $\mathfrak M$. Thus $n$ is one of the numbers
\[
0,1,2,\ldots,\infty.
\]
\end{namedstatement}

In what follows, $\mathfrak U$ will not in fact be entirely arbitrary, but will always have the following properties. First, it is a set of residue classes modulo $\mathfrak M$: if
\[
f=g\pmod{\mathfrak M},
\]
then either neither or both of $f,g$ belong to $\mathfrak U$. Second, whenever $f\in\mathfrak U$, then $af\in\mathfrak U$ for every $a\ne0$. It is immediately seen that, among such sets $\mathfrak U$, only the empty set and $\mathfrak M$ have dimension $0$. It is also clear that the dimension of a closed linear manifold $\mathfrak U$ modulo $(0)$ agrees with the dimension defined in \hyperref[chap:V]{Chapter V}.

The $+$-, $-$-, and $0$-classes are sets of this kind modulo $\mathfrak A$, as are unions of any of them. Since the first two do not contain $0$, whereas the third does, in the zero-dimensional case they must be empty and equal to $\mathfrak A$, respectively.

\begin{namedstatement}{Theorem}{32}\label{thm:32}
The $+$-class, and also its union with the $0$-class, is $m$-dimensional modulo $\mathfrak A$; the $-$-class, and also its union with the $0$-class, is $n$-dimensional modulo $\mathfrak A$; the $0$-class is $\min(m,n)$-dimensional modulo $\mathfrak A$.
\end{namedstatement}

\begin{proof}
Let $p\le m$ be finite, as $m$ may be $\infty$, and let $f_1,\ldots,f_p$ be linearly independent elements of the $m$-dimensional manifold $\mathfrak H-\mathfrak E$. Their entire linear manifold lies in $\mathfrak H-\mathfrak E$, and hence, except for $0$, in the $+$-class, and therefore automatically outside $\mathfrak A$. Thus the $+$-class has at least $m$ dimensions modulo $\mathfrak A$. To prove the first two assertions, it remains to show that its union with the $0$-class has at most $m$ dimensions.

For $m=\infty$ this is trivial. If $m$ is finite, we must show that, whenever $f_1,\ldots,f_{m+1}$ are linearly independent extension elements, some nonzero combination
\[
a_1f_1+\cdots+a_{m+1}f_{m+1}\ne0
\]
belongs either to $\mathfrak A$ or to the $-$-class.

Write
\[
f_\mu=f_\mu^0+f_\mu^++f_\mu^-
\qquad(\mu=1,\ldots,m+1).
\]
The vectors $f_\mu^+$ all lie in the $m$-dimensional manifold $\mathfrak H-\mathfrak E$, so they are not linearly independent. Choose $a_1,\ldots,a_{m+1}$, not all zero, such that
\[
a_1f_1^++\cdots+a_{m+1}f_{m+1}^+=0.
\]
Then
\[
a_1f_1+\cdots+a_{m+1}f_{m+1}=g^0+g^-,
\qquad
g^0\in\mathfrak A,\quad
g^-\in\mathfrak H-\mathfrak F.
\]
For $g^-=0$ this vector belongs to $\mathfrak A$; for $g^-\ne0$, it belongs to the $-$-class by \hyperref[thm:31]{Theorem 31}. This proves the first two assertions. The next two follow in the same way by interchanging $+$ and $-$.

It remains to prove the final assertion. Since the dimension of the $0$-class modulo $\mathfrak A$ is at most both $m$ and $n$, it suffices to show that it is at least $\min(m,n)$. Thus, for every finite
\[
p\le\min(m,n),
\]
we must exhibit linearly independent vectors $f_1,\ldots,f_p$ such that every nonzero combination
\[
a_1f_1+\cdots+a_pf_p
\]
belongs to the $0$-class but not to $\mathfrak A$.

There are $p$ linearly independent elements in $\mathfrak H-\mathfrak E$. We may ``orthogonalize'' them immediately (cf. \hyperref[thm:8]{Theorem 8}), and hence may assume them normalized and orthogonal:
\[
\varphi_1,\ldots,\varphi_p.
\]
Likewise, let
\[
\psi_1,\ldots,\psi_p
\]
be normalized orthogonal elements of $\mathfrak H-\mathfrak F$. Put
\[
f_1=\varphi_1+\psi_1,\ldots,f_p=\varphi_p+\psi_p.
\]
Then, unless $a_1=\cdots=a_p=0$,
\begin{align*}
a_1f_1+\cdots+a_pf_p
={}&(a_1\varphi_1+\cdots+a_p\varphi_p)\\
 &+(a_1\psi_1+\cdots+a_p\psi_p).
\end{align*}
By \hyperref[thm:30]{Theorem 30} this vector certainly does not belong to $\mathfrak A$, and in particular is nonzero. Moreover,
\[
|a_1\varphi_1+\cdots+a_p\varphi_p|^2
=|a_1\psi_1+\cdots+a_p\psi_p|^2
=|a_1|^2+\cdots+|a_p|^2.
\]
Hence it belongs to the $0$-class by \hyperref[thm:31]{Theorem 31}. Thus $f_1,\ldots,f_p$ are linearly independent and have all the required properties.
\end{proof}

\begin{namedstatement}{Theorem}{33}\label{thm:33}
The operator $R$ is maximal if and only if $m=0$ or $n=0$, and hypermaximal if and only if $m=n=0$. Equivalently, $R$ is maximal if and only if $\mathfrak E=\mathfrak H$ or $\mathfrak F=\mathfrak H$, and hypermaximal if and only if
\[
\mathfrak E=\mathfrak F=\mathfrak H.
\]
\end{namedstatement}

\begin{proof}
By \hyperref[thm:11]{Theorem 11}, maximality means that the $0$-class equals $\mathfrak A$; by \hyperref[def:9]{Definition 9}, hypermaximality means in addition that the $+$- and $-$-classes are empty. By the remark preceding \hyperref[thm:32]{Theorem 32}, the first condition means that the $0$-class is zero-dimensional, while the additional condition means that the $+$- and $-$-classes are zero-dimensional. Finally, by \hyperref[thm:32]{Theorem 32}, these conditions say
\[
\min(m,n)=0,
\]
that is, $m=0$ or $n=0$, and, respectively, $m=0$ and $n=0$.

The second formulation follows immediately from the first.
\end{proof}

The contrast between maximality and hypermaximality is now clear, as is the role of $m,n$. We add the following observation concerning \hyperref[thm:32]{Theorem 32}. If $R$ is replaced by
\[
aR+b1
\qquad(a,b\ \text{real},\ a\ne0),
\]
which is again a closed linear H.O. with the same domain as $R$, then $\mathfrak E,\mathfrak F$ change in an opaque way, and at first we cannot see what happens to $m,n$. But $\mathfrak A$, the set of extension elements, and the $0$-, $+$-, and $-$-classes plainly do not change at all when $a>0$; when $a<0$, only the last two classes are interchanged. It follows from \hyperref[thm:32]{Theorem 32} that $m,n$ likewise remain unchanged or are merely interchanged.\footnote{\label{fn:46}Let $\overline R$ be the following operator: $\overline Rf$ is defined for every extension element $f$ of $R$, and its value is the assigned vector $f^*$. The operator $\overline R$ is plainly closed and linear and is an extension of $R$. (If $R$ is hypermaximal, then plainly $\overline R=R$; otherwise $\overline R$ is not even a H.O., since $m>0$ or $n>0$, and hence the $+$- or the $-$-class is nonempty, while in these classes
\[
(\overline Rf,f)=(f^*,f)
\]
is not real.) By \hyperref[thm:28]{Theorem 28}, $\mathfrak H-\mathfrak E$ and $\mathfrak H-\mathfrak F$ are the sets of solutions of
\[
\overline Rf=if,\qquad\text{respectively}\qquad \overline Rf=-if.
\]
Thus $m$ and $n$ are the maximum numbers of linearly independent solutions of these equations.

By the preceding observation, the maximum number of linearly independent solutions of
\[
a\overline Rf+bf=if
\]
is $m$, respectively $n$, according as $a>0$, respectively $a<0$. This equation says
\[
\overline Rf=\rho f,\qquad
\rho=-\frac ba+\frac1a i.
\]
We may therefore say: let $\rho$ be nonreal. For the equation
\[
\overline Rf=\rho f
\]
the maximum number of linearly independent solutions is $m$ when $\operatorname{Im}\rho>0$, and $n$ when $\operatorname{Im}\rho<0$. (For real $\rho$, the maximum number is, as usual, the ``multiplicity of the point eigenvalue'' $\rho$. When $\rho$ is nonreal, the equation $Rf=\rho f$ must have no solution other than $f=0$, since otherwise
\[
(Rf,f)=\rho(f,f)
\]
would not be real.)

These facts are related to certain results of Carleman; cf. the citation in \hyperref[fn:source-4]{note 4}.}

\section*{VIII. The Extension Process}\label{chap:VIII}
\addcontentsline{toc}{section}{VIII. The Extension Process}

Already at the end of \hyperref[chap:VI]{Chapter VI} we obtained an overview of all possible extensions of a closed linear H.O. $R$; we can now decide when and how an extension to a maximal or hypermaximal operator is possible.

\begin{namedstatement}{Theorem}{34}\label{thm:34}
A closed linear H.O. $R$ with deficiency indices $m,n$ can be extended to one with deficiency indices $m',n'$ if and only if there exists
\[
p\in\{0,1,2,\ldots,\infty\}
\]
such that
\[
m'+p=m,\qquad n'+p=n.
\]
\end{namedstatement}

\begin{proof}
Consider the extension process of \hyperref[thm:27]{Theorem 27}, with $U,\mathfrak A,\mathfrak E,\mathfrak F$ having their usual meanings and $\mathfrak M,\mathfrak N$ as in that theorem. The manifolds
\[
\mathfrak H-\mathfrak E,\qquad
\mathfrak H-\mathfrak F
\]
are $m$- and $n$-dimensional, respectively. The question is whether they have closed linear submanifolds $\mathfrak M,\mathfrak N$ of the same dimension, say $p$, such that
\[
(\mathfrak H-\mathfrak E)-\mathfrak M,
\qquad
(\mathfrak H-\mathfrak F)-\mathfrak N
\]
are $m'$- and $n'$-dimensional, respectively.

More generally, the question is this: when does a $k$-dimensional closed linear manifold $\mathfrak K$ have a $p$-dimensional part, that is, a closed linear manifold $\mathfrak P$, such that $\mathfrak K-\mathfrak P$ has $k'$ dimensions? We shall see that this happens exactly when
\[
k=k'+p.
\]
Substituting
\[
\mathfrak K=\mathfrak H-\mathfrak E,\quad
\mathfrak P=\mathfrak M,\quad
k=m,\quad k'=m',
\]
and then
\[
\mathfrak K=\mathfrak H-\mathfrak F,\quad
\mathfrak P=\mathfrak N,\quad
k=n,\quad k'=n',
\]
gives the theorem.

The condition is necessary. Indeed, suppose $\mathfrak P$ and $\mathfrak K-\mathfrak P$ are $p$- and $k'$-dimensional, respectively. They are then spanned, as closed linear manifolds, by normalized orthogonal systems
\[
\varphi_1,\ldots,\varphi_p,
\qquad
\psi_1,\ldots,\psi_{k'},
\]
respectively. Hence
\[
\mathfrak K=\mathfrak P+(\mathfrak K-\mathfrak P)
\]
is spanned by
\[
\varphi_1,\ldots,\varphi_p,\psi_1,\ldots,\psi_{k'}.
\]
These systems are mutually orthogonal because they belong to $\mathfrak P$ and $\mathfrak K-\mathfrak P$. Thus $\mathfrak K$ has dimension $p+k'$, and therefore $k=k'+p$. Here either $k'$ or $p$ may also be $\infty$.

The condition is also sufficient. Suppose $\mathfrak K$ has dimension $k=k'+p$. It is then spanned by a normalized orthogonal system
\[
\varphi_1,\ldots,\varphi_p,\psi_1,\ldots,\psi_{k'},
\]
again allowing $k'$ or $p$ to be $\infty$. Let $\mathfrak P$ be the closed linear manifold spanned by $\varphi_1,\ldots,\varphi_p$. Then $\mathfrak K-\mathfrak P$ is spanned by $\psi_1,\ldots,\psi_{k'}$, and the proof is complete.
\end{proof}

\begin{namedstatement}{Theorem}{35}\label{thm:35}
Let $R$ be a closed linear H.O. with deficiency indices $m,n$. The operator $R$ can always be extended to a maximal H.O. If $m\ne n$, none of its extensions is hypermaximal. If $m=n$ and these numbers are finite, every maximal extension is hypermaximal; if $m=n=\infty$, then among the maximal extensions there are both hypermaximal ones and ones that are not hypermaximal.
\end{namedstatement}

\begin{proof}
By \hyperref[thm:33]{Theorem 33}, maximality means that at least one deficiency index is $0$, hypermaximality means that both are $0$, and maximality without hypermaximality means that exactly one is $0$. Applying \hyperref[thm:34]{Theorem 34} gives precisely the assertions of the theorem.
\end{proof}

\begin{namedstatement}{Corollary}{to Theorem 35}\label{cor:35}
If $R$ is not already maximal, then its maximal extensions are possible in continuum many different ways; and, whenever such extensions exist, the same is true separately of its hypermaximal extensions and of its maximal but nonhypermaximal extensions.
\end{namedstatement}

\begin{proof}
By \hyperref[thm:27]{Theorem 27}, our extension method amounted to enlarging $\mathfrak E$ and $\mathfrak F$ by two closed linear manifolds $\mathfrak M,\mathfrak N$, which, when $R$ is not maximal, had positive dimension.\translatornote{The source prints $\mathfrak H-\mathfrak E$ and $\mathfrak H-\mathfrak F$ here; Theorem 27 shows that $\mathfrak E$ and $\mathfrak F$ are meant.} But a length-preserving mapping $U$ from $\mathfrak M$ onto $\mathfrak N$ also enters, and it can be chosen in continuum many ways: by \hyperref[thm:26]{Theorem 26}, every normalized orthogonal system
\[
\psi_1,\ldots,\psi_p
\]
spanning $\mathfrak N$ gives one such $U$, and there are continuum many of these systems, since, for example, $\psi_1$ may be replaced by any $\alpha\psi_1$ with $|\alpha|=1$.\footnote{\label{fn:47}There cannot be more than continuum many extensions, since there are only continuum many closed operators in $\mathfrak H$, as follows easily from the separability of $\mathfrak H$.}
\end{proof}

If $R$ is merely an H.O., then, as we know, its maximal and hypermaximal extensions are the same as those of the closed linear H.O. $\widetilde R$, to which the results just obtained are applicable. In particular, if $R$ has only one maximal extension, then $\widetilde R$ itself must already be maximal; the same holds for hypermaximal extensions.

We now turn to a closer investigation of hypermaximal H.O. (\hyperref[chap:IX]{Chapter IX}), and then of maximal but nonhypermaximal H.O. (\hyperref[chap:X]{Chapter X}).

\section*{IX. Hypermaximal Operators and the Eigenvalue Representation}\label{chap:IX}
\addcontentsline{toc}{section}{IX. Hypermaximal Operators and the Eigenvalue Representation}

If $R$ is hypermaximal, then
\[
\mathfrak E=\mathfrak F=\mathfrak H;
\]
that is, its Cayley transform $U$ is unitary: it is meaningful everywhere and has an inverse $U^{-1}$ that is likewise unitary and meaningful everywhere. Thus the hypermaximal H.O. $R$ correspond one-to-one to the unitary operators $U$ satisfying the condition of \hyperref[thm:24]{Theorem 24}, namely, those for which the set of all $\varphi-U\varphi$ is everywhere dense.

Already in the proof of \hyperref[thm:24]{Theorem 24} we observed that, if the $\varphi-U\varphi$ are everywhere dense, then $U\varphi=\varphi$ implies $\varphi=0$. Here the converse also holds. For if the $\varphi-U\varphi$ are not everywhere dense, then they span a closed linear manifold $\mathfrak M\ne\mathfrak H$ (they themselves already form a linear manifold), and hence $\mathfrak H-\mathfrak M\ne(0)$. Let $f\ne0$ be an element of $\mathfrak H-\mathfrak M$. Every $\varphi-U\varphi$ belongs to $\mathfrak M$ and is therefore orthogonal to $f$:
\begin{align*}
(f,\varphi-U\varphi)
&=(f,\varphi)-(f,U\varphi)\\
&=(Uf,U\varphi)-(f,U\varphi)
=(Uf-f,U\varphi)
\end{align*}
must vanish for every $\varphi$; consequently $Uf-f=0$, or $Uf=f$.

The unitary operators in question can therefore also be characterized as those for which $Uf=f$ has no solution other than $0$ (the number $1$ is not a ``point eigenvalue''). Thus we control the hypermaximal H.O. $R$ once we know these $U$. This, however, is a problem concerning bounded operators ($U$ preserves length and is therefore continuous), and it can be solved by the methods of Hilbert's theory of bounded forms, with an addition that also belongs to that circle of ideas. We carry out the necessary considerations, partly following F.~Riesz's simple method, in \hyperref[app:II]{Appendix II}, and state only the result here.

\begin{namedstatement}{Definition}{17}\label{def:17}
Let $a<x<b$ be an open interval (below, $a,b$ will once be $0,1$ and once be $-\infty,\infty$). By a \emph{resolution of the identity} (abbreviated: r.i.) we mean a family of P.O.'s $E(\lambda)$, with $\lambda$ ranging over all of $(a,b)$, having the following properties:

\begin{enumerate}
\renewcommand{\labelenumi}{$\alpha)$}
\item\label{def17:alpha}
If $\lambda\le\mu$, then $E(\lambda)$ is a part of $E(\mu)$.

\renewcommand{\labelenumi}{$\beta)$}
\item\label{def17:beta}
If $\lambda\ge\lambda_0$ and $\lambda\to\lambda_0$, then, for every $f$,
\[
E(\lambda)f\longrightarrow E(\lambda_0)f.
\]

\renewcommand{\labelenumi}{$\gamma)$}
\item\label{def17:gamma}
If $\lambda\to a$, respectively $\lambda\to b$, then
\[
E(\lambda)f\longrightarrow0,
\qquad\text{respectively}\qquad
E(\lambda)f\longrightarrow f.
\]
\end{enumerate}
\end{namedstatement}

If now $(a,b)=(0,1)$ and $E(\rho)$ is an r.i., then for every $f$ there is exactly one $f^*$ such that, for every $g$,
\[
(f^*,g)=\int_0^1e^{2\pi i\rho}\,d\bigl(E(\rho)f,g\bigr)
\]
holds.\footnote{\label{fn:48}This is a Stieltjes integral, and it is absolutely convergent. We write $\rho$ in place of $\lambda$.} The operator $U$ defined by $Uf=f^*$ is unitary, and $Uf=f$ has no solution other than $0$. Conversely, every $U$ with these properties is produced in exactly one way by an r.i. $E(\rho)$ through the procedure just sketched (cf., as stated, \hyperref[app:II]{Appendix II}).

We now prove:

\begin{namedstatement}{Theorem}{36}\label{thm:36}
Let $(a,b)=(-\infty,\infty)$, and let $F(\lambda)$ be an r.i. If the integral
\[
\int_{-\infty}^{\infty}\lambda^2\,d\,|F(\lambda)f|^2
\]
is finite,\footnote{\label{fn:49}As $\lambda$ increases, $F(\lambda)$ increases, and hence so does $|F(\lambda)f|^2$. Moreover, $\lambda^2$ is never negative. By its nature, therefore, this integral either converges and is finite or diverges properly to $+\infty$.} then there is exactly one $f^*$ such that, for every $g$,
\[
(f^*,g)
=\int_{-\infty}^{\infty}\lambda\,d\bigl(F(\lambda)f,g\bigr),
\]
where the integral on the right is absolutely convergent. We define an operator $R$ for these $f$ by $Rf=f^*$.

This $R$ is a hypermaximal H.O., and every hypermaximal H.O. $R$ can be produced in this way by means of exactly one r.i.
\end{namedstatement}

\begin{proof}
The r.i. $E(\rho)$ for the interval $(0,1)$ and the r.i. $F(\lambda)$ for the interval $(-\infty,\infty)$ correspond one-to-one by the following substitutions.\footnote{\label{fn:50}The substitutions
\[
\rho=-\frac1\pi\operatorname{arccot}\lambda,
\qquad
\lambda=-\cot\pi\rho
\]
put the intervals $0<\rho<1$ and $-\infty<\lambda<\infty$ into one-to-one correspondence. For $\operatorname{arccot}$, the value between $-\pi$ and $0$ is always to be taken.}
\begin{align*}
E(\rho)&\longmapsto
F(\lambda)=E\left(-\frac1\pi\operatorname{arccot}\lambda\right),\\
F(\lambda)&\longmapsto
E(\rho)=F(-\cot\pi\rho).
\end{align*}
It therefore suffices to show the following. If $E(\rho)$ is the r.i. of a unitary operator $U$ (cf.\ the discussion preceding the theorem), then $F(\lambda)$ really does produce an operator $R$ by the prescription of the theorem, and this operator is the Cayley transform of $U$. Everything we know about the relation between Cayley transforms, and everything stated before the theorem about hypermaximal H.O., unitary operators, and r.i., then proves all the assertions.

Let, therefore, $F(\lambda)$ be an r.i., let
\[
E(\rho)=F(-\cot\pi\rho),
\]
and let $U$ be the corresponding unitary operator. We first show that, when
\[
\int_{-\infty}^{\infty}\lambda^2\,d\,|F(\lambda)f|^2=C^2
\]
is finite, the $f^*$ of the theorem can actually be constructed.

Let $\lambda<\lambda'<\mu$. Since $F(\lambda)$ is a part of $F(\mu)$, $F(\mu)-F(\lambda)$ is a P.O., and therefore
\begin{align*}
\bigl|\bigl((F(\mu)-F(\lambda))f,g\bigr)\bigr|
&=\bigl|\bigl((F(\mu)-F(\lambda))f,(F(\mu)-F(\lambda))g\bigr)\bigr|\\
&\le
\bigl|(F(\mu)-F(\lambda))f\bigr|\,
\bigl|(F(\mu)-F(\lambda))g\bigr|.
\end{align*}
Moreover, in general,
\begin{align*}
\bigl|(F(\mu)-F(\lambda))h\bigr|^2
&=\bigl(h,(F(\mu)-F(\lambda))h\bigr)\\
&=(h,F(\mu)h)-(h,F(\lambda)h)\\
&=|F(\mu)h|^2-|F(\lambda)h|^2.
\end{align*}
Consequently,\translatornote{The source writes $E(\lambda)$ rather than $F(\lambda)$ throughout this intermediate estimate. We use $F$ consistently with the statement of the theorem and with the surrounding proof.}
\begin{align*}
\bigl|\bigl((F(\mu)-F(\lambda))f,g\bigr)\bigr|
&\le
\sqrt{|F(\mu)f|^2-|F(\lambda)f|^2}\,
\sqrt{|F(\mu)g|^2-|F(\lambda)g|^2},\\
\left|\lambda'
\bigl\{(F(\mu)f,g)-(F(\lambda)f,g)\bigr\}\right|
&\le
\sqrt{\lambda'^2\bigl\{|F(\mu)f|^2-|F(\lambda)f|^2\bigr\}}\,
\sqrt{|F(\mu)g|^2-|F(\lambda)g|^2}.
\end{align*}
Now the two integrals
\[
\int_{-\infty}^{\infty}\lambda^2\,d\,|F(\lambda)f|^2=C^2,
\qquad
\int_{-\infty}^{\infty}d\,|F(\lambda)g|^2
=
\left.|F(\lambda)g|^2\right|_{\lambda=-\infty}^{\lambda=\infty}
=|g|^2
\]
converge absolutely. Hence the familiar inequality
\[
\sqrt{a_1b_1}+\cdots+\sqrt{a_pb_p}
\le
\sqrt{(a_1+\cdots+a_p)(b_1+\cdots+b_p)}
\qquad
(a_1,\ldots,a_p,b_1,\ldots,b_p\ge0)
\]
has, first, the consequence that
\[
\int_{-\infty}^{\infty}\lambda\,d\bigl(F(\lambda)f,g\bigr)
\]
also converges absolutely, and, second, that\footnote{\label{fn:51}It suffices to recall the definition of the Stieltjes integral. The inequality just used is merely a modified form of Schwarz's inequality.}
\[
\left|
\int_{-\infty}^{\infty}\lambda\,d\bigl(F(\lambda)f,g\bigr)
\right|
\le
\sqrt{
\left(\int_{-\infty}^{\infty}\lambda^2\,d\,|F(\lambda)f|^2\right)
\left(\int_{-\infty}^{\infty}d\,|F(\lambda)g|^2\right)}
=C|g|.
\]

For a moment, with $f$ fixed and $g$ variable, denote this integral by $L(g)$. Then $L(g)$ has the properties
\begin{align*}
L(a_1g_1+\cdots+a_ng_n)
&=\overline{a_1}L(g_1)+\cdots+\overline{a_n}L(g_n),\\
|L(g)|&\le C|g|;
\end{align*}
it is conjugate-linear and continuous. By a theorem of F.~Riesz, there is therefore exactly one $f^*$ such that
\[
(f^*,g)=L(g)
\]
always holds.\footnote{\label{fn:52}The theorem is proved as follows. Since all the $(f^*,g)$ are prescribed, there cannot be two distinct $f^*$, so it suffices to prove existence. Let $\varphi_1,\varphi_2,\ldots$ be a complete orthonormal system, and put
\[
L(\varphi_n)=a_n
\qquad(n=1,2,\ldots).
\]
By \hyperref[thm:7]{Theorem 7}, $\beta)$,
\[
g=\sum_{n=1}^{\infty}(g,\varphi_n)\varphi_n;
\]
hence, by the conjugate-linearity and continuity of $L$,
\[
L(g)=\sum_{n=1}^{\infty}\overline{(g,\varphi_n)}\,a_n,
\]
and this series must converge. If $\sum_{n=1}^{\infty}|a_n|^2$ is finite, then
\[
f^*=\sum_{n=1}^{\infty}a_n\varphi_n
\]
converges by \hyperref[thm:5]{Theorem 5}, and $(f^*,\varphi_n)=a_n$. It then follows from \hyperref[thm:7]{Theorem 7}, $\gamma)$, that
\[
L(g)
=\sum_{n=1}^{\infty}(f^*,\varphi_n)\overline{(g,\varphi_n)}
=(f^*,g),
\]
which is the assertion.

It remains to prove that $\sum_{n=1}^{\infty}|a_n|^2$ is finite. Since $|L(g)|\le C|g|$, taking
\[
g=x_1\varphi_1+\cdots+x_m\varphi_m
\]
gives
\[
\left|\sum_{n=1}^{m}a_n\overline{x_n}\right|
\le C\sqrt{\sum_{n=1}^{m}|x_n|^2}.
\]
Set $x_n=a_n$ for $n=1,\ldots,m$. Then
\[
\sum_{n=1}^{m}|a_n|^2\le C^2.
\]
Thus $\sum_{n=1}^{\infty}|a_n|^2$ is finite, and is at most $C^2$.}
This is precisely the assertion concerning $f^*$.

The operator $R$ can therefore actually be constructed. Let $R^*$ be the Cayley transform of $U$; it remains to prove that $R=R^*$. It is enough to show that $R$ is an extension of $R^*$. For whenever $Rf$ and $Rg$ have meaning, interchanging $f$ and $g$ changes $(Rf,g)$ into its complex conjugate, by the definition in \hyperref[thm:36]{Theorem 36}; hence
\[
(Rf,g)=(f,Rg).
\]
Thus $R$ is an H.O. (its domain contains, by assumption, that of $R^*$ and is therefore everywhere dense) and an extension of the maximal $R^*$, which is maximal by the unitarity of $U$; consequently $R=R^*$.

Suppose, then, that $R^*f$ has meaning. We must prove that $Rf$ has meaning and equals $R^*f$; equivalently, that
\[
\int_{-\infty}^{\infty}\lambda^2\,d\,|F(\lambda)f|^2
\]
is finite and that, for every $g$,
\[
(R^*f,g)=(Rf,g)
=\int_{-\infty}^{\infty}\lambda\,d\bigl(F(\lambda)f,g\bigr).
\]
The fact that $R^*f$ has meaning says, moreover, that
\[
f=\varphi-U\varphi;
\]
this will be our starting point.

Under the transformation $\lambda=-\cot\pi\rho$, the two integrals
\[
\int_{-\infty}^{\infty}\lambda^2\,d\,|F(\lambda)f|^2,
\qquad
\int_{-\infty}^{\infty}\lambda\,d\bigl(F(\lambda)f,g\bigr)
\]
may be written as
\[
\int_0^1\cot^2\pi\rho\,d\,|E(\rho)f|^2,
\qquad
\int_0^1-\cot\pi\rho\,d\bigl(E(\rho)f,g\bigr).
\]
To substitute $f=\varphi-U\varphi$ here, we must first calculate several integrals:
\begin{align*}
(\varphi\pm U\varphi,g)
&=\int_0^1d\bigl(E(\rho)\varphi,g\bigr)
\pm\int_0^1e^{2\pi i\rho}\,d\bigl(E(\rho)\varphi,g\bigr)\\
&=\int_0^1(1\pm e^{2\pi i\rho})\,d\bigl(E(\rho)\varphi,g\bigr).
\end{align*}

It follows that\footnote{\label{fn:53}For $\rho'>\rho$, the expression
\[
\bigl(E(\min(\rho,\rho'))\varphi,g\bigr)
\]
is constant, and hence its Stieltjes differential, and therefore the corresponding integral, is $0$.}
\begin{align*}
(\varphi-U\varphi,E(\rho)g)
&=\int_0^1(1-e^{2\pi i\rho'})\,
d\bigl(E(\rho')\varphi,E(\rho)g\bigr)\\
&=\int_0^1(1-e^{2\pi i\rho'})\,
d\bigl(E(\rho)E(\rho')\varphi,g\bigr)\\
&=\int_0^1(1-e^{2\pi i\rho'})\,
d\bigl(E(\min(\rho,\rho'))\varphi,g\bigr)\\
&=\int_0^\rho(1-e^{2\pi i\rho'})\,
d\bigl(E(\rho')\varphi,g\bigr).
\end{align*}
Therefore, using complex conjugation in the second line\footnote{\label{fn:54}Take the complex conjugate of the preceding formula with $g=\varphi$.} and the general Stieltjes formula in the next step\footnote{\label{fn:55}The general formula for Stieltjes integrals is
\[
\int_A^Bp(\rho)\,
d\left[\int_A^\rho q(\rho')\,dr(\rho')\right]
=
\int_A^Bp(\rho)q(\rho)\,dr(\rho).
\]}, we obtain
\begin{align*}
|E(\rho)(\varphi-U\varphi)|^2
&=(\varphi-U\varphi,E(\rho)(\varphi-U\varphi))\\
&=\int_0^\rho(1-e^{2\pi i\rho'})\,
d\bigl(E(\rho')\varphi,\varphi-U\varphi\bigr)\\
&=\int_0^\rho(1-e^{2\pi i\rho'})\,
d\left[
\int_0^{\rho'}(1-e^{-2\pi i\rho''})\,
d\bigl(E(\rho'')\varphi,\varphi\bigr)
\right]\\
&=\int_0^\rho
(1-e^{2\pi i\rho'})(1-e^{-2\pi i\rho'})\,
d\bigl(E(\rho')\varphi,\varphi\bigr)\\
&=\int_0^\rho4\sin^2(\pi\rho')\,d\,|E(\rho')\varphi|^2.
\end{align*}
It follows, using the formula of note \ref{fn:55} once more, that
\begin{align*}
\int_0^1\cot^2(\pi\rho)\,
d\,|E(\rho)(\varphi-U\varphi)|^2
&=
\int_0^1\cot^2(\pi\rho)\,
d\left[
\int_0^\rho4\sin^2(\pi\rho')\,d\,|E(\rho')\varphi|^2
\right]\\
&=
\int_0^1\cot^2(\pi\rho)\,4\sin^2(\pi\rho)\,
d\,|E(\rho)\varphi|^2\\
&=
\int_0^14\cos^2(\pi\rho)\,d\,|E(\rho)\varphi|^2.
\end{align*}
This is finite, since it is bounded above by the finite quantity
\[
\int_0^14\,d\,|E(\rho)\varphi|^2=4|\varphi|^2.
\]
Furthermore,
\begin{align*}
\int_0^1\!-\cot(\pi\rho)\,
d\bigl(E(\rho)(\varphi-U\varphi),g\bigr)
&=
\int_0^1-\cot(\pi\rho)\,
d\left[
\int_0^\rho(1-e^{2\pi i\rho'})\,
d\bigl(E(\rho')\varphi,g\bigr)
\right]\\
&=
\int_0^1\!-\cot(\pi\rho)(1-e^{2\pi i\rho})\,
d\bigl(E(\rho)\varphi,g\bigr)\\
&=
\int_0^1i(1+e^{2\pi i\rho})\,
d\bigl(E(\rho)\varphi,g\bigr)\\
&=(i(\varphi+U\varphi),g)\\
&=(R^*(\varphi-U\varphi),g).
\end{align*}
This proves everything required.
\end{proof}

We have obtained for hypermaximal H.O. the normal form analogous to Hilbert's spectral representation or eigenvalue representation for bounded forms (cf.\ the citation in note \ref{fn:source-7}). In doing so, however, we have already exhausted all eigenvalue representations for these operators, so there certainly can be none for maximal but nonhypermaximal H.O. (Admittedly, we have not yet constructed an example of such an H.O.) We therefore refrain from further investigation of hypermaximal H.O., which would offer scarcely any new points of view, and turn to maximal but nonhypermaximal H.O.

\section*{X. Maximal but Nonhypermaximal Operators}\label{chap:X}
\addcontentsline{toc}{section}{X. Maximal but Nonhypermaximal Operators}

Such an H.O. $R$, with deficiency indices $m,n$, is characterized by either
\[
m=0,\quad n>0,
\qquad\text{or}\qquad
m>0,\quad n=0.
\]
Since replacing $R$ by $-R$ interchanges $m,n$, it is enough to consider $m=0,n>0$. Form $U,\mathfrak E,\mathfrak F$. Then
\[
\mathfrak E=\mathfrak H,
\qquad
\mathfrak G_0=\mathfrak H-\mathfrak F
\]
is $n$-dimensional, and $U$ maps $\mathfrak H$ isometrically onto $\mathfrak F$.

Let $\mathfrak G_\nu$ be the image of $\mathfrak G_0$ under $U^\nu$ ($\nu=0,1,2,\ldots$); it too is a closed linear manifold. For $\nu>0$, $\mathfrak G_\nu$ is a part of $\mathfrak F$ and is therefore orthogonal to $\mathfrak G_0$. Applying the mapping $U^k$ ($k=0,1,2,\ldots$), the images $\mathfrak G_{k+\nu}$ and $\mathfrak G_k$ are orthogonal, since $(f,g)$ is invariant under this mapping; cf.\ \hyperref[thm:10]{Theorem 10}. Thus, in general,
\[
\mathfrak G_k\perp\mathfrak G_l
\qquad(k,l=0,1,2,\ldots;\ k\ne l).
\]
Let $\mathfrak G_0,\mathfrak G_1,\mathfrak G_2,\ldots$ span the closed linear manifold $\mathfrak M$, and put
\[
\mathfrak H-\mathfrak M=\mathfrak N.
\]

We already know that $U$ maps $\mathfrak G_\nu$ one-to-one and isometrically onto $\mathfrak G_{\nu+1}$. We now show that it likewise maps $\mathfrak N$ onto itself. As is easily seen, $\mathfrak N$ is the set of all elements of $\mathfrak H$ orthogonal to every $\mathfrak G_\nu$. Because $(f,g)$ is invariant, its image is the set of all elements in the image of $\mathfrak H$ that are orthogonal to every image of the $\mathfrak G_\nu$. These are the elements of
\[
\mathfrak F=\mathfrak H-\mathfrak G_0
\]
orthogonal to every $\mathfrak G_{\nu+1}$, hence the elements of $\mathfrak H$ orthogonal to $\mathfrak G_0$ and to all $\mathfrak G_{\nu+1}$. This is indeed $\mathfrak N$ again.

Let $\varphi_1^0,\ldots,\varphi_p^0$ be an orthonormal system spanning the closed linear manifold $\mathfrak G_0$, where possibly $p=\infty$, and put
\[
U^\nu\varphi_q^0=\varphi_q^\nu
\qquad
(\nu=0,1,2,\ldots;\ q=1,\ldots,p).
\]
Then $\varphi_1^\nu,\ldots,\varphi_p^\nu$ is also an orthonormal system and spans $\mathfrak G_\nu$. For $k\ne l$, $\varphi_{q'}^k$ lies in $\mathfrak G_k$ and $\varphi_{q''}^l$ in $\mathfrak G_l$, and the two are therefore orthogonal. Thus all
\[
\varphi_q^\nu
\qquad
(\nu=0,1,2,\ldots;\ q=1,\ldots,p)
\]
form an orthonormal system; moreover, they span the same closed linear manifold as the $\mathfrak G_\nu$, namely $\mathfrak M$.

For each $q=1,\ldots,p$, let
\[
\varphi_q^0,\varphi_q^1,\varphi_q^2,\ldots
\]
span the closed linear manifold $\mathfrak M_q$. Then $\mathfrak M_1,\ldots,\mathfrak M_p$ span $\mathfrak M$, and $\mathfrak M_q$ is a part of $\mathfrak M$. Since
\[
U\varphi_q^\nu=\varphi_q^{\nu+1},
\]
$U$ maps $\mathfrak M_q$ into itself. In summary:

\[
\mathfrak M_1,\ldots,\mathfrak M_p,\mathfrak N
\]
span $\mathfrak H$, just as $\mathfrak M,\mathfrak N$ do, and $U$ maps them respectively into $\mathfrak M_1,\ldots,\mathfrak M_p,\mathfrak N$, that is, into subsets of themselves. Moreover, $\mathfrak M_1,\ldots,\mathfrak M_p$, as parts of $\mathfrak M$, are orthogonal to $\mathfrak N$ and also to one another, because their spanning vectors $\varphi_{q'}^k,\varphi_{q''}^l$ have $q'\ne q''$. It follows that each of
\[
\mathfrak M_1,\ldots,\mathfrak M_p,\mathfrak N
\]
reduces $U$, and therefore, by \hyperref[thm:25]{Theorem 25}, also its Cayley transform $R$.

Consider the parts of $R$ and $U$ lying in $\mathfrak M_1,\ldots,\mathfrak M_p,\mathfrak N$. Denote them, respectively, by
\[
R_1,\ldots,R_p,S,
\qquad\text{and}\qquad
U_1,\ldots,U_p,V.
\]
By \hyperref[thm:21]{Theorem 21}, $R$ is characterized by $R_1,\ldots,R_p,S$. Now all the closed linear manifolds $\mathfrak M_1,\ldots,\mathfrak M_p$ are infinite-dimensional, since $\mathfrak M_q$ is spanned by the orthonormal system
\[
\varphi_q^0,\varphi_q^1,\varphi_q^2,\ldots.
\]
Thus they are all Hilbert spaces, meaning that they satisfy conditions A through E of \hyperref[chap:I]{Chapter I}, with the definitions of $af,f+g,(f,g)$ unchanged,\footnote{\label{fn:56}Every closed linear manifold plainly satisfies A through C and E, while D is equivalent to infinite dimensionality. A finite-dimensional closed linear manifold consists only of $0$ when it has dimension $0$. If it has $N=1,2,\ldots$ dimensions, it is spanned by an orthonormal system $\varphi_1,\ldots,\varphi_N$ and therefore consists of all
\[
x_1\varphi_1+\cdots+x_N\varphi_N;
\]
that is, it is the $N$-dimensional complex Euclidean space of all $x_1,\ldots,x_N$.} and hence all the concepts developed for $\mathfrak H$ may be used in them without qualification. Thus $R_q$ is an H.O. in $\mathfrak M_q$. In particular, its domain is everywhere dense in $\mathfrak M_q$, since it consists of the projections of the everywhere-dense domain of $R$ in $\mathfrak H$. The operator $U_q$ preserves length and is the Cayley transform of $R_q$.

Three cases are possible in $\mathfrak N$. First, $\mathfrak N$ may be infinite-dimensional, and hence a Hilbert space; then $S$ is again an H.O., $V$ is unitary, since it maps $\mathfrak N$ onto $\mathfrak N$, and $S$ is its Cayley transform, hence $S$ is hypermaximal. Second, $\mathfrak N$ may be $N$-dimensional, where $N=1,2,\ldots$, and hence an $N$-dimensional complex Euclidean space (cf.\ note \ref{fn:56}); then $S$ is an H.O. in it, that is, an ordinary finite-dimensional Hermitian operator. Third, $\mathfrak N$ may be zero-dimensional, that is, $\mathfrak N=(0)$; then it need not be considered at all.

Thus nothing new can occur in $\mathfrak N$ and $S$, and we therefore turn to $\mathfrak M_q,R_q,U_q$ ($q=1,\ldots,p$). They all behave alike: $R_q$ is the Cayley transform of $U_q$, and in $\mathfrak M_q$ there is a complete orthonormal system
\[
\varphi_q^0,\varphi_q^1,\varphi_q^2,\ldots
\]
such that
\[
U_q\varphi_q^\nu=U\varphi_q^\nu=\varphi_q^{\nu+1}.
\]
Thus
\[
U_q\left(\sum_{\nu=0}^{\infty}x_\nu\varphi_q^\nu\right)
=
\sum_{\nu=0}^{\infty}x_\nu\varphi_q^{\nu+1},
\qquad
\sum_{\nu=0}^{\infty}|x_\nu|^2<\infty.
\]
This completely determines $U_q$, and hence $R_q$, once $\varphi_q^0,\varphi_q^1,\varphi_q^2,\ldots$ are regarded as known. Using the isomorphism of $\mathfrak M_q$ with $\mathfrak H$, we may also say that $R_q$ must be isomorphic to the following operator $\overline R$ in $\mathfrak H$: $\overline R$ is the Cayley transform of $\overline U$, where, for a fixed complete orthonormal system $\psi_0,\psi_1,\psi_2,\ldots$ in $\mathfrak H$,
\[
\overline U\left(\sum_{\nu=0}^{\infty}x_\nu\psi_\nu\right)
=
\sum_{\nu=0}^{\infty}x_\nu\psi_{\nu+1},
\qquad
\sum_{\nu=0}^{\infty}|x_\nu|^2<\infty.
\]

It remains first to show that $\overline R$ really exists, that is, that the evidently length-preserving $\overline U$ satisfies the condition of \hyperref[thm:24]{Theorem 24}. If it does, then, since $\overline U$ is defined on all of $\mathfrak H$, so that $\mathfrak E=\mathfrak H$ and $m=0$, $\overline R$ is a maximal H.O. The range $\mathfrak F$ of $\overline U$ is the set of all
\[
\sum_{\nu=0}^{\infty}x_\nu\psi_{\nu+1},
\]
that is, the closed linear manifold spanned by $\psi_1,\psi_2,\ldots$. Hence $\mathfrak H-\mathfrak F$ is the set of all $a\psi_0$, and is one-dimensional.\translatornote{The source prints $a\psi_1$ here; the orthogonal complement of the span of $\psi_1,\psi_2,\ldots$ is plainly the span of $\psi_0$.} Consequently $\overline R$ has deficiency indices $0,1$; it is the first example of a maximal but nonhypermaximal H.O.

It remains to prove its existence, that is, the everywhere-denseness of the set of all $\varphi-\overline U\varphi$. Since this set is a linear manifold, it suffices to show that all elements of a complete orthonormal system are its accumulation points, for example all $\psi_\nu$. Now
\begin{align*}
&\left(
\psi_\nu+\frac{l-1}{l}\psi_{\nu+1}
\,+\cdots+\frac1l\psi_{\nu+l-1}
\right)
-\overline U\left(
\psi_\nu+\frac{l-1}{l}\psi_{\nu+1}
\,+\cdots+\frac1l\psi_{\nu+l-1}
\right)\\
&\quad=
\left(
\psi_\nu+\frac{l-1}{l}\psi_{\nu+1}
\,+\cdots+\frac1l\psi_{\nu+l-1}
\right)
-
\left(
\psi_{\nu+1}+\frac{l-1}{l}\psi_{\nu+2}
\,+\cdots+\frac1l\psi_{\nu+l}
\right)\\
&\quad=
\psi_\nu-\frac1l(\psi_{\nu+1}+\cdots+\psi_{\nu+l}).
\end{align*}
This differs from $\psi_\nu$ by
\[
\frac1l(\psi_{\nu+1}+\cdots+\psi_{\nu+l}),
\]
whose squared norm and norm are, respectively,
\[
\left|\frac1l(\psi_{\nu+1}+\cdots+\psi_{\nu+l})\right|^2
=\frac1l,
\qquad
\left|\frac1l(\psi_{\nu+1}+\cdots+\psi_{\nu+l})\right|
=\frac1{\sqrt l}.
\]
With a suitable $l$, this is arbitrarily small. Everything is therefore proved.

We formulate the result:

\begin{namedstatement}{Theorem}{37}\label{thm:37}
Let $\psi_0,\psi_1,\psi_2,\ldots$ be a complete orthonormal system. There is then exactly one H.O. $\overline R$ whose Cayley transform $\overline U$ is defined by
\[
\overline U\left(\sum_{\nu=0}^{\infty}x_\nu\psi_\nu\right)
=
\sum_{\nu=0}^{\infty}x_\nu\psi_{\nu+1},
\qquad
\sum_{\nu=0}^{\infty}|x_\nu|^2<\infty.
\]
The operator $\overline R$ has deficiency indices $0,1$ and is therefore maximal but not hypermaximal.
\end{namedstatement}

\begin{proof}
This was just proved.
\end{proof}

Below we shall derive several further important general properties of $\overline R$; Part VI of the ``Remarks'' is devoted to its explicit representation.\translatornote{The ``Remarks'' cited here are not included in the present article, so there is no internal cross-reference target.} First we return to our earlier developments and prove:

\begin{namedstatement}{Theorem}{38}\label{thm:38}
Let $R$ be a maximal but nonhypermaximal H.O.; thus its deficiency indices are either
\[
0,n\quad(n=1,2,\ldots,\infty)
\]
or
\[
m,0\quad(m=1,2,\ldots,\infty).
\]
Put $p=n$, respectively $p=m$. Then $\mathfrak H$ can be decomposed into $p+1$ closed linear manifolds
\[
\mathfrak M_1,\ldots,\mathfrak M_p,\mathfrak N,
\]
which are pairwise orthogonal and together span the closed linear manifold $\mathfrak H$; when $p=\infty$, the sequence $\mathfrak M_1,\mathfrak M_2,\ldots$ does not terminate. They have the following properties:

\begin{enumerate}
\renewcommand{\labelenumi}{$\alpha)$}
\item\label{thm38:alpha}
Each of $\mathfrak M_1,\ldots,\mathfrak M_p,\mathfrak N$ reduces $R$; denote the parts of $R$ lying in them by
\[
R_1,\ldots,R_p,S,
\]
respectively.

\renewcommand{\labelenumi}{$\beta)$}
\item\label{thm38:beta}
Each $\mathfrak M_q$ ($q=1,2,\ldots,p$) is infinite-dimensional and is a Hilbert space, that is, isomorphic to $\mathfrak H$. In the case of deficiency indices $0,n$, all $R_q$ are isomorphic to $\overline R$; in the case of deficiency indices $m,0$, all $R_q$ are isomorphic to $-\overline R$, where $\overline R$ is the operator of \hyperref[thm:37]{Theorem 37}.

\renewcommand{\labelenumi}{$\gamma)$}
\item\label{thm38:gamma}
There are three possibilities for $\mathfrak N$. First, it may be infinite-dimensional and a Hilbert space, hence isomorphic to $\mathfrak H$; then $S$ is a hypermaximal H.O. Second, it may be $N$-dimensional, where $N=1,2,\ldots$, and hence an $N$-dimensional complex Euclidean space; then $S$ is an ordinary finite-dimensional Hermitian operator. Third, it may be zero-dimensional, that is, $\mathfrak N=(0)$; then it and $S$ may be omitted.
\end{enumerate}

These $R_1,\ldots,R_p,S$ determine $R$ by \hyperref[thm:21]{Theorem 21}.
\end{namedstatement}

\begin{proof}
For $m=0,n>0$ this was proved in the first part of this chapter; one need only note \hyperref[thm:37]{Theorem 37}. For $m>0,n=0$, it follows by replacing $R$ by $-R$, thereby interchanging $m,n$.
\end{proof}

In finite-dimensional spaces every Hermitian operator can be put into diagonal form. As long as the space has more than one dimension, the operator is therefore reducible. In $\mathfrak H$, hypermaximal H.O. are likewise reducible, as follows.

Let $R$ be hypermaximal, and let $F(\lambda)$ be its r.i. from \hyperref[thm:36]{Theorem 36}. Suppose first that there is a $\lambda$ such that
\[
F(\lambda)\ne0,1.
\]
Then the closed linear manifold $\mathfrak M$ whose P.O. is
\[
P_{\mathfrak M}=F(\lambda)
\]
is neither $(0)$ nor $\mathfrak H$. This $\mathfrak M$ reduces $R$. First,
\begin{align*}
\int_{-\infty}^{\infty}\lambda'^2\,
d\,|F(\lambda')F(\lambda)f|^2
&=
\int_{-\infty}^{\infty}\lambda'^2\,
d\,|F(\min(\lambda',\lambda))f|^2\\
&=
\int_{-\infty}^{\lambda}\lambda'^2\,
d\,|F(\lambda')f|^2,
\end{align*}
by the observation of note \ref{fn:53}; and this is finite whenever
\[
\int_{-\infty}^{\infty}\lambda'^2\,d\,|F(\lambda')f|^2
\]
is finite. Second, in that case,\translatornote{From these displays through the following scalar case, the source switches from the previously declared $F$ back to $E$. We retain $F$ consistently.}
\begin{align*}
(F(\lambda)Rf,g)
&=(Rf,F(\lambda)g)\\
&=\int_{-\infty}^{\infty}\lambda'\,
d\bigl(F(\lambda')f,F(\lambda)g\bigr)\\
&=\int_{-\infty}^{\infty}\lambda'\,
d\bigl(F(\lambda)F(\lambda')f,g\bigr)\\
&=\int_{-\infty}^{\infty}\lambda'\,
d\bigl(F(\lambda')F(\lambda)f,g\bigr)\\
&=(RF(\lambda)f,g)
\end{align*}
for every $g$. Thus
\[
F(\lambda)Rf=RF(\lambda)f.
\]
Suppose, on the other hand, that every $F(\lambda)$ equals either $0$ or $1$. Then, by $\alpha)$ of \hyperref[def:17]{Definition 17}, there is a $\lambda_0$ such that the former holds for $\lambda<\lambda_0$ and the latter for $\lambda>\lambda_0$. The number $\lambda_0$ is finite by $\gamma)$, and at $\lambda=\lambda_0$ the latter holds by $\beta)$.\translatornote{The source interchanges the references to $\beta)$ and $\gamma)$ in this sentence; the references are corrected here.} Therefore
\[
(Rf,g)
=\int_{-\infty}^{\infty}\lambda'\,
d\bigl(F(\lambda')f,g\bigr)
=\lambda_0(f,g),
\qquad
Rf=\lambda_0f.
\]
Thus $\lambda_0\,1$ is an extension of $R$, and hence, by maximality,
\[
R=\lambda_0\,1.
\]
But this operator is plainly reduced by every closed linear manifold.

Reducibility is therefore a vestige of transformability to diagonal form, and as such is most closely connected with the possibility of an eigenvalue representation. Irreducibility of a matrix already signifies the occurrence of exceptional cases in elementary-divisor theory, something excluded by Hermitian character in finite dimensions and, as we have seen, also in $\mathfrak H$ for hypermaximal operators. From this point of view, the following property of $\overline R$ is particularly noteworthy.

\begin{namedstatement}{Theorem}{39}\label{thm:39}
A maximal H.O. is irreducible if and only if, for a suitable choice of the complete orthonormal system $\psi_0,\psi_1,\psi_2,\ldots$ in \hyperref[thm:37]{Theorem 37}, it coincides with $\overline R$ or with $-\overline R$. Both alternatives can never occur for the same H.O.
\end{namedstatement}

\begin{proof}
Let $R$ be maximal and irreducible. By the preceding discussion it is certainly not hypermaximal, so \hyperref[thm:38]{Theorem 38} applies. Since the $\mathfrak M_1$ occurring there reduces $R$ and is nonzero, it must equal $\mathfrak H$. Thus $p=1$ and $\mathfrak N$ is absent. Hence $R=R_1$, and by $\beta)$ of that theorem it is $\overline R$ or $-\overline R$, with a suitable choice of $\psi_0,\psi_1,\psi_2,\ldots$. In the former case the deficiency indices of $R$ are $0,1$, while in the latter they are $1,0$; the two alternatives are therefore incompatible.

It remains to show that $\overline R$, and hence also $-\overline R$, is irreducible; equivalently, by \hyperref[thm:25]{Theorem 25}, that $\overline U$ is irreducible. Suppose that $\overline U$ is reduced by $\mathfrak M$, and hence also by
\[
\mathfrak N=\mathfrak H-\mathfrak M.
\]
The everywhere-defined $\overline U$ maps each of these manifolds into a subset of itself. Put
\[
\mathfrak M'=\overline U(\mathfrak M),
\qquad
\mathfrak N'=\overline U(\mathfrak N).
\]
If $\mathfrak M'=\mathfrak M$ and $\mathfrak N'=\mathfrak N$, then the range of $\overline U$ would contain both $\mathfrak M$ and $\mathfrak N$, and hence also
\[
\mathfrak M+\mathfrak N=\mathfrak H,
\]
which is not the case. Thus $\mathfrak M'$ or $\mathfrak N'$ is a proper subset of $\mathfrak M$ or $\mathfrak N$, respectively, and
\[
\mathfrak M-\mathfrak M'\ne(0)
\qquad\text{or}\qquad
\mathfrak N-\mathfrak N'\ne(0).
\]
Let $f\ne0$ belong to $\mathfrak M-\mathfrak M'$, respectively to $\mathfrak N-\mathfrak N'$. It is then orthogonal to $\mathfrak M'$, respectively $\mathfrak N'$. Being an element of $\mathfrak M$, respectively $\mathfrak N$, it is also orthogonal to $\mathfrak N$ and $\mathfrak N'$, respectively to $\mathfrak M$ and $\mathfrak M'$. It is therefore orthogonal to the entire range of $\overline U$,
\[
\mathfrak M'+\mathfrak N'.
\]
That is, it is orthogonal to every $\psi_1,\psi_2,\ldots$, and hence
\[
f=a\psi_0
\qquad(a\ne0).
\]
Consequently $\psi_0$ also belongs to $\mathfrak M-\mathfrak M'$, respectively to $\mathfrak N-\mathfrak N'$, and thus to $\mathfrak M$, respectively $\mathfrak N$. But $\mathfrak M$ and $\mathfrak N$ contain their $\overline U$-images; whichever of them contains $\psi_0$ therefore also contains $\psi_1,\psi_2,\ldots$, and hence all of $\mathfrak H$. Thus either
\[
\mathfrak M=\mathfrak H,
\]
or
\[
\mathfrak H-\mathfrak M=\mathfrak N=\mathfrak H,
\qquad
\mathfrak M=(0).
\]
This proves everything.
\end{proof}

Together with $\overline R$, the operator
\[
a\overline R+b1
\qquad(a,b\ \text{real},\ a\ne0)
\]
is also maximal and irreducible, and hence, up to the choice of coordinate system, equals either $\overline R$ or $-\overline R$. For $a>0$ its deficiency indices are $0,1$, so in that case $\overline R$ is the possibility; for $a<0$ they are $1,0$, so in that case $-\overline R$ is the possibility.

Finally, we mention that the decomposition of the H.O. $R$ in \hyperref[thm:38]{Theorem 38} admits a kind of converse, which makes it possible to construct closed linear Hermitian operators $R$ having any prescribed deficiency indices $m,n$. Since $m=n=0$ is realized by every hypermaximal H.O. (for example, by $1$), we may assume $m+n>0$.

There then exist $m+n$ infinite-dimensional closed linear manifolds
\[
\mathfrak M_1,\ldots,\mathfrak M_m,
\mathfrak N_1,\ldots,\mathfrak N_n
\]
which are pairwise orthogonal and together span $\mathfrak H$ as a closed linear manifold.\footnote{\label{fn:57}Let $\varphi_1,\varphi_2,\ldots$ be a complete normalized orthogonal system. Divide it into $m+n$ subsequences
\[
\psi_1^1,\psi_2^1,\ldots;\ \ldots;\
\psi_1^m,\psi_2^m,\ldots;\
\chi_1^1,\chi_2^1,\ldots;\ \ldots;\
\chi_1^n,\chi_2^n,\ldots
\]
(either $m$ or $n$ may also equal $\infty$), and let
$\mathfrak M_1,\ldots,\mathfrak M_m,\mathfrak N_1,\ldots,\mathfrak N_n$
be the closed linear manifolds spanned by these respective subsequences.}
Each is isomorphic to $\mathfrak H$, so we can realize $-\overline R$ in the first $m$ of them and $\overline R$ in the remaining $n$; denote the resulting operators by
\[
R_1,\ldots,R_m,S_1,\ldots,S_n.
\]
It is easily verified that, when combined according to the prescription of \hyperref[thm:21]{Theorem 21}, these yield a closed linear H.O. $R$, and that $R$ has deficiency indices $m,n$.

\section*{XI. Semiboundedness}\label{chap:XI}
\addcontentsline{toc}{section}{XI. Semiboundedness}

We shall give several properties that every nonhypermaximal (but maximal) H.O. must lack. (Most of the results of the following theorem will soon also be obtained by another route.)

\begin{namedstatement}{Theorem}{40}\label{thm:40}
A maximal but nonhypermaximal H.O. can be neither semibounded above nor semibounded below (and hence, a fortiori, cannot be bounded); it can neither be defined everywhere nor take every value; and it cannot be real in any complete normalized orthogonal system (cf. \hyperref[fn:source-20]{note 20}).
\end{namedstatement}

\begin{proof}
For such an operator there is a closed linear manifold $\mathfrak M$ reducing it such that the part lying in $\mathfrak M$ is $\overline R$ or $-\overline R$, by \hyperref[thm:38]{Theorem 38}. If it possessed any of the first four properties just named, then $\overline R$ or $-\overline R$ would have to possess the same property. But $\overline R$ (and hence also $-\overline R$) has none of them, as will be shown in its detailed discussion in \hyperref[app:IV]{Appendix IV}, \S~4.

As for the last property: if an operator is real, in the sense of \hyperref[fn:source-20]{note 20}, in a complete normalized orthogonal system, then for it there is no distinction between $i$ and $-i$, and hence its deficiency indices are equal. For a maximal but nonhypermaximal H.O., however, they are unequal; cf. \hyperref[thm:33]{Theorem 33}.
\end{proof}

We now turn to the systematic investigation of semibounded operators.

\begin{namedstatement}{Theorem}{41}\label{thm:41}
A H.O. $R$ that takes every value is hypermaximal. It takes each value exactly once.
\end{namedstatement}

\begin{proof}
Let $f$ be an extension element of $R$, and let $f^*$ be assigned to it. By assumption there is a $g$ for which $Rg$ is defined and $Rg=f^*$. For every $h$ for which $Rh$ is defined,
\[
(f,Rh)=(f^*,h)=(Rg,h)=(g,Rh).
\]
Since $Rh$ takes every value, it follows that $f=g$; hence $Rf$ is defined. Thus $R$ is hypermaximal. At the same time we have proved that $Rf=Rg$ implies $f=g$, by putting $f^*=Rf$.
\end{proof}

\begin{namedstatement}{Theorem}{42}\label{thm:42}
A linear H.O. $R$ satisfying
\[
(f,Rf)\ge |f|^2
\]
for every $f$ in its domain can be extended to a definite hypermaximal H.O.
\end{namedstatement}

\begin{proof}
First form the closed linear H.O. $\widetilde R$. Plainly it too satisfies
\[
(f,\widetilde Rf)\ge|f|^2.
\]
Let $\mathfrak M$ be the range of $\widetilde R$; it is plainly a linear manifold. We have
\[
|f|^2\le(f,\widetilde Rf)\le|f|\,|\widetilde Rf|,
\qquad
|\widetilde Rf|\ge|f|,
\]
and
\[
|\widetilde Rf-\widetilde Rg|
=|\widetilde R(f-g)|\ge|f-g|.
\]
Thus $\widetilde Rf=\widetilde Rg$ implies $f=g$, so $\widetilde R$ has an inverse $\widetilde R^{-1}$. This inverse is plainly also closed and linear, is defined on $\mathfrak M$, and is continuous because
\[
|f-g|\ge
|\widetilde R^{-1}f-\widetilde R^{-1}g|.
\]
Consequently $\mathfrak M$ is closed, and we may form the closed linear manifold
\[
\mathfrak N=\mathfrak H-\mathfrak M.
\]

Let $f\in\mathfrak N$. Whenever $\widetilde Rg$ is defined, it belongs to $\mathfrak M$, and hence
\[
(f,\widetilde Rg)=0.
\]
Thus $f$ is an extension element to which $0$ is assigned. If $\widetilde Rf$ is defined, therefore, $\widetilde Rf=0$, and then $f=0$. Hence the domain $\mathfrak A$ of $\widetilde R$, which is a linear manifold, and $\mathfrak N$ are linearly independent.

We may therefore define a new operator $S$ as follows: $Sf$ is defined when
\[
f=g+h,\qquad g\in\mathfrak A,\quad h\in\mathfrak N,
\]
the decomposition being unique, and its value is
\[
Sf=\widetilde Rg.
\]
The operator $S$ is an extension of $\widetilde R$ and is a H.O. Indeed, if
\[
f=g+h,\qquad f'=g'+h',
\]
with $g,g'\in\mathfrak A$ and $h,h'\in\mathfrak N$, then
\begin{align*}
(f,Sf')
 &=(g+h,\widetilde Rg')
  =(g,\widetilde Rg')\\
 &=(\widetilde Rg,g')
  =(\widetilde Rg,g'+h')
  =(Sf,f').
\end{align*}

The manifold $\mathfrak N$ reduces $S$. By definition, $Sf$ is always defined on $\mathfrak N$ and equals $0$ there; and since every $Sf$ equals some $\widetilde Rg$ and therefore lies in $\mathfrak M$, we also have
\[
P_{\mathfrak N}Sf=0=SP_{\mathfrak N}f.
\]
Hence $\mathfrak M=\mathfrak H-\mathfrak N$ also reduces $S$. Denote the parts of $S$ lying in $\mathfrak M,\mathfrak N$ by $S'$ and $S''$, respectively.

The range of $S$ is $\mathfrak M$. Since
\[
Sf=P_{\mathfrak M}Sf=SP_{\mathfrak M}f,
\]
$S$ already takes each of its values in $\mathfrak M$, so $S'$ also has range $\mathfrak M$. As the one-to-one linear image of the infinite-dimensional domain $\mathfrak A$ (which is infinite-dimensional because it is everywhere dense), $\mathfrak M$ is likewise infinite-dimensional and hence is a Hilbert space. We may therefore apply \hyperref[thm:41]{Theorem 41}: $S'$ is hypermaximal, while $S''$ is in fact $0$.

It follows immediately that $S$ itself must be hypermaximal. Moreover, it is definite. Indeed, if
\[
f=g+h,\qquad g\in\mathfrak A,\quad h\in\mathfrak N,
\]
then
\[
(f,Sf)=(g+h,\widetilde Rg)
=(g,\widetilde Rg)\ge|g|^2\ge0.
\]
This proves everything.
\end{proof}

\begin{namedstatement}{Theorem}{43}\label{thm:43}
Suppose a linear H.O. $R$ satisfies either
\[
(f,Rf)\ge C|f|^2
\]
for every $f$ in its domain, or
\[
(f,Rf)\le C|f|^2
\]
for every such $f$; these are precisely the conditions of semiboundedness. Then $R$ can be extended to a hypermaximal H.O. $S$ satisfying, respectively,
\[
(f,Sf)\ge C'|f|^2
\qquad\text{or}\qquad
(f,Sf)\le C'|f|^2.
\]
Here $C'$ may be chosen to be any number less than $C$, respectively greater than $C$.\footnote{\label{fn:58}The theorem probably remains true for $C'=C$, but no proof of this is available.}
\end{namedstatement}

\begin{proof}
For $C'=0$, $C=1$, and the $\ge$ sign, this is \hyperref[thm:42]{Theorem 42}. All other cases reduce to this one by considering, in place of $R,S$,
\[
\frac{1}{C-C'}R-\frac{C'}{C-C'}1,
\qquad
\frac{1}{C-C'}S-\frac{C'}{C-C'}1.
\]
\end{proof}

\section*{XII. Pathology of Unbounded Operators}\label{chap:XII}
\addcontentsline{toc}{section}{XII. Pathology of Unbounded Operators}

We begin with several auxiliary theorems.

\begin{namedstatement}{Theorem}{44}\label{thm:44}
If $R$ is a linear H.O. that is not semibounded above, then, for every $n=1,2,\ldots$ and every $C$, there exists an $n$-dimensional linear manifold on which $Rf$ is always defined and throughout which
\[
(f,Rf)\ge C|f|^2.
\]
\end{namedstatement}

\begin{proof}
Let $f_1,\ldots,f_k$ be arbitrary elements for which $Rf_1,\ldots,Rf_k$ are defined. We first show that there must be a $\varphi$ such that
\[
\varphi\perp f_1,\ldots,f_k,\qquad
|\varphi|=1,\qquad
R\varphi\ \text{is defined},\qquad
(\varphi,R\varphi)\ge D.
\]
Suppose the contrary. Since only the linear manifold spanned by $f_1,\ldots,f_k$ matters, we may ``orthogonalize'' them (cf. \hyperref[thm:8]{Theorem 8}), replacing them by a normalized orthogonal system
\[
\varphi_1,\ldots,\varphi_l\qquad(l\le k).
\]
Let $f$ be arbitrary, subject only to $Rf$ being defined. By adjoining $f$ and orthogonalizing, we can write
\[
f=\sum_{\mu=1}^{l}a_\mu\varphi_\mu+b\varphi,
\]
where
\[
|\varphi|=1,\qquad
\varphi\perp\varphi_1,\ldots,\varphi_l.
\]
Since $Rf,R\varphi_1,\ldots,R\varphi_l$ are defined, $R\varphi$ is also defined; by our contrary assumption,
\[
(\varphi,R\varphi)\le D.
\]
Now
\[
|f|^2=\sum_{\mu=1}^{l}|a_\mu|^2+|b|^2,
\]
and
\begin{align*}
(f,Rf)
={}&\sum_{\mu,\nu=1}^{l}
a_\mu\overline{a_\nu}(\varphi_\mu,R\varphi_\nu)\\
&+2\operatorname{Re}\sum_{\mu=1}^{l}
b\overline{a_\mu}(\varphi,R\varphi_\mu)
 +|b|^2(\varphi,R\varphi)\\
\le{}&
\sum_{\mu,\nu=1}^{l}
|a_\mu|\,|a_\nu|\,|\varphi_\mu|\,|R\varphi_\nu|\\
&+2\sum_{\mu=1}^{l}
|b|\,|a_\mu|\,|\varphi|\,|R\varphi_\mu|
+|b|^2D.
\end{align*}
If
\[
D'=\max_{\mu=1,\ldots,l}\bigl(|R\varphi_\mu|,D\bigr),
\]
then
\begin{align*}
(f,Rf)
&\le D'\left(\sum_{\mu=1}^{l}|a_\mu|+|b|\right)^2\\
&\le(l+1)D'
\left(\sum_{\mu=1}^{l}|a_\mu|^2+|b|^2\right)
=(l+1)D'|f|^2.
\end{align*}
Thus $R$ would be semibounded above, contrary to the hypothesis.

We go one step further. Let $g_1,\ldots,g_k$ now be entirely arbitrary. Then there exists a $\varphi$ for which $R\varphi$ is defined and
\[
|\varphi|=1,\qquad
(\varphi,R\varphi)\ge D,\qquad
|(\varphi,g_\mu)|\le\varepsilon
\quad(\mu=1,\ldots,k;\ \varepsilon>0).
\]
(Thus it is not quite exactly orthogonal.) Indeed, the vectors $f$ for which $Rf$ is defined are everywhere dense. Choose $f_1,\ldots,f_k$ with
\[
|g_\mu-f_\mu|\le\varepsilon
\qquad(\mu=1,\ldots,k),
\]
and choose $\varphi$ as in the result just proved. Then
\[
|\varphi|=1,\qquad
(\varphi,R\varphi)\ge D,
\]
and
\[
|(\varphi,g_\mu)|
=|(\varphi,g_\mu-f_\mu)|
\le|\varphi|\,|g_\mu-f_\mu|
\le\varepsilon.
\]

We can now prove the assertion itself. Choose successively:
\[
\begin{array}{l}
\varphi_1,\quad R\varphi_1\text{ defined},\quad
|\varphi_1|=1,\quad(\varphi_1,R\varphi_1)\ge D;\\
\varphi_2,\quad R\varphi_2\text{ defined},\quad
|\varphi_2|=1,\quad(\varphi_2,R\varphi_2)\ge D,\quad
|(R\varphi_1,\varphi_2)|\le1;\\
\hfill\vdots\\
\varphi_n,\quad R\varphi_n\text{ defined},\quad
|\varphi_n|=1,\quad(\varphi_n,R\varphi_n)\ge D,\\
\hfill
|(R\varphi_1,\varphi_n)|\le1,\ldots,
|(R\varphi_{n-1},\varphi_n)|\le1.
\end{array}
\]
Thus all $R\varphi_\mu$ are defined,
\[
|\varphi_\mu|=1,\qquad
(\varphi_\mu,R\varphi_\mu)\ge D,
\]
and
\[
|(R\varphi_\mu,\varphi_\nu)|\le1
\qquad(\mu<\nu),
\]
hence, by the Hermitian property, for every $\mu\ne\nu$. It follows that
\begin{align*}
\left|\sum_{\mu=1}^{n}a_\mu\varphi_\mu\right|^2
&=\sum_{\mu,\nu=1}^{n}
a_\mu\overline{a_\nu}(\varphi_\mu,\varphi_\nu)\\
&\le\sum_{\mu,\nu=1}^{n}
|a_\mu|\,|a_\nu|\,|\varphi_\mu|\,|\varphi_\nu|\\
&\le\left(\sum_{\mu=1}^{n}|a_\mu|\right)^2
\le n\sum_{\mu=1}^{n}|a_\mu|^2,
\end{align*}
while
\begin{align*}
\left(\sum_{\mu=1}^{n}a_\mu\varphi_\mu,\,
R\left\{\sum_{\mu=1}^{n}a_\mu\varphi_\mu\right\}\right)
&=\sum_{\mu,\nu=1}^{n}
a_\mu\overline{a_\nu}(\varphi_\mu,R\varphi_\nu)\\
&=\sum_{\mu=1}^{n}|a_\mu|^2(\varphi_\mu,R\varphi_\mu)
 +\sum_{\substack{\mu,\nu=1\\\mu\ne\nu}}^{n}
a_\mu\overline{a_\nu}(\varphi_\mu,R\varphi_\nu)\\
&\ge D\sum_{\mu=1}^{n}|a_\mu|^2
-\sum_{\substack{\mu,\nu=1\\\mu\ne\nu}}^{n}
|a_\mu|\,|a_\nu|\\
&=(D+1)\sum_{\mu=1}^{n}|a_\mu|^2
-\left(\sum_{\mu=1}^{n}|a_\mu|\right)^2\\
&\ge(D+1-n)\sum_{\mu=1}^{n}|a_\mu|^2.
\end{align*}
If $D>n-1$, the second inequality shows that
\[
\sum_{\mu=1}^{n}a_\mu\varphi_\mu=0
\]
implies
\[
\sum_{\mu=1}^{n}|a_\mu|^2=0,
\qquad
a_1=\cdots=a_n=0.
\]
Thus $\varphi_1,\ldots,\varphi_n$ are linearly independent, and the linear manifold of all $\sum_{\mu=1}^{n}a_\mu\varphi_\mu$ is $n$-dimensional. Moreover, the two inequalities give
\[
\left(\sum_{\mu=1}^{n}a_\mu\varphi_\mu,\,
R\left\{\sum_{\mu=1}^{n}a_\mu\varphi_\mu\right\}\right)
\ge
\frac{D+1-n}{n}
\left|\sum_{\mu=1}^{n}a_\mu\varphi_\mu\right|^2.
\]
Taking also
\[
D\ge nC+n-1,
\]
we therefore have throughout this linear manifold
\[
(f,Rf)\ge C|f|^2.
\]
This proves the theorem.
\end{proof}

\begin{namedstatement}{Theorem}{45}\label{thm:45}
Let $R$ be a linear H.O. that is not semibounded above, let $f_1,\ldots,f_k$ be arbitrary elements, and let $C$ be any number. Then there exists a $\varphi$ for which $R\varphi$ is defined such that
\[
\varphi\perp f_1,\ldots,f_k,\qquad
|\varphi|=1,\qquad
(\varphi,R\varphi)\ge C.
\]
\end{namedstatement}

\begin{proof}
\footnote{\label{fn:59}When $Rf_1,\ldots,Rf_k$ were defined, we already obtained this at the beginning of the proof of \hyperref[thm:44]{Theorem 44}; but we must free ourselves from that assumption.}
Choose the linear manifold in \hyperref[thm:44]{Theorem 44} to be $(k+1)$-dimensional. In it the equations
\[
(f_1,\varphi)=0,\ldots,(f_k,\varphi)=0,
\]
which represent at most $k$ independent linear conditions, have a solution $f\ne0$. Multiplication by a scalar allows us to arrange $|f|=1$. This $f$ has all the required properties.
\end{proof}

\begin{namedstatement}{Theorem}{46}\label{thm:46}
Every linear H.O. $R$ that is not semibounded above is an extension of a definite linear H.O. $S$.
\end{namedstatement}

\begin{proof}
It suffices to exhibit an everywhere-dense sequence $f_1,f_2,\ldots$ such that all $Rf_\mu$ are defined and
\[
\left(\sum_{\mu=1}^{n}x_\mu f_\mu,\,
R\left\{\sum_{\mu=1}^{n}x_\mu f_\mu\right\}\right)\ge0
\qquad(n=1,2,\ldots).
\]
Indeed, the operator $S$ defined on the linear manifold spanned by $f_1,f_2,\ldots$---an everywhere-dense, and hence nonclosed, manifold---and agreeing there with $R$ is plainly a linear definite H.O., and $R$ is its extension.

Let $g_1,g_2,\ldots$ be an everywhere-dense sequence such that all $Rg_\mu$ are defined. Thus $g_1,g_2,\ldots$ may be taken to be a subsequence dense in the domain $\mathfrak A$ of $R$; the domain itself is everywhere dense, and such a sequence exists because $\mathfrak A$, as a subset of $\mathfrak H$, is separable; cf. \hyperref[fn:33]{note 33}. Let $C_1,C_2,\ldots$ be a sequence of numbers to be specified later.

Using \hyperref[thm:45]{Theorem 45}, choose a sequence $\varphi_1,\varphi_2,\ldots$ with the following properties:
\[
\begin{array}{lll}
R\varphi_1\text{ defined},&
|\varphi_1|=1,&
(\varphi_1,R\varphi_1)\ge C_1;\\
R\varphi_2\text{ defined},&
|\varphi_2|=1,&
(\varphi_2,R\varphi_2)\ge C_2,\quad
(R\varphi_1,\varphi_2)=0;\\
R\varphi_3\text{ defined},&
|\varphi_3|=1,&
(\varphi_3,R\varphi_3)\ge C_3,\quad
(R\varphi_1,\varphi_3)=(R\varphi_2,\varphi_3)=0;\\
&\vdots&
\end{array}
\]
Thus all $R\varphi_\mu$ are defined,
\[
|\varphi_\mu|=1,\qquad
(\varphi_\mu,R\varphi_\mu)\ge C_\mu,
\qquad
(R\varphi_\mu,\varphi_\nu)=0
\quad(\mu<\nu),
\]
and therefore the last equality holds whenever $\mu\ne\nu$. Put
\[
f_1=g_1+\varphi_1,\qquad
f_2=g_2+\frac12\varphi_2,\qquad
f_3=g_3+\frac13\varphi_3,\qquad\ldots,
\]
and claim that, with a suitable choice of $C_1,C_2,\ldots$, the sequence $f_1,f_2,\ldots$ has all the desired properties.

First, it is everywhere dense:
\[
|f_n-g_n|=\left|\frac1n\varphi_n\right|=\frac1n\longrightarrow0.
\]
For every $f$ there is a subsequence $g_{N_1},g_{N_2},\ldots$ of $g_1,g_2,\ldots$ such that $g_{N_j}\to f$, and then also $f_{N_j}\to f$.

We now compute $(f,Rf)$ for the linear combinations of $f_1,f_2,\ldots$. For every $n$,
\begin{align*}
&\left(\sum_{\mu=1}^{n}a_\mu f_\mu,\,
R\left\{\sum_{\mu=1}^{n}a_\mu f_\mu\right\}\right)\\
&\quad=
\left(
\sum_{\mu=1}^{n}a_\mu g_\mu
+\sum_{\mu=1}^{n}\frac{a_\mu}{\mu}\varphi_\mu,\,
\sum_{\mu=1}^{n}a_\mu Rg_\mu
+\sum_{\mu=1}^{n}\frac{a_\mu}{\mu}R\varphi_\mu
\right)\\
&\quad=
\sum_{\mu,\nu=1}^{n}
a_\mu\overline{a_\nu}(g_\mu,Rg_\nu)
+2\operatorname{Re}\sum_{\mu,\nu=1}^{n}
\frac{a_\mu\overline{a_\nu}}{\mu}
(\varphi_\mu,Rg_\nu)\\
&\qquad
+\sum_{\mu=1}^{n}
\frac{|a_\mu|^2}{\mu^2}
(\varphi_\mu,R\varphi_\mu)
+\sum_{\substack{\mu,\nu=1\\\mu\ne\nu}}^{n}
\frac{a_\mu\overline{a_\nu}}{\mu\nu}
(\varphi_\mu,R\varphi_\nu).
\end{align*}
Consequently,
\begin{align*}
&\left(\sum_{\mu=1}^{n}a_\mu f_\mu,\,
R\left\{\sum_{\mu=1}^{n}a_\mu f_\mu\right\}\right)\\
&\quad\ge
-\sum_{\mu,\nu=1}^{n}
|a_\mu|\,|a_\nu|\,|g_\mu|\,|Rg_\nu|
-2\sum_{\mu,\nu=1}^{n}
|a_\mu|\,|a_\nu|\,
\frac{|\varphi_\mu|\,|Rg_\nu|}{\mu}
+\sum_{\mu=1}^{n}|a_\mu|^2\frac{C_\mu}{\mu^2}\\
&\quad=
\sum_{\mu=1}^{n}\frac{C_\mu}{\mu^2}|a_\mu|^2
-\sum_{\mu,\nu=1}^{n}
\left(|g_\mu|\,|Rg_\nu|+\frac{2|Rg_\nu|}{\mu}\right)
|a_\mu|\,|a_\nu|.
\end{align*}
Put
\[
\overline A_{\mu\nu}
=|g_\mu|\,|Rg_\nu|+\frac{2|Rg_\nu|}{\mu}.
\]
This depends on $g_1,g_2,\ldots$, but not on $f_1,f_2,\ldots$ or on $\varphi_1,\varphi_2,\ldots$, that is, not on $C_1,C_2,\ldots$.\translatornote{The source changes the factor $2/\mu$ in the expansion to $2/\nu$ in the following bound and then drops the factor $2$ in the equality and in the definition of $\overline A_{\mu\nu}$.  The three occurrences have been made consistent with the preceding expansion.} We estimate the subtracted term:
\begin{align*}
\sum_{\mu,\nu=1}^{n}
\overline A_{\mu\nu}|a_\mu|\,|a_\nu|
={}&
\sum_{\mu=1}^{n}\overline A_{\mu\mu}|a_\mu|^2\\
&+\sum_{\substack{\mu,\nu=1\\\mu<\nu}}^{n}
(\overline A_{\mu\nu}+\overline A_{\nu\mu})
|a_\mu|\,|a_\nu|\\
\le{}&
\sum_{\mu=1}^{n}\overline A_{\mu\mu}|a_\mu|^2\\
&+\sum_{\substack{\mu,\nu=1\\\mu<\nu}}^{n}
\left[
2^{\nu-\mu-1}
(\overline A_{\mu\nu}+\overline A_{\nu\mu})^2|a_\mu|^2
+\frac{1}{2^{\nu-\mu}}|a_\nu|^2
\right]\\
={}&
\sum_{\mu=1}^{n}
\left\{
\overline A_{\mu\mu}
+\sum_{\rho=1}^{\mu-1}
2^{\mu-\rho-1}
(\overline A_{\mu\rho}+\overline A_{\rho\mu})^2
+\sum_{\rho=\mu+1}^{n}\frac{1}{2^{\rho-\mu}}
\right\}|a_\mu|^2\\
\le{}&
\sum_{\mu=1}^{n}
\left\{
\overline A_{\mu\mu}+1
+\sum_{\sigma=0}^{\mu-2}
2^\sigma
(\overline A_{\mu,\mu-\sigma-1}
 +\overline A_{\mu-\sigma-1,\mu})^2
\right\}|a_\mu|^2.
\end{align*}
Denote the expression in braces by $\overline B_\mu$. Then
\[
\left(\sum_{\mu=1}^{n}a_\mu f_\mu,\,
R\left\{\sum_{\mu=1}^{n}a_\mu f_\mu\right\}\right)
\ge
\sum_{\mu=1}^{n}
\left(\frac{C_\mu}{\mu^2}-\overline B_\mu\right)|a_\mu|^2.
\]
It therefore suffices to set
\[
C_\mu=\mu^2\overline B_\mu,
\]
and the proof is complete.\footnote{\label{fn:60}For the purposes of the work announced in \hyperref[fn:source-22]{note 22}, we record here a certain sharpening of this estimate, although we do not need it here. We have
\begin{align*}
\left|\sum_{\mu=1}^{n}a_\mu f_\mu\right|^2
&=
\left|
\sum_{\mu=1}^{n}a_\mu g_\mu
+\sum_{\mu=1}^{n}\frac{a_\mu}{\mu}\varphi_\mu
\right|^2\\
&=
\sum_{\mu,\nu=1}^{n}
a_\mu\overline{a_\nu}(g_\mu,g_\nu)
+2\operatorname{Re}\sum_{\mu,\nu=1}^{n}
\frac{a_\mu\overline{a_\nu}}{\mu}(\varphi_\mu,g_\nu)\\
&\qquad
+\sum_{\mu,\nu=1}^{n}
\frac{a_\mu\overline{a_\nu}}{\mu\nu}
(\varphi_\mu,\varphi_\nu)\\
&\le
\sum_{\mu,\nu=1}^{n}
\left\{
|g_\mu|\,|g_\nu|
+\frac{2|\varphi_\mu|\,|g_\nu|}{\mu}
+\frac{|\varphi_\mu|\,|\varphi_\nu|}{\mu\nu}
\right\}
|a_\mu|\,|a_\nu|\\
&=
\sum_{\mu,\nu=1}^{n}
\left(
|g_\mu|\,|g_\nu|
+\frac{2|g_\nu|}{\mu}
+\frac{1}{\mu\nu}
\right)
|a_\mu|\,|a_\nu|.
\end{align*}
Denote the expression in parentheses by $\overline C_{\mu\nu}$. As in the text,
\[
\sum_{\mu,\nu=1}^{n}
\overline C_{\mu\nu}|a_\mu|\,|a_\nu|
\le
\sum_{\mu=1}^{n}
\left\{
\overline C_{\mu\mu}+1
+\sum_{\sigma=0}^{\mu-2}
2^\sigma
(\overline C_{\mu,\mu-\sigma-1}
 +\overline C_{\mu-\sigma-1,\mu})^2
\right\}|a_\mu|^2.
\]
Call this brace-enclosed expression $\overline D_\mu$. Then
\[
\left|\sum_{\mu=1}^{n}a_\mu f_\mu\right|^2
\le\sum_{\mu=1}^{n}\overline D_\mu|a_\mu|^2.
\]
If we now put
\[
C_\mu=\mu^2(\overline B_\mu+\mu\overline D_\mu),
\]
then, whenever $a_1=\cdots=a_{p-1}=0$ and $n\ge p$,
\begin{align*}
&\left(\sum_{\mu=p}^{n}a_\mu f_\mu,\,
R\left\{\sum_{\mu=p}^{n}a_\mu f_\mu\right\}\right)
-p\left|\sum_{\mu=p}^{n}a_\mu f_\mu\right|^2\\
&\qquad\ge
\sum_{\mu=p}^{n}
\left(
\frac{C_\mu}{\mu^2}
-\overline B_\mu-p\overline D_\mu
\right)|a_\mu|^2
\ge0.
\end{align*}
Thus, on the linear manifold spanned by
\[
f_p,f_{p+1},f_{p+2},\ldots,
\]
which is again everywhere dense and hence nonclosed, one even has
\[
(f,Rf)\ge p|f|^2
\]
for every $p=1,2,\ldots$.}%
\textsuperscript{,}\translatornote{In note 60 the source drops the factor
\(2\) from the bound on the mixed term, although it is present in the
preceding exact expansion. The two occurrences have been restored here;
consequently the definitions of \(\overline C_{\mu\nu}\) and
\(\overline D_\mu\) use the corrected bound.}
\end{proof}

We have already obtained a rather ``pathological'' result, which is easily sharpened further.

\begin{namedstatement}{Theorem}{47}\label{thm:47}
Let $R$ be a linear H.O. that is not semibounded above or, respectively, not semibounded below. Then there is a hypermaximal H.O. $S$ satisfying throughout
\[
(f,Sf)\ge C|f|^2,
\qquad\text{respectively}\qquad
(f,Sf)\le C|f|^2,
\]
such that $R$ and $S$ are both extensions of the same linear H.O. Here $C$ may be chosen arbitrarily.
\end{namedstatement}

\begin{proof}
For $C=-1$ and failure of semiboundedness above, this follows from \hyperref[thm:46]{Theorem 46} and \hyperref[thm:43]{Theorem 43}, in which one takes $C=0$ and $C'=-1$. Every other value of $C$ follows from this case by considering, in place of $R,S$,
\[
R-(C+1)1,\qquad S-(C+1)1,
\]
or, respectively,
\[
-R+(C-1)1,\qquad -S+(C-1)1.
\]
\end{proof}

The strangest consequences for unbounded Hermitian operators follow from this. The relation of a H.O. to its matrix, and of that matrix to its sections, is thereby qualitatively different from the bounded case. In particular, the ``section spectra'' of unbounded matrices have very little to do with the spectrum of the matrix itself and do not, in any sense, ``approximate'' it as they do in Hilbert's theory. We shall say something about the operator--matrix relation in \hyperref[app:III]{Appendix III}; the detailed discussion of the properties of unbounded matrices will be given in the work announced in \hyperref[fn:source-22]{note 22}.

We give several further characteristic properties of unbounded Hermitian operators. In contrast to \hyperref[thm:41]{Theorems 41} through \hyperref[thm:47]{47}, we now restrict ourselves to closed linear operators.

\begin{namedstatement}{Theorem}{48}\label{thm:48}
Let $R$ be a closed linear H.O. Then the following four assertions are equivalent:
\[
\begin{array}{ll}
\alpha)&R\text{ is bounded},\\
\beta)&R\text{ is bounded and hypermaximal},\\
\gamma)&R\text{ is defined everywhere},\\
\delta)&\text{there is no closed linear H.O. of which }R
\text{ is a proper extension}.
\end{array}
\]\footnote{\label{fn:61}The equivalence of $\alpha)$ and $\gamma)$ is the well-known theorem of Toeplitz, \emph{Mathematische Annalen} \textbf{69} (1911), p.~321.}
\end{namedstatement}

\begin{proof}
Suppose first that $R$ is not semibounded above. Choose $C$ so that
\[
(f,Rf)<C|f|^2
\]
occurs, and then, by \hyperref[thm:47]{Theorem 47}, choose a hypermaximal H.O. $S$ such that
\[
(f,Sf)\ge C|f|^2
\]
always holds. Plainly $S$ cannot be an extension of $R$. Now $R,S$ are extensions of a linear H.O. $T$, and, since they are closed and linear, they are also extensions of the closed linear H.O. $\widetilde T$. Since $S$ is not an extension of $R$, we must have $\widetilde T\ne R$; thus $\delta)$ is not satisfied.

If $R$ is not semibounded below, then $\delta)$ fails for $-R$, and hence also for $R$. Thus unboundedness implies the negation of $\delta)$; equivalently, $\delta)$ implies $\alpha)$.

Assertion $\alpha)$ implies $\gamma)$: for $R$ is closed and continuous, so its domain is closed; but the domain is everywhere dense, and hence must equal $\mathfrak H$. Assertion $\gamma)$ implies $\delta)$: if $R$ were a proper extension of a closed linear H.O. $S$, then $S$ would not be maximal and so, by the corollary to \hyperref[thm:35]{Theorem 35}, could be extended to several maximal Hermitian operators, for example to an $R'\ne R$. Whenever $R'f$ and $Sg$ are defined,
\[
(R'f,g)=(f,R'g)=(f,Sg)=(f,Rg)=(Rf,g).
\]
Since the set of such $g$ is everywhere dense, $R'f=Rf$. Thus $R$ would be an extension of $R'$, with $R'\ne R$, even though $R'$ was to be maximal.

We have the logical scheme
\[
\alpha)\ \Longrightarrow\ \gamma)\ \Longrightarrow\ \delta)\
\Longrightarrow\ \alpha),
\]
so $\alpha),\gamma),\delta)$ are logically equivalent. Assertion $\beta)$ implies $\alpha)$, while $\beta)$ follows from $\alpha)$ and $\gamma)$; hence it too is equivalent to them.
\end{proof}

\begin{namedstatement}{Theorem}{49}\label{thm:49}
Let $R$ be an unbounded closed linear H.O. Then it is an extension of a closed linear H.O. that can be extended to continuum many different maximal Hermitian operators.
\end{namedstatement}

\begin{proof}
By \hyperref[thm:48]{Theorem 48}, $\delta)$, $R$ must be a proper extension of a closed linear H.O. $S$. This $S$ is not maximal and therefore, by the corollary to \hyperref[thm:35]{Theorem 35}, has the required property.
\end{proof}

\section*{Appendix I. Function Spaces}\label{app:I}
\addcontentsline{toc}{section}{Appendix I. Function Spaces}

We shall prove here the assertion announced in \hyperref[chap:I]{Chapter I}: all the function spaces mentioned in \hyperref[intro:I]{Introduction I} and \hyperref[intro:II]{II} are abstract Hilbert spaces, that is, they satisfy conditions \hyperref[axiom:A]{A} through \hyperref[axiom:E]{E} of \hyperref[chap:I]{Chapter I}. For the ``ordinary Hilbert space'' of all sequences
\[
(x_1,x_2,\ldots)
\qquad\text{with}\qquad
\sum_{n=1}^{\infty}|x_n|^2<\infty,
\]
this proof is unnecessary as soon as at least one system satisfying A through E is known: every such system is, by the concluding considerations of \hyperref[chap:I]{Chapter I}, isomorphic to the ordinary Hilbert space, which must therefore also possess properties A through E. Hence we need consider only the function spaces of \hyperref[intro:I]{Introduction I}.

Let $\Omega$ be one of the metric spaces named there, or indeed any more general structure of which we require only the following. On $\Omega$ there are to be notions of ``measurability,'' a ``measure,'' and an ``integral,'' denoted
\[
\int_\Omega f(P)\,dv,
\]
in the Lebesgue sense. We shall not spell out precisely what this ``Lebesgue sense'' is to mean; it would be easy to describe, but somewhat lengthy. The considerations below will in any event reveal which properties of these notions matter. We also assume that $\Omega$ is separable.

The function space $\mathfrak H'$ consists of all measurable complex-valued functions $f(P)$ on $\Omega$ for which
\[
\int_\Omega|f(P)|^2\,dv
\]
is finite; two functions differing only on a set of measure $0$ are not regarded as distinct. The meanings of $af$ and $f+g$ are clear, and we put
\[
(f,g)=\int_\Omega f(P)\overline{g(P)}\,dv.
\]
These definitions are permissible. If $f,g\in\mathfrak H'$, so that
\[
\int_\Omega|f(P)|^2\,dv,\qquad
\int_\Omega|g(P)|^2\,dv
\]
are finite, then $af$ and $f+g$ also belong to $\mathfrak H'$. Indeed,
\[
\int_\Omega|af(P)|^2\,dv
\]
is plainly finite, while
\[
\int_\Omega|f(P)+g(P)|^2\,dv
\]
is finite because
\[
|f(P)+g(P)|^2
\le2|f(P)|^2+2|g(P)|^2.
\]
Moreover,
\[
\int_\Omega f(P)\overline{g(P)}\,dv
\]
is meaningful, and indeed absolutely convergent, because
\[
|f(P)\overline{g(P)}|
\le\frac12|f(P)|^2+\frac12|g(P)|^2.
\]

We now verify conditions A through E in order.

\medskip
\noindent\textbf{A, B.}
These are plainly satisfied.

\medskip
\noindent\textbf{C.}
We must find a sequence
\[
f_1(P),f_2(P),\ldots
\]
which is everywhere dense with respect to the distance
\[
|f-g|
=\sqrt{\int_\Omega|f(P)-g(P)|^2\,dv}.
\]
The measure of $\Omega$ may be infinite, but in any case $\Omega$ can be exhausted by an increasing sequence
\[
\Pi_1,\Pi_2,\ldots
\]
of subsets of finite measure. If $f$ is any function in $\mathfrak H'$, define $f_N$, for $N=1,2,\ldots$, by
\[
f_N(P)=
\begin{cases}
f(P),&P\in\Pi_N\ \text{and}\ |f(P)|\le N,\\
0,&\text{otherwise}.
\end{cases}
\]
For every $P$,
\[
f_N(P)\longrightarrow f(P),
\]
and all the integrals
\[
\int_\Omega|f_N(P)|^2\,dv
\]
are bounded above by $\int_\Omega|f(P)|^2\,dv$. Hence
\[
\int_\Omega|f_N(P)-f(P)|^2\,dv\longrightarrow0;
\]
that is, $f_N\to f$ in our distance. Thus the functions which are nonzero only on a set of finite measure and which are also bounded are everywhere dense. It therefore suffices to find a sequence dense among them.

Let $f$ be a function of this class. Let $M$ be the measure of the set on which it is nonzero, and let $A$ be the supremum of its absolute value. In the closed disk
\[
\{z\in\mathbb C:|z|\le A\},
\]
choose finitely many Gaussian rational numbers
\(\rho_1,\ldots,\rho_k\) (numbers in \(\mathbb Q+i\mathbb Q\))
forming an \(\varepsilon\)-net, where \(\varepsilon>0\) is otherwise
arbitrary.\translatornote{The source partitions the real interval
\([-A,A]\) using rational numbers, although \(\mathfrak H'\) consists of
complex-valued functions. We instead use a finite \(\varepsilon\)-net of
Gaussian rationals in the disk \(|z|\le A\); the resulting bound below is
\(M\varepsilon^2\).} Put
\[
f_\varepsilon(P)=
\begin{cases}
\rho_h,
&f(P)\ne0\ \text{and }\rho_h
\text{ is nearest to }f(P)\text{ among }
\rho_1,\ldots,\rho_k,\\
0,&\text{otherwise}.
\end{cases}
\]
It follows immediately that
\[
\int_\Omega|f_\varepsilon(P)-f(P)|^2\,dv
\le M\varepsilon^2.
\]
Thus $f_\varepsilon\to f$ as $\varepsilon\to0$. Consequently, those functions which take only Gaussian-rational values, take only finitely many distinct values, and for which the set corresponding to each nonzero value has finite measure are also everywhere dense. The desired sequence need only be dense among these functions.

If $\Sigma$ is any subset of $\Omega$ of finite measure, let $f_\Sigma$ be the function equal to $1$ on $\Sigma$ and $0$ elsewhere. The preceding class is simply the class of all functions
\[
\rho_1f_{\Sigma_1}+\cdots+\rho_kf_{\Sigma_k},
\]
where $\rho_1,\ldots,\rho_k$ are Gaussian-rational numbers. If we can find a sequence of sets
\[
\Sigma^{(1)},\Sigma^{(2)},\ldots
\]
which is everywhere dense among all sets of finite measure---the distance between two sets $\Sigma',\Sigma''$ being the distance between their functions $f_{\Sigma'},f_{\Sigma''}$, that is, the measure of their ``difference set''
\[
(\Sigma'+\Sigma'')-(\Sigma'\cdot\Sigma'')\footnote{\label{fn:62}Here $+$, $-$, and $\cdot$ denote union, set difference, and intersection, respectively.}
\]
---then the functions
\[
\rho_1f_{\Sigma^{(n_1)}}+\cdots+
\rho_kf_{\Sigma^{(n_k)}}
\]
with
\[
\begin{gathered}
k=1,2,\ldots,\qquad
\rho_1,\ldots,\rho_k\in\mathbb Q+i\mathbb Q,\\
n_1,\ldots,n_k=1,2,\ldots,
\end{gathered}
\]
would be everywhere dense in $\mathfrak H'$, and we would have reached our goal. It remains to find such a sequence of sets.

Let $\Sigma$ be a set of finite measure $M$. As is well known, $M$ is the infimum of the measures of all open point sets containing $\Sigma$. Let $\Sigma'$ be such a set with measure at most $M+\varepsilon$. Then the ``difference set'' $\Sigma'-\Sigma$ has measure at most $\varepsilon$. Thus the open sets are everywhere dense.

Since $\Omega$ is separable, there is in $\Omega$ a sequence of open sets
\[
\Lambda_1,\Lambda_2,\ldots
\]
forming an ``equivalent neighborhood system,'' in the sense that every open set $\Sigma'$ is the union of all the $\Lambda_n$ contained in it---say, of
\[
\Lambda_{n_1},\Lambda_{n_2},\ldots.
\]
Let the measure of $\Sigma'$ be $M'$. Then the measures of
\[
\Lambda_{n_1},\quad
\Lambda_{n_1}+\Lambda_{n_2},\quad
\Lambda_{n_1}+\Lambda_{n_2}+\Lambda_{n_3},\quad\ldots
\]
converge to $M'$. If, for example, the measure of

% Source PDF page 63 (printed page 65): continuation of Appendix I
\(\Lambda_{n_1}+\cdots+\Lambda_{n_k}\) is already at least
\(M'-\varepsilon\), then the ``difference set''
\[
  \Sigma'-(\Lambda_{n_1}+\cdots+\Lambda_{n_k})
\]
has measure at most \(\varepsilon\).  Thus the sets
\[
  \Lambda_{n_1}+\cdots+\Lambda_{n_k}
  \qquad
  (k=1,2,\ldots;\ n_1,\ldots,n_k=1,2,\ldots)
\]
are likewise everywhere dense,\footnote{\label{fn:63}Naturally, we take only those
of finite measure.} and there are only countably many of them.  This
brings us to our goal.

\medskip
\noindent
\textbf{D.} Let \(k=1,2,\ldots\).  Choose \(k\) pairwise disjoint sets
\(K_1,\ldots,K_k\) in \(\Omega\), each of finite, nonzero measure.  Let
\(f_h(P)=1\) on \(K_h\) and \(f_h(P)=0\) elsewhere
\((h=1,\ldots,k)\).  Then \(f_1,\ldots,f_k\) are plainly linearly
independent (and belong to \(\mathfrak H'\)).

\medskip
\noindent
\textbf{E.} Let \(f_1,f_2,\ldots\) be a sequence of functions satisfying
Cauchy's convergence criterion: for every \(\varepsilon>0\) there is an
\(N=N(\varepsilon)\) such that, for \(m,n\geq N\),
\[
  \int_\Omega |f_m(P)-f_n(P)|^2\,dv\leq\varepsilon .
\]
Following a proof of Weyl, choose
\[
  N_1<N_2<N_3<\cdots,
  \qquad
  N_p\geq N\!\left(\frac1{8^p}\right).
\]
Then
\[
  \int_\Omega
  |f_{N_{p+1}}(P)-f_{N_p}(P)|^2\,dv\leq\frac1{8^p}.
\]
Consequently, the set on which
\[
  |f_{N_{p+1}}(P)-f_{N_p}(P)|\geq\frac1{2^p}
\]
has measure at most \(2^{-p}\).  Hence
\[
  |f_{N_{p+1}}(P)-f_{N_p}(P)|\leq\frac1{2^p},\qquad
  |f_{N_{p+2}}(P)-f_{N_{p+1}}(P)|\leq\frac1{2^{p+1}},\ldots
\]
everywhere except on a set of measure at most
\[
  \frac1{2^p}+\frac1{2^{p+1}}+\cdots=\frac1{2^{p-1}}.
\]
Where these inequalities hold, the sequence
\(f_{N_p}(P),f_{N_{p+1}}(P),\ldots\) has a limit.  It therefore has a
limit everywhere except on a set of measure at most \(2^{1-p}\).
Since this is true for every \(p\), the exceptional set has measure
zero.

Define
\[
  f(P)=\lim_{p\to\infty}f_{N_p}(P)
\]
where this limit exists, that is, everywhere except on a null set, and
put \(f(P)=0\) otherwise.  If \(m\geq N(\varepsilon)\), then, for all
sufficiently large \(p\) (namely, as soon as
\(N_p\geq N(\varepsilon)\)),
\[
  \int_\Omega |f_m(P)-f_{N_p}(P)|^2\,dv\leq\varepsilon .
\]
We may therefore let \(p\to\infty\), obtaining
\[
  \int_\Omega |f_m(P)-f(P)|^2\,dv\leq\varepsilon .
\]
Thus \(f_m-f\), and hence \(f\), belongs to \(\mathfrak H'\), and
secondly \(f_m\to f\).  This proves everything asserted.

\section*{Appendix II. The Eigenvalue Problem for Unitary Operators}
\label{app:II}
\addcontentsline{toc}{section}{Appendix II. The Eigenvalue Problem for Unitary Operators}

\subsection*{1.}
\label{appII:1}
\addcontentsline{toc}{subsection}{Appendix II.1}

Throughout what follows we consider everywhere-defined, bounded
(continuous) linear operators \(R\) on \(\mathfrak H\); thus
\hyperref[thm:12]{Theorem 12(\(\alpha\)), (\(\beta\))} applies to them.  Hermitian
character, however, will not always be assumed.  For fixed \(f\), the
functional
\[
  L(g)=(f,Rg)
\]
has the properties
\[
  L(a_1g_1+\cdots+a_ng_n)
  =\overline{a_1}L(g_1)+\cdots+\overline{a_n}L(g_n),
  \qquad
  |L(g)|\leq C\,|f|\,|g|=D\,|g|.
\]
The theorem of F.~Riesz (cf.\ \hyperref[fn:52]{Note~52}) is therefore applicable:
\(L(g)=(f^*,g)\).  We denote the uniquely determined \(f^*\) by
\(R^*f\).  Then \(R^*\) is an everywhere-defined and plainly linear
operator, and
\[
  (f,Rg)=(R^*f,g),\qquad (Rf,g)=(f,R^*g).
\]
The first equation is the definition; the second follows by
interchanging \(f,g\) and taking complex conjugates.  In particular,
\[
  |R^*f|^2=(R^*f,R^*f)=(f,RR^*f)
  \leq C\,|f|\,|R^*f|,
  \qquad |R^*f|\leq C\,|f|,
\]
so \(R^*\) too is bounded.\footnote{\label{fn:64}For matrices, this star means taking
the conjugate transpose; for integral-operator kernels it has the
analogous meaning, and for differential operators it means taking the
adjoint.}

One verifies at once that
\begin{align*}
  R^{**}&=R,
  & (R\pm S)^*&=R^*\pm S^*,\\
  (RS)^*&=S^*R^*,
  & (aR)^*&=\overline a\,R^*,\\
  0^*&=0,
  & 1^*&=1.
\end{align*}
\footnote{\label{fn:65}The operators in question are everywhere defined, so the
operations \(+\), \(-\), and multiplication are meaningful without
further qualification; cf.\ \hyperref[fn:37]{Note~37}.}
For bounded Hermitian operators, \(R=R^*\) is characteristic.  Every
unitary operator \(U\) belongs to our class (it is bounded because
\(|Uf|=|f|\)), and within this class it is characterized by
\[
  UU^*=U^*U=1,\qquad\text{that is,}\qquad U^{-1}=U^*.
\]
\footnote{\label{fn:66}If \(U\) is unitary, it has an inverse \(U^{-1}\), and
\[
  (f,U^{-1}g)=(Uf,UU^{-1}g)=(Uf,g)=(f,U^*g)
\]
for all \(f,g\); hence \(U^{-1}=U^*\).  Conversely, if \(U^{-1}=U^*\),
then
\((Uf,Ug)=(f,U^*Ug)=(f,g)\), so that \(|Uf|=|f|\).  Thus \(U\) is
isometric; since it is everywhere defined and assumes every value
(\(U^{-1}\) is everywhere defined), it is unitary.}

\subsection*{2.}
\label{appII:2}
\addcontentsline{toc}{subsection}{Appendix II.2}

Now suppose that \(R\) is Hermitian.  We first attack its eigenvalue
problem by the method of F.~Riesz.  Let
\[
  p(x)=\sum_{\nu=0}^{n}a_\nu x^\nu
\]
be a polynomial with real coefficients.\translatornote{The source begins
this sum at \(\nu=1\), although the immediately following definition of
\(p(R)\) begins at \(\nu=0\) and explicitly includes \(R^0=1\).  The
lower limit has therefore been normalized to \(0\) here.}  Then
\[
  p(R)=\sum_{\nu=0}^{n}a_\nu R^\nu,\qquad R^0=1,
\]
is a Hermitian operator of our class.  Addition, subtraction, and
multiplication of these \(p(R)\) are performed exactly as for the
polynomials \(p(x)\); in particular, all of them commute.  Choose \(C\)
such that \(|Rf|\leq C|f|\) for every \(f\).

\begin{namedstatement}{Theorem}{\(1^*\)}
\label{thm:appII-1}
If \(p(x)\geq0\) for \(-C\leq x\leq C\), then \(p(R)\) is positive.
\end{namedstatement}

\begin{proof}
Let \(\mathfrak M\) be any finite-dimensional, and hence closed, linear
manifold, spanned for example by the orthonormal system
\(\varphi_1,\ldots,\varphi_k\).  The operator
\[
  R'=P_{\mathfrak M}RP_{\mathfrak M}
\]
is Hermitian and satisfies
\[
  |R'f|=|P_{\mathfrak M}RP_{\mathfrak M}f|
  \leq |RP_{\mathfrak M}f|
  \leq C|P_{\mathfrak M}f|
  \leq C|f|.
\]
It is therefore of our class and is plainly reduced by
\(\mathfrak M\) (its part in \(\mathfrak H-\mathfrak M\) is zero).
Its part in \(\mathfrak M\) is a \(k\)-dimensional Hermitian operator,
say with eigenvalues \(\lambda_1,\ldots,\lambda_k\).  Since
\(|R'f|\leq C|f|\), all \(|\lambda_j|\leq C\).  The operator \(p(R')\)
is also reduced by \(\mathfrak M\), and its part there has the
eigenvalues \(p(\lambda_1),\ldots,p(\lambda_k)\), all nonnegative.
Thus, if \(f\in\mathfrak M\),
\[
  (f,p(R')f)\geq0.
\]

Now let \(f\) be arbitrary, and take for \(\mathfrak M\) the closed
linear manifold spanned by \(f,Rf,\ldots,R^nf\).  Then \(f\in\mathfrak
M\), \(R'f=Rf,\ldots,R'^nf=R^nf\), and therefore
\(p(R')f=p(R)f\).  Hence \((f,p(R)f)\geq0\).  Since \(f\) was
arbitrary, \(p(R)\) is positive.
\end{proof}

\begin{namedstatement}{Theorem}{\(2^*\)}
\label{thm:appII-2}
If \(|p(x)|\leq\varepsilon\) for \(-C\leq x\leq C\), then
\[
  |p(R)f|\leq\varepsilon |f|.
\]
\end{namedstatement}

\begin{proof}
\hyperref[thm:appII-1]{Theorem \(1^*\)} applies to the polynomials
\(\varepsilon\pm p(x)\); hence the operators
\(\varepsilon 1\pm p(R)\) are positive, and
\[
  (f,\{\varepsilon1\pm p(R)\}f)\geq0,\qquad
  |(f,p(R)f)|\leq\varepsilon |f|^2.
\]
By \hyperref[thm:12]{Theorem 12} (the equivalence of \(\gamma\) and
\(\alpha\)), this means precisely that
\(|p(R)f|\leq\varepsilon|f|\).
\end{proof}

\begin{namedstatement}{Theorem}{\(3^*\)}
\label{thm:appII-3}
Let \(P(x)\) be continuous on \(-C\leq x\leq C\), and let
\(p_1(x),p_2(x),\ldots\) be polynomials converging there uniformly to
\(P(x)\).  Then \(p_n(R)f\) has a limit \(f^*\) depending only on
\(R,f\), and \(P(x)\).
\end{namedstatement}

\begin{proof}
This follows immediately from \hyperref[thm:appII-2]{Theorem \(2^*\)}.
\end{proof}

We therefore define \(P(R)f=f^*\).  This gives an everywhere-defined
linear Hermitian operator, because the \(p_n(R)\) have those
properties; it is bounded by \hyperref[thm:48]{Theorem 48(\(\alpha\)),
\((\gamma)\)}, or directly because the \(p_n(x)\), and hence by
\hyperref[thm:appII-2]{Theorem \(2^*\)} the \(p_n(R)\), are uniformly
bounded.

\begin{namedstatement}{Theorem}{\(4^*\)}
\label{thm:appII-4}
The \(P(R)\) are added, subtracted, and multiplied just as the
\(P(x)\) are.  Theorems \(1^*\) and \(2^*\) hold for them as well; and
when \(P(x)\) is a polynomial, the old and new definitions agree.
\end{namedstatement}

\begin{proof}
The last assertion follows by taking
\(p_1(x)=p_2(x)=\cdots=P(x)\).  The remaining assertions hold for the
polynomials \(p(x)\), and hence, by continuity, also for \(P(x)\).
\end{proof}

\subsection*{3.}
\label{appII:3}
\addcontentsline{toc}{subsection}{Appendix II.3}

We now depart from F.~Riesz's route.

\begin{namedstatement}{Theorem}{\(5^*\)}
\label{thm:appII-5}
Let \(\mathfrak M_R\) be the set of all \(f\) for which \(Rf=0\), and
let \(\mathfrak N_R\) be the set of all \(Rf\) and their limit points.
Then
\[
  \mathfrak M_R=\mathfrak H-\mathfrak N_R.
\]
\end{namedstatement}

\begin{proof}
The vectors \(Rf\) form a linear manifold.  It is enough to show that
\(\mathfrak M_R\) consists of all \(g\) orthogonal to this manifold.
Orthogonality to it means
\[
  (g,Rf)=0\quad\text{for every }f,
\]
that is, \((Rg,f)=0\), and hence \(Rg=0\).
\end{proof}

\begin{namedstatement}{Theorem}{\(6^*\)}
\label{thm:appII-6}
If \(R,S\) are positive, then \(\mathfrak M_{R+S}\) is a part of
\(\mathfrak M_R\), and \(\mathfrak N_R\) is a part of
\(\mathfrak N_{R+S}\).
\end{namedstatement}

\begin{proof}
From \((R+S)g=0\) it follows that
\[
  0\leq(g,Rg)\leq(g,(R+S)g)=0,
\]
and hence \((g,Rg)=0\).  For positive \(R\) one has
\[
  |(f,Rg)|^2\leq(f,Rf)(g,Rg).
\]
This is proved just like the corresponding inequality for \(R=1\) in
\hyperref[thm:1]{Theorem 1}.  Thus \((g,Rg)=0\) implies
\((f,Rg)=0\) for every \(f\), and therefore \(Rg=0\).  The preceding
relation says that \(\mathfrak M_{R+S}\) is a part of
\(\mathfrak M_R\); \hyperref[thm:appII-5]{Theorem \(5^*\)} then shows
that \(\mathfrak N_R\) is a part of \(\mathfrak N_{R+S}\).
\end{proof}

We define
\[
  f_a(x)=\operatorname{Max}(x-a,0),\qquad
  \mathfrak M(R,a)=\mathfrak M_{f_a(R)},\qquad
  \mathfrak N(R,a)=\mathfrak N_{f_a(R)}.
\]

\begin{namedstatement}{Theorem}{\(7^*\)}
\label{thm:appII-7}
Throughout \(\mathfrak M(R,a)\), respectively \(\mathfrak N(R,a)\),
one has
\[
  (f,Rf)\leq a|f|^2,
  \qquad\text{respectively}\qquad
  (f,Rf)\geq a|f|^2.
\]
\end{namedstatement}

\begin{proof}
Put \(g_a(x)=\operatorname{Max}(-(x-a),0)\).  Since
\[
  f_a(x)g_a(x)=0,\qquad f_a(x)-g_a(x)=x-a,
\]
we have, with \(S'=f_a(R)\) and \(S''=g_a(R)\),
\[
  S'S''=S''S'=0,\qquad S'-S''=R-a1
\]
by \hyperref[thm:appII-4]{Theorem \(4^*\)}.  Both \(S'\) and \(S''\)
are positive.  In
\(\mathfrak M(R,a)=\mathfrak M_{S'}\) we have \(S'f=0\), so
\[
  Rf=af-S''f,\qquad
  (f,Rf)=a|f|^2-(f,S''f)\leq a|f|^2.
\]
In \(\mathfrak N(R,a)=\mathfrak N_{S'}\), by
\hyperref[thm:appII-5]{Theorem \(5^*\)} and \(S''S'=0\), we lie in
\(\mathfrak M_{S''}\).  Thus \(S''f=0\), and
\[
  Rf=af+S'f,\qquad
  (f,Rf)=a|f|^2+(f,S'f)\geq a|f|^2.
\]
\end{proof}

\begin{namedstatement}{Theorem}{\(8^*\)}
\label{thm:appII-8}
If \(a\leq b\), then \(\mathfrak M(R,a)\) is a part of
\(\mathfrak M(R,b)\), and \(\mathfrak N(R,b)\) is a part of
\(\mathfrak N(R,a)\).  Moreover, for \(a<-C\), respectively \(a>C\),
\[
  \mathfrak M(R,a)=(0),\quad \mathfrak M(R,a)=\mathfrak H,
  \qquad
  \mathfrak N(R,a)=\mathfrak H,\quad \mathfrak N(R,a)=(0),
\]
respectively.
\end{namedstatement}

\begin{proof}
Since \(f_b(x)\geq0\) and \(f_a(x)-f_b(x)\geq0\), the first assertion
follows from \hyperref[thm:appII-6]{Theorem \(6^*\)}, applied to
\(f_b(R)\) and \(f_a(R)-f_b(R)\).\translatornote{The source prints
\(f_b-f_a\geq0\) and applies the theorem to \(f_a\) and \(f_b-f_a\).
Since \(a\leq b\) gives \(f_a\geq f_b\), both signs have been
corrected.}  For the second, if a nonzero \(f\)
belonged, when \(a<-C\), to \(\mathfrak M(R,a)\), or, when \(a>C\),
to \(\mathfrak N(R,a)\), then
\[
  (f,Rf)\leq a|f|^2<-C|f|^2,
  \qquad\text{respectively}\qquad
  (f,Rf)\geq a|f|^2>C|f|^2,
\]
contrary to the definition of \(C\) in
\hyperref[thm:12]{Theorem 12(\(\gamma\))}.  Hence the corresponding
manifold is \((0)\), and its orthogonal complement is
\(\mathfrak H-(0)=\mathfrak H\).
\end{proof}

\begin{namedstatement}{Theorem}{\(9^*\)}
\label{thm:appII-9}
Quite generally,
\[
  (Rf,g)=\int a\,d\bigl(P_{\mathfrak M(R,a)}f,g\bigr),
\]
where the integral may be taken between any two limits smaller than
\(-C\) and greater than \(C\).
\end{namedstatement}

\begin{proof}
For \(f=g\), the assertion follows from
\hyperref[thm:appII-7]{Theorems \(7^*\)} and
\hyperref[thm:appII-8]{\(8^*\)}.\footnote{\label{fn:67}The operator \(R\) commutes
with \(f_a(R)\), by Theorem \(4^*\), and is therefore reduced by
\(\mathfrak M(R,a)\): from \(f_a(R)f=0\) follows
\(f_a(R)Rf=Rf_a(R)f=0\), and one applies \hyperref[thm:20]{Theorem 20}.  Thus \(R\)
commutes with \(P_{\mathfrak M(R,a)}\).

Let
\[
  x_0<x_1<\cdots<x_p,\qquad x_0<-C,\quad x_p>C.
\]
Then
\begin{align*}
 (Rf,f)
 &=\sum_{\nu=1}^{p}
   \bigl(R(P_{\mathfrak M(R,x_\nu)}
      -P_{\mathfrak M(R,x_{\nu-1})})f,f\bigr)\\
 &=\sum_{\nu=1}^{p}
   \bigl(RP_{\mathfrak M(R,x_\nu)-\mathfrak M(R,x_{\nu-1})}f,f\bigr)\\
 &=\sum_{\nu=1}^{p}
   \bigl(RP_{\mathfrak M(R,x_\nu)\mathfrak N(R,x_{\nu-1})}f,f\bigr)\\
 &=\sum_{\nu=1}^{p}
   \bigl(RP_{\mathfrak M(R,x_\nu)\mathfrak N(R,x_{\nu-1})}f,
         P_{\mathfrak M(R,x_\nu)\mathfrak N(R,x_{\nu-1})}f\bigr).
\end{align*}
The last step uses the commutativity just noted.  By Theorem \(7^*\),
each summand lies between
\[
 x_{\nu-1}
 |P_{\mathfrak M(R,x_\nu)-\mathfrak M(R,x_{\nu-1})}f|^2
 \quad\text{and}\quad
 x_\nu
 |P_{\mathfrak M(R,x_\nu)-\mathfrak M(R,x_{\nu-1})}f|^2,
\]
and
\[
 |P_{\mathfrak M(R,x_\nu)-\mathfrak M(R,x_{\nu-1})}f|^2
 =|P_{\mathfrak M(R,x_\nu)}f|^2
  -|P_{\mathfrak M(R,x_{\nu-1})}f|^2.
\]
Consequently,
\begin{align*}
 (Rf,f)
 &\leq\sum_{\nu=1}^{p}x_\nu
 \bigl(
 |P_{\mathfrak M(R,x_\nu)}f|^2
 -|P_{\mathfrak M(R,x_{\nu-1})}f|^2
 \bigr),\\
 (Rf,f)
 &\geq\sum_{\nu=1}^{p}x_{\nu-1}
 \bigl(
 |P_{\mathfrak M(R,x_\nu)}f|^2
 -|P_{\mathfrak M(R,x_{\nu-1})}f|^2
 \bigr).
\end{align*}
These are the upper and lower Stieltjes sums.  By the definition of
the Stieltjes integral, they yield
\[
  (Rf,f)=\int a\,d|P_{\mathfrak M(R,a)}f|^2
        =\int a\,d(P_{\mathfrak M(R,a)}f,f).
\]}
Substituting \((f+g)/2\) and \((f-g)/2\) and subtracting gives the
equation for the real parts in general; replacing \(f,g\) by \(if,g\)
then gives the imaginary parts.
\end{proof}

Theorem \(9^*\) is essentially Hilbert's result for bounded Hermitian
operators.  We proceed to two Hermitian operators \(R,S\) of our class.

\begin{namedstatement}{Theorem}{\(10^*\)}
\label{thm:appII-10}
If \(R,S\) commute, then so do
\[
  P_{\mathfrak M(R,a)}
  \quad\text{and}\quad
  P_{\mathfrak M(S,b)}.
\]
\end{namedstatement}

\begin{proof}
Since \(R,S\) commute, so do all \(P(R),Q(S)\): this is immediate for
polynomials and follows by continuity for general continuous functions.
In particular, \(f_a(R)\) commutes with \(f_b(S)\).  Hence
\(\mathfrak M(R,a)\) reduces \(f_b(S)\), by the argument in the preceding
footnote, and so does
\(\mathfrak N(R,a)=\mathfrak H-\mathfrak M(R,a)\).
Thus \(\mathfrak M(S,b)\) is obtained by first taking, within
\(\mathfrak M(R,a)\), the \(f\) for which \(f_b(S)f=0\), then doing the
same within \(\mathfrak N(R,a)\), and finally adding the two closed
linear manifolds.  The commutativity of the two projection operators
follows at once.
\end{proof}

\subsection*{4.}
\label{appII:4}
\addcontentsline{toc}{subsection}{Appendix II.4}

We now turn to the so-called normal operators \(A\) of our class,
characterized by
\[
  AA^*=A^*A.
\]
The Hermitian operators
\[
  R=\frac12(A+A^*),\qquad
  S=\frac1{2i}(A-A^*)
\]
then commute, and
\[
  A=R+iS,\qquad A^*=R-iS.
\]
Form the projection operators
\[
  E(A,a,b)
  =P_{\mathfrak M(R,a)}P_{\mathfrak M(S,b)}.
\]

\begin{namedstatement}{Theorem}{\(11^*\)}
\label{thm:appII-11}
The operator \(A\) and the \(E(A,a,b)\) satisfy the following
conditions:\footnote{\label{fn:68}Cf.\ \hyperref[intro:IX]{Introduction IX} and
\hyperref[fn:58]{Note~58}.}

\(\alpha\) The operators \(E(A,a,b)\) and \(E(A,a',b')\) commute, and
\[
  E(A,a,b)E(A,a',b')
  =E\bigl(A,\min(a,a'),\min(b,b')\bigr).
\]

\(\beta\) For sufficiently small \(a\) or \(b\),
\(E(A,a,b)=0\).  For sufficiently large \(a\), respectively \(b\), it
is independent of \(a\), respectively \(b\); and for sufficiently
large \(a,b\), it equals \(1\).

Moreover,
\begin{align*}
  (Af,g)&=\iint(a+ib)\,d\bigl(E(A,a,b)f,g\bigr),\\
  (A^*f,g)&=\iint(a-ib)\,d\bigl(E(A,a,b)f,g\bigr),
\end{align*}
where the double integrals extend over a sufficiently large rectangle.
\end{namedstatement}

\begin{proof}
By \hyperref[thm:appII-9]{Theorems \(9^*\)} and
\hyperref[thm:appII-10]{\(10^*\)},
\[
  (Rf,g)=\iint a\,d\bigl(E(A,a,b)f,g\bigr),\qquad
  (Sf,g)=\iint b\,d\bigl(E(A,a,b)f,g\bigr).
\]
Everything now follows immediately from
\hyperref[thm:appII-8]{Theorem \(8^*\)},
\hyperref[thm:appII-9]{Theorem \(9^*\)},
\hyperref[thm:appII-10]{Theorem \(10^*\)}, and the definition of
\(E(A,a,b)\).
\end{proof}

Unitary operators \(U\) are normal, since \(UU^*=U^*U=1\), so
\hyperref[thm:appII-11]{Theorem \(11^*\)} applies to them.  Their
\(E(U,a,b)\), however, have further properties, which we now derive.

First, quite generally,
\begin{align*}
 (Af,E(A,a,b)g)
 &=\iint(a'+ib')\,d\bigl(E(A,a',b')f,E(A,a,b)g\bigr)\\
 &=\iint(a'+ib')\,d\bigl(E(A,a,b)E(A,a',b')f,g\bigr)\\
 &=\iint(a'+ib')\,d\bigl(
     E(A,\min(a,a'),\min(b,b'))f,g\bigr)\\
 &=\int_{-\infty}^{a}\int_{-\infty}^{b}
      (a'+ib')\,d\bigl(E(A,a',b')f,g\bigr),
\end{align*}
and
\begin{align*}
 (Af,Ag)
 &=\iint(a-ib)\,d\bigl(Af,E(A,a,b)g\bigr)\\
 &=\iint(a-ib)\,
    d\!\left\{\int_{-\infty}^{a}\int_{-\infty}^{b}
       (a'+ib')\,d\bigl(E(A,a',b')f,g\bigr)\right\}\\
 &=\iint(a-ib)(a+ib)\,d\bigl(E(A,a,b)f,g\bigr)\\
 &=\iint(a^2+b^2)\,d\bigl(E(A,a,b)f,g\bigr),
\end{align*}
where the interchange in the second calculation is as in \hyperref[fn:55]{Note~55}.
Also
\[
  (f,g)=\iint d\bigl(E(A,a,b)f,g\bigr).
\]
For unitary \(A\), therefore, \((f,g)=(Af,Ag)\) gives
\[
  \iint(1-a^2-b^2)\,d\bigl(E(A,a,b)f,g\bigr)=0,
\]
and, on putting \(f=g\),
\[
  \iint(1-a^2-b^2)\,d|E(A,a,b)f|^2=0.
\]

\subsection*{5.}
\label{appII:5}
\addcontentsline{toc}{subsection}{Appendix II.5}

If \(\mathcal R\) is a rectangle
\[
  a'\leq x\leq a'',\qquad b'\leq y\leq b''
\]
with vertices
\((a',b'),(a'',b'),(a',b''),(a'',b'')\), define the projection
operator\translatornote{The source reverses \(\mathfrak M\) and
\(\mathfrak N\) in this rectangle projection, the following
parenthesis, and Note~69.  Since \(E(A,a,b)\) is increasing in each
endpoint, the interval factors are
\(\mathfrak N(R,a')\mathfrak M(R,a'')\) and
\(\mathfrak N(S,b')\mathfrak M(S,b'')\), as corrected here.}
\begin{align*}
 E(A,\mathcal R)
 &=E(A,a'',b'')-E(A,a'',b')-E(A,a',b'')+E(A,a',b')\\
 &=P_{\mathfrak N(R,a')}\bigl(1-P_{\mathfrak N(R,a'')}\bigr)
   P_{\mathfrak N(S,b')}\bigl(1-P_{\mathfrak N(S,b'')}\bigr)\\
 &=P_{\mathfrak N(R,a')\mathfrak M(R,a'')
                   \mathfrak N(S,b')\mathfrak M(S,b'')}.
\end{align*}
Suppose that \(\mathcal R\) has no point in common with the unit circle.
Then throughout \(\mathcal R\) either
\(1-a^2-b^2\geq\varepsilon\), or
\(1-a^2-b^2\leq-\varepsilon\), for some \(\varepsilon>0\).
For every \(f\) such that \(E(A,\mathcal R)f=f\) (that is, \(f\) lies
in
\(\mathfrak N(R,a')\mathfrak M(R,a'')
  \mathfrak N(S,b')\mathfrak M(S,b'')\))\footnote{\label{fn:69}Since \(f\) lies in both
\(\mathfrak N(R,a')\) and \(\mathfrak M(R,a'')\),
\(P_{\mathfrak M(R,a)}f\) equals \(0\) for \(a\leq a'\) and \(f\) for
\(a\geq a''\); thus \(E(A,a,b)f\) is independent of \(a\) outside this
interval.  The same is true of \(b\) outside \(b'\leq b\leq b''\).
Accordingly, the integral over the whole plane may be identified with
the integral over \(\mathcal R\).}, \mbox{we therefore have}
\begin{align*}
 \left|\iint_{\mathcal R}(1-a^2-b^2)\,
           d|E(A,a,b)f|^2\right|
 &=\left|\iint(1-a^2-b^2)\,d|E(A,a,b)f|^2\right|\\
 &\geq\varepsilon
   \left|\iint d|E(A,a,b)f|^2\right|
  =\varepsilon|f|^2.
\end{align*}
Thus \(f=0\), and \(E(A,\mathcal R)=0\).

If \(\mathcal R,\mathcal S\) are disjoint rectangles (their boundaries
may meet),\footnote{\label{fn:70}Write, for example,
\[
 \mathcal R:\ a'\leq x\leq a'',\ b'\leq y\leq b'',\qquad
 \mathcal S:\ c'\leq x\leq c'',\ d'\leq y\leq d''.
\]
One of \(a''\leq c'\), \(c''\leq a'\), \(b''\leq d'\), or
\(d''\leq b'\) must hold.  In the first case
\(\mathfrak M(R,a'')\) is a part of \(\mathfrak M(R,c')\), so
\(\mathfrak M(R,a'')\) is orthogonal to \(\mathfrak N(R,c')\).  Hence
\(P_{\mathfrak M(R,a'')}P_{\mathfrak N(R,c')}=0\), and a fortiori
\(E(A,\mathcal R)E(A,\mathcal S)=0\).  The other three cases are
identical.}\textsuperscript{,}\translatornote{In Note 70 the source instead states
\(\mathfrak N(R,a'')\perp\mathfrak M(R,c')\) and writes the
corresponding product of projections.  The complementary factors have
been corrected here; the corrected orthogonality is the one needed for
the conclusion.} then
\[
  E(A,\mathcal R)E(A,\mathcal S)=0.
\]
If two rectangles \(\mathcal R,\mathcal S\) together make exactly one
rectangle \(\mathcal T\),\footnote{\label{fn:71}Either \(a''=c'\) or
\(c''=a'\), while simultaneously \(b'=d'\) and \(b''=d''\); or
\(b''=d'\) or \(d''=b'\), while simultaneously \(a'=c'\) and
\(a''=c''\).  The assertion follows directly from the definition of
\(E(A,\mathcal R)\).} then
\[
  E(A,\mathcal R)+E(A,\mathcal S)=E(A,\mathcal T).
\]

Repeated use of these facts yields the following.  If \(\mathcal U\)
is a union of finitely many rectangles
\(\mathcal R_1,\ldots,\mathcal R_n\), then
\[
  E(A,\mathcal U)
  =E(A,\mathcal R_1)+\cdots+E(A,\mathcal R_n)
\]
is a projection operator; it is independent of the rectangular
decomposition of \(\mathcal U\).  If \(\mathcal U,\mathcal V\) are two
such sets having at most boundary points in common, then
\[
  E(A,\mathcal U)E(A,\mathcal V)=0,\qquad
  E(A,\mathcal U+\mathcal V)
  =E(A,\mathcal U)+E(A,\mathcal V).
\]
Moreover, if \(\mathcal U\) does not meet the unit circle, then
\(E(A,\mathcal U)=0\).  It follows easily that if
\(\mathcal U,\mathcal V\) have the same points in common with the unit
circle, and neither has a corner on that circle, then
\[
  E(A,\mathcal U)=E(A,\mathcal V).
\]
We can now, so to speak, transform the Stieltjes double integral to
polar coordinates.

\subsection*{6.}
\label{appII:6}
\addcontentsline{toc}{subsection}{Appendix II.6}

Let \(0\leq\rho<\sigma\leq1\).  There certainly exist finite unions of
rectangles \(\mathcal U\) whose intersection with the unit circle is
exactly the arc from \(2\pi\rho\) to \(2\pi\sigma\), and none of whose
corners lies on the unit circle.  For all such \(\mathcal U\),
\(E(A,\mathcal U)\) is the same projection operator; denote it by
\(F(\rho,\sigma)\).  For \(0\leq\rho<\sigma<\tau\leq1\),
\[
  F(\rho,\sigma)+F(\sigma,\tau)=F(\rho,\tau).
\]
Putting \(F(\rho)=F(0,\rho)\), we therefore have
\[
  F(\rho,\sigma)=F(\sigma)-F(\rho).
\]
Since plainly \(F(0,1)=1\), it follows that \(F(1)=1\).  The \(F(\rho)\)
are projection operators, and because \(F(\sigma)-F(\rho)\) is a
projection operator for \(\rho\leq\sigma\), \(F(\rho)\) is a part of
\(F(\sigma)\).  Define the hitherto undefined \(F(0)\) to be \(0\).

Partition the integration region in
\[
  (Af,f)=\iint(a+ib)\,d|E(A,a,b)f|^2
\]
(a large rectangle) into rectangles
\(\mathcal R_1,\ldots,\mathcal R_k\) as follows.  Each has diameter at
most \(\varepsilon>0\), and none has a corner on the unit circle.
Order them so that \(\mathcal R_1,\ldots,\mathcal R_h\), where
\(h\leq k\), meet the circle in the respective arcs
\[
  2\pi\rho_0,2\pi\rho_1;\ \ldots;\
  2\pi\rho_{h-1},2\pi\rho_h,
  \qquad
  0=\rho_0<\rho_1<\cdots<\rho_h=1.
\]
Choose a point \(\xi_\nu\) from each \(\mathcal R_\nu\).  Then
\[
  (Af,f)
  =\sum_{\nu=1}^{k}\iint_{\mathcal R_\nu}
      (a+ib)\,d|E(A,a,b)f|^2
\]
differs by an arbitrarily small amount\footnote{\label{fn:72}For every rectangle
\(\mathcal S\),
\(\iint_{\mathcal S}d|E(A,a,b)f|^2\geq0\), since it equals
\(|E(A,\mathcal S)f|^2\).  Thus the difference is at most
\[
  \varepsilon\iint d|E(A,a,b)f|^2
  =\varepsilon|f|^2.
\]}
from
\begin{align*}
 \sum_{\nu=1}^{k}\xi_\nu
    \iint_{\mathcal R_\nu}d|E(A,a,b)f|^2
 &=\sum_{\nu=1}^{k}\xi_\nu(E(A,\mathcal R_\nu)f,f)\\
 &=\sum_{\nu=1}^{h}\xi_\nu|F(\rho_{\nu-1},\rho_\nu)f|^2\\
 &=\sum_{\nu=1}^{h}\xi_\nu
   \bigl(|F(\rho_\nu)f|^2-|F(\rho_{\nu-1})f|^2\bigr).
\end{align*}
We may choose
\[
  \xi_\nu=e^{2\pi i\rho_\nu'},
  \qquad
  \rho_{\nu-1}<\rho_\nu'<\rho_\nu
  \quad(\nu=1,\ldots,h),
\]
so the last expression becomes
\[
  \sum_{\nu=1}^{h}e^{2\pi i\rho_\nu'}
   \bigl(|F(\rho_\nu)f|^2-|F(\rho_{\nu-1})f|^2\bigr).
\]
The points \(\rho_0,\rho_1,\ldots,\rho_h\) are at most
\(\varepsilon\) apart.  Hence
\[
  (Af,f)=\int_0^1e^{2\pi i\rho}\,d|F(\rho)f|^2.
\]
\footnote{\label{fn:73}The integral on the right converges; for example, prescribe
\(\rho_0,\rho_1,\ldots,\rho_h\) and adapt the rectangular partition to
them.}

\begin{namedstatement}{Theorem}{\(12^*\)}
\label{thm:appII-12}
For every unitary operator \(U\) there is a family of projection
operators \(E(\lambda)\), \(0\leq\lambda\leq1\), with the following
properties:

\(\alpha\) If \(\lambda\leq\mu\), then \(E(\lambda)\) is a part of
\(E(\mu)\); moreover \(E(0)=0\) and \(E(1)=1\).

\(\beta\) If \(\lambda\geq\lambda_0\) and
\(\lambda\to\lambda_0\), then, for every \(f\),
\[
  E(\lambda)f\longrightarrow E(\lambda_0)f.
\]

For all \(f,g\),
\[
  (Uf,g)=\int_0^1e^{2\pi i\lambda}\,
     d\bigl(E(\lambda)f,g\bigr).
\]
\end{namedstatement}

\begin{proof}
Consider the \(F(\lambda)\) just constructed.  They have property
\(\alpha\), but may fail to have \(\beta\).  Let
\[
  0\leq\lambda_0<1,\qquad
  \lambda_1>\lambda_2>\cdots>\lambda_0,\qquad
  \lambda_n\to\lambda_0.
\]
The \(F(\lambda_n)\) form a decreasing sequence of projection
operators.  By \hyperref[thm:19]{Theorem 19}, there is a projection
operator \(E\) such that
\[
  F(\lambda_n)f\longrightarrow Ef
\]
for every \(f\).  This \(E\) depends only on \(\lambda_0\): for any two
such sequences, suitable subsequences can be interlaced to give
another sequence of the same kind, for which the convergence must also
hold.  Denote this \(E\) by \(E(\lambda_0)\).

Since \(F(\lambda_0)\) is a part of every
\(F(\lambda_1),F(\lambda_2),\ldots\), it is a part of
\(E(\lambda_0)\).  If \(\lambda_0<\mu_0\), choose
\(\lambda_1=\mu_0\); then \(E(\lambda_0)\) is a part of
\(F(\mu_0)=F(\lambda_1)\), since it is a part of all the
\(F(\lambda_n)\).  Thus
\[
  F(\lambda)\ \text{is a part of}\ E(\lambda),\qquad
  E(\lambda)\ \text{is a part of}\ F(\mu)\quad(\lambda<\mu),
\]
and consequently \(E(\lambda)\) is a part of \(E(\mu)\) for
\(\lambda\leq\mu\).

Now let \(\lambda>\lambda_0\), \(\lambda\to\lambda_0\).  The projection
\(E(\lambda)-E(\lambda_0)\) is a part of
\(F(\lambda)-E(\lambda_0)\).  Hence
\[
 |E(\lambda)f-E(\lambda_0)f|
 \leq|F(\lambda)f-E(\lambda_0)f|,
\]
and the right-hand side tends to zero by the definition of
\(E(\lambda_0)\).  Thus the \(E(\lambda)\) satisfy \(\beta\) and also
\(\alpha\), apart from the requirements \(E(0)=0\), \(E(1)=1\).  We
enforce those by definition; \(\beta\) may then fail at \(\lambda=0\).
Denote the old value of \(E(0)\) by \(\overline E\).

For \(F(\lambda)\) we had
\[
  (Uf,f)=\int_0^1 e^{2\pi i\lambda}\,d|F(\lambda)f|^2.
\]
The functions \(|F(\lambda)f|^2\) and \(|E(\lambda)f|^2\) are monotone
in \(\lambda\), and, for \(\lambda<\mu\),
\[
  |F(\lambda)f|^2
  \leq|E(\lambda)f|^2
  \leq|F(\mu)f|^2.
\]
It follows readily that the same integral relation holds with
\(E(\lambda)\).  This is the asserted final relation when \(f=g\).
Substitution of \((f+g)/2\) and \((f-g)/2\), followed by subtraction,
gives its real part for arbitrary \(f,g\); replacing \(f,g\) by \(if,g\)
gives its imaginary part.

It remains to repair the possible failure of \(\beta\) at
\(\lambda=0\).  For \(0<\lambda<1\), replace \(E(\lambda)\) by
\(E(\lambda)-\overline E\).  Everything else remains unchanged and
this last point is put in order.
\end{proof}

\subsection*{7.}
\label{appII:7}
\addcontentsline{toc}{subsection}{Appendix II.7}

\hyperref[thm:appII-12]{Theorem \(12^*\)} was the principal
difficulty.  Some easy uniqueness theorems now follow.

\begin{namedstatement}{Theorem}{\(13^*\)}
\label{thm:appII-13}
If projection operators \(E(\lambda)\) are chosen in accordance with
conditions \(\alpha,\beta\) of \hyperref[thm:appII-12]{Theorem \(12^*\)}, then there is exactly
one unitary operator \(U\) related to them as in that theorem.
\end{namedstatement}

\begin{proof}
There is at most one such \(U\), since
\hyperref[thm:appII-12]{Theorem \(12^*\)} determines every
\((Uf,g)\), and hence every \(Uf\).  It remains to prove existence.
The argument used in the proof of \hyperref[thm:36]{Theorem 36}, and
\(|e^{2\pi i\lambda}|=1\), give
\[
  \left|
  \int_0^1e^{2\pi i\lambda}\,d(E(\lambda)f,g)
  \right|
  \leq |f|\,|g|.
\]
Call the integral on the left \(L(g)\), with \(f\) fixed and \(g\)
variable.  Then
\[
  L(a_1g_1+\cdots+a_ng_n)
  =\overline{a_1}L(g_1)+\cdots+\overline{a_n}L(g_n),
  \qquad
  |L(g)|\leq|f|\,|g|=C|g|.
\]
F.~Riesz's theorem applies (cf.\ \hyperref[fn:53]{Note~53}): there is exactly one \(f^*\)
such that \(L(g)=(f^*,g)\).  Define \(U\) by \(Uf=f^*\).  This is plainly
an everywhere-defined linear operator; only its unitarity remains to be
proved.

First \(|(Uf,g)|\leq|f||g|\), so \(U\) is bounded and \(U^*\) exists.
We have
\[
  (Uf,g)=\int_0^1e^{2\pi i\lambda}\,d(E(\lambda)f,g),
  \qquad
  (U^*f,g)=\int_0^1e^{-2\pi i\lambda}\,d(E(\lambda)f,g).
\]
The first equation is the definition; the second follows by
interchanging \(f,g\) and taking complex conjugates.  Just as at the
end of \hyperref[appII:4]{Section 4}, these formulas give
\[
  (Uf,Ug)=(f,g),\qquad (U^*f,U^*g)=(f,g).
\]
Consequently
\[
  (U^*Uf,g)=(UU^*f,g)=(f,g)
\]
for all \(f,g\), so \(U^*U=UU^*=1\); that is, \(U\) is unitary.
\end{proof}

\begin{namedstatement}{Theorem}{\(14^*\)}
\label{thm:appII-14}
For a given unitary operator \(U\), only one family \(E(\lambda)\) of
projection operators can be related to it as in Theorem \(12^*\).
\end{namedstatement}

\begin{proof}
The formula
\[
  (U^nf,g)=\int_0^1e^{2\pi ni\lambda}\,
             d(E(\lambda)f,g)
\]
is trivial for \(n=0\), is the definition for \(n=1\), and follows from
\(U^{-1}=U^*\) for \(n=-1\).  If it has been proved for \(m,n\), then it
also holds for \(m-n\), since the method used at the end of
\hyperref[appII:4]{Section 4} gives
\[
  (U^mf,U^ng)
  =\int_0^1e^{2(m-n)\pi i\lambda}\,d(E(\lambda)f,g),
\]
whereas
\[
  (U^mf,U^ng)
  =((U^n)^*U^mf,g)
  =(U^{*n}U^mf,g)
  =(U^{-n}U^mf,g)
  =(U^{m-n}f,g).
\]
Thus the formula holds for every
\(n=0,\pm1,\pm2,\ldots\).

If \(F(\lambda)\) is a second family with the same properties, then the
same equation holds with \(F\) in place of \(E\).  Hence
\[
  \int_0^1p(\lambda)\,
    d\bigl((E(\lambda)f,g)-(F(\lambda)f,g)\bigr)=0
\]
first for every \(p(\lambda)=e^{2n\pi i\lambda}\), and then for their
linear combinations, the trigonometric polynomials.  By continuity it
holds for every continuous \(p(\lambda)\) with \(p(0)=p(1)\).  Because
of condition \(\beta\) in \hyperref[thm:appII-12]{Theorem \(12^*\)}, it remains valid for
functions having finitely many jumps, provided the value at a jump is
always taken on its right.  Put, for \(0\leq\lambda_0<1\),
\[
  p(\lambda)=
  \begin{cases}
    1,&0<\lambda\leq\lambda_0,\\
    0,&\text{otherwise}.
  \end{cases}
\]
Then
\[
 \int_0^{\lambda_0}
  d\bigl((E(\lambda)f,g)-(F(\lambda)f,g)\bigr)
 =(E(\lambda_0)f,g)-(F(\lambda_0)f,g)=0.
\]
Since this holds for all \(f,g\),
\(E(\lambda_0)=F(\lambda_0)\).  Here \(\lambda_0\neq1\); at
\(\lambda_0=1\) the assertion follows from condition \(\alpha\) of
\hyperref[thm:appII-12]{Theorem \(12^*\)}.
\end{proof}

\subsection*{8.}
\label{appII:8}
\addcontentsline{toc}{subsection}{Appendix II.8}

Theorems \(12^*,13^*,14^*\) show that the correspondence in
Theorem \(12^*\) between a unitary operator \(U\) and a family of
projection operators \(E(\lambda)\) satisfying \(\alpha,\beta\) is
one-to-one, and exhausts all \(U\) and all \(E(\lambda)\).  To be a
resolution of the identity, these \(E(\lambda)\) still require
\[
  E(\lambda)f\longrightarrow f
  \qquad(\lambda<1,\ \lambda\to1);
\]
cf.\ \hyperref[def:17]{Definition 17}.  We shall therefore have proved
everything announced in \hyperref[chap:IX]{Chapter IX} once we show
that this property is equivalent to the assertion that
\(\varphi=U\varphi\) occurs only for \(\varphi=0\).

First suppose \(\varphi=U\varphi\), \(\varphi\neq0\).  Then
\begin{align*}
 \Re\bigl((\varphi,\varphi)-(U\varphi,\varphi)\bigr)
 &=\Re\int_0^1(1-e^{2\pi i\lambda})\,
             d(E(\lambda)\varphi,\varphi)\\
 &=\int_0^1 2\sin^2(\pi\lambda)\,
             d|E(\lambda)\varphi|^2.
\end{align*}
The left-hand side vanishes, while, for every \(\varepsilon>0\), the
right-hand side is at least
\begin{align*}
 \int_{\varepsilon}^{1-\varepsilon}
   2\sin^2(\pi\lambda)\,d|E(\lambda)\varphi|^2
 &\geq2\sin^2(\pi\varepsilon)
    \int_{\varepsilon}^{1-\varepsilon}
      d|E(\lambda)\varphi|^2\\
 &=2\sin^2(\pi\varepsilon)
   \bigl(|E(1-\varepsilon)\varphi|^2
         -|E(\varepsilon)\varphi|^2\bigr).
\end{align*}
For \(\varepsilon\leq\tfrac12\), the last expression is nonnegative,
and therefore zero:
\[
  |E(1-\varepsilon)\varphi|=|E(\varepsilon)\varphi|.
\]
As \(\varepsilon\to0\), the right-hand side tends to zero, and hence so
does the left.  Thus
\[
  E(1-\varepsilon)\varphi\longrightarrow0\neq\varphi.
\]

Conversely, suppose \(E(\lambda)f\not\to f\) as
\(\lambda<1,\lambda\to1\).  By \hyperref[thm:19]{Theorem 19}, there is
a projection operator \(E\) such that \(E(\lambda)g\to Eg\) for every
\(g\); in particular \(Ef\neq f\).  For \(f'=f-Ef\), therefore,
\[
  f'\neq0,\qquad E(\lambda)f'\longrightarrow Ef'=0.
\]
Now \(|E(\lambda)f'|\) is nonnegative and monotonically nondecreasing
in \(\lambda\).  Since it tends to zero as \(\lambda\to1\), it must
vanish identically.  Hence
\[
  E(\lambda)f'=0\quad(\lambda<1),\qquad E(1)f'=f'.
\]
It follows that
\[
  (Uf',g)=\int_0^1e^{2\pi i\lambda}\,
             d(E(\lambda)f',g)=(f',g)
\]
for every \(g\), and therefore \(Uf'=f'\).  This completes the proof.

\section*{Appendix III. Operators and Matrices}
\label{app:III}
\addcontentsline{toc}{section}{Appendix III. Operators and Matrices}

\subsection*{1.}
\label{appIII:1}
\addcontentsline{toc}{subsection}{Appendix III.1}

It is desirable to translate our results on Hermitian operators into
the customary terminology of matrices.  We shall see that, in the
unbounded case, the terminology of operators is considerably more
practical.  In the bounded case, by Hilbert's results, the two amount to
the same thing.  We shall also see that closed linear Hermitian
operators correspond to Hermitian matrices whose row sums of absolute
squares are all finite, though not necessarily bounded.  These are the
so-called \emph{square-summable matrices}.

The more detailed theory of these matrices is to be developed in the
work announced in \hyperref[fn:source-22]{Note~22}.  Here we give only what is needed for
orientation and for the purposes of \hyperref[app:IV]{Appendix IV}.

A matrix
\[
  A=\{a_{\mu\nu}\},\qquad \mu,\nu=1,2,\ldots,
\]
will be called \emph{Hermitian square-summable} (abbreviated
H.\ square-summable) if
\[
  a_{\mu\nu}=\overline{a_{\nu\mu}},
  \qquad
  \sum_{\nu=1}^{\infty}|a_{\mu\nu}|^2<\infty
  \quad(\mu=1,2,\ldots).
\]
\footnote{\label{fn:74}Then every series
\(\sum_{\rho=1}^{\infty}a_{\mu\rho}a_{\rho\nu}\) converges absolutely,
so the matrix \(A^2\) can be formed.}
If \(\varphi_1,\varphi_2,\ldots\) is a complete normalized orthogonal
system, define an operator \(R\) on the vectors \(\varphi_\mu\), and
only on those vectors, by
\[
  R\varphi_\mu
  =\sum_{\nu=1}^{\infty}a_{\mu\nu}\varphi_\nu;
\]
the series converges.  The closed linear manifold spanned by
\(\varphi_1,\varphi_2,\ldots\) is \(\mathfrak H\), and
\[
  (\varphi_\mu,R\varphi_\nu)
  =\overline{a_{\nu\mu}}
  =a_{\mu\nu}
  =(R\varphi_\mu,\varphi_\nu).
\]
Thus \(R\) is Hermitian.  We call \(R\) the \emph{elementary Hermitian
operator} of \(A,\{\varphi_1,\varphi_2,\ldots\}\), and call its linear
extension \(\widehat R\), and its closed linear extension
\(\widetilde R\), the \emph{linear} and \emph{closed linear Hermitian
operators}, respectively, of
\(A,\{\varphi_1,\varphi_2,\ldots\}\).

The operator \(\widehat Rf\) is defined precisely for finite sums
\(f=\sum_{\mu=1}^{n}x_\mu\varphi_\mu\), and
\[
 \widehat Rf
 =\sum_{\mu=1}^{n}x_\mu
      \left(\sum_{\nu=1}^{\infty}a_{\mu\nu}\varphi_\nu\right)
 =\sum_{\nu=1}^{\infty}
      \left(\sum_{\mu=1}^{n}a_{\mu\nu}x_\mu\right)\varphi_\nu.
\]

The operator \(\widetilde R\) is not so easy to characterize directly,
but its extension elements are.  For \(f\) to be an extension element,
with \(f^*\) assigned to it, one must have
\[
 (f^*,\varphi_\nu)
 =(f,R\varphi_\nu)
 =\sum_{\mu=1}^{\infty}(f,\varphi_\mu)
       (\varphi_\mu,R\varphi_\nu)
 =\sum_{\mu=1}^{\infty}a_{\mu\nu}(f,\varphi_\mu),
 \quad \nu=1,2,\ldots,
\]
since \(R\) and \(\widetilde R\) have the same extension elements.
Write
\[
  (f,\varphi_\mu)=x_\mu,\qquad
  f=\sum_{\mu=1}^{\infty}x_\mu\varphi_\mu.
\]
There is an \(f^*\) satisfying
\[
  (f^*,\varphi_\nu)
  =\sum_{\mu=1}^{\infty}a_{\mu\nu}x_\mu
\]
if and only if
\[
  \sum_{\nu=1}^{\infty}
  \left|\sum_{\mu=1}^{\infty}a_{\mu\nu}x_\mu\right|^2<\infty,
\]
and in that case it is
\[
  f^*=\sum_{\nu=1}^{\infty}
      \left(\sum_{\mu=1}^{\infty}a_{\mu\nu}x_\mu\right)\varphi_\nu.
\]
Thus
\[
  f=\sum_{\mu=1}^{\infty}x_\mu\varphi_\mu,
  \qquad \sum_{\mu=1}^{\infty}|x_\mu|^2<\infty,
\]
is an extension element if and only if, for
\[
  y_\nu=\sum_{\mu=1}^{\infty}a_{\mu\nu}x_\mu,
  \tag{75}\label{eq:appIII-75}
\]
one has \(\sum_{\nu=1}^{\infty}|y_\nu|^2<\infty\); then
\[
  f^*=\sum_{\nu=1}^{\infty}y_\nu\varphi_\nu.
\]
\footnote{\label{fn:75}Both \(\sum_\mu|a_{\mu\nu}|^2\) and
\(\sum_\mu|x_\mu|^2\) are finite, so the series
\(\sum_\mu a_{\mu\nu}x_\mu\) converge absolutely.}
If, in fact, \(\widehat Rf\) is defined, then \(f\) is certainly such
an extension element.

Furthermore,
\begin{align*}
 (f^*,f)
 &=\sum_{\nu=1}^{\infty}y_\nu\overline{x_\nu}
  =\sum_{\nu=1}^{\infty}
      \left(\sum_{\mu=1}^{\infty}a_{\mu\nu}x_\mu\right)
      \overline{x_\nu},\\
 (f,f^*)
 &=\sum_{\mu=1}^{\infty}x_\mu\overline{y_\mu}
  =\sum_{\mu=1}^{\infty}x_\mu
      \left(\sum_{\nu=1}^{\infty}
        a_{\mu\nu}\overline{x_\nu}\right).
\end{align*}
Thus the \(0\)-class is characterized by the interchangeability of the
two infinite summations.

\subsection*{2.}
\label{appIII:2}
\addcontentsline{toc}{subsection}{Appendix III.2}

We now show that, for every closed linear Hermitian operator \(R\),
there are an H.\ square-summable matrix
\(A=\{a_{\mu\nu}\}\) and a complete normalized orthogonal system
\(\varphi_1,\varphi_2,\ldots\) such that \(R\) is the closed linear
Hermitian operator of \(A,\{\varphi_1,\varphi_2,\ldots\}\).  It is
enough to specify the \(\varphi_\nu\), for
\[
  a_{\mu\nu}=(R\varphi_\mu,\varphi_\nu)
\]
then determines \(A\).  If every \(R\varphi_\mu\) is defined, \(A\) is
H.\ square-summable: \(a_{\mu\nu}=\overline{a_{\nu\mu}}\) is clear, and
\[
  \sum_{\nu=1}^{\infty}|a_{\mu\nu}|^2
  =\sum_{\nu=1}^{\infty}|(R\varphi_\mu,\varphi_\nu)|^2
  =|R\varphi_\mu|^2<\infty.
\]

Let \(S\) be the elementary operator of
\(A,\{\varphi_1,\varphi_2,\ldots\}\).  Then
\[
  S\varphi_\mu
  =\sum_{\nu=1}^{\infty}a_{\mu\nu}\varphi_\nu
  =\sum_{\nu=1}^{\infty}(R\varphi_\mu,\varphi_\nu)\varphi_\nu
  =R\varphi_\mu;
\]
thus \(S\) is the restriction of \(R\) to
\(\varphi_1,\varphi_2,\ldots\).  We must choose this system so that
\(\widetilde S=R\).  In other words, for every \(f\) for which \(Rf\)
is defined, there must be a sequence \(f_1,f_2,\ldots\) of finite
linear combinations of \(\varphi_1,\varphi_2,\ldots\) such that
\[
  f_n\to f,\qquad Rf_n\to Rf,
\]
and of course every \(R\varphi_n\) must be defined.

This is achieved if there is a sequence \(g_1,g_2,\ldots\), with every
\(Rg_n\) defined, such that for each \(f\) in the domain of \(R\) a
subsequence \(g_{N_1},g_{N_2},\ldots\) satisfies
\[
  g_{N_n}\to f,\qquad Rg_{N_n}\to Rf.
\]
This sequence is automatically everywhere dense.  It is then enough
to orthogonalize \(g_1,g_2,\ldots\), as in
\hyperref[thm:8]{Theorem 8}, to obtain
\(\varphi_1,\varphi_2,\ldots\).

Such a sequence is constructed as follows.  Let \(h_1,h_2,\ldots\) be
everywhere dense but otherwise arbitrary.  For every triple of
positive integers \(m,n,k\) for which there is an \(f\) in the domain
of \(R\) satisfying
\[
  |h_m-f|\leq\frac1k,\qquad
  |h_n-Rf|\leq\frac1k,
\]
choose one such \(f\), denoted \(f_{m,n,k}\).  Written as a single
sequence, these \(f_{m,n,k}\) have all the required properties.
Indeed, given \(f,Rf\) and \(\varepsilon>0\), choose \(k\) with
\(2/k\leq\varepsilon\), and choose \(h_m,h_n\) such that
\[
  |h_m-f|\leq\frac1k,\qquad
  |h_n-Rf|\leq\frac1k.
\]
Then \(f_{m,n,k}\) is defined, and
\[
  |h_m-f_{m,n,k}|\leq\frac1k,\qquad
  |h_n-Rf_{m,n,k}|\leq\frac1k.
\]
Consequently
\[
  |f-f_{m,n,k}|\leq\varepsilon,\qquad
  |Rf-Rf_{m,n,k}|\leq\varepsilon.
\]
Since \(\varepsilon\) was arbitrary, the construction is complete.

\subsection*{3.}
\label{appIII:3}
\addcontentsline{toc}{subsection}{Appendix III.3}

For a given \(R\), there may be several pairs
\[
  A=\{a_{\mu\nu}\},\{\varphi_1,\varphi_2,\ldots\},
  \qquad
  B=\{b_{\mu\nu}\},\{\psi_1,\psi_2,\ldots\},
\]
whose closed linear Hermitian operator is \(R\); both \(A\) and \(B\)
are to be H.\ square-summable.  The two complete normalized orthogonal
systems determine a unitary infinite matrix
\[
  \{u_{\mu\nu}\},\qquad u_{\mu\nu}=(\varphi_\mu,\psi_\nu).
\]
Indeed,
\begin{align}
 \sum_{\rho=1}^{\infty}u_{\mu\rho}\overline{u_{\nu\rho}}
 &=\sum_{\rho=1}^{\infty}
   (\varphi_\mu,\psi_\rho)(\psi_\rho,\varphi_\nu)
  =(\varphi_\mu,\varphi_\nu)=\delta_{\mu\nu},\notag\\
 \sum_{\rho=1}^{\infty}\overline{u_{\rho\mu}}u_{\rho\nu}
 &=\sum_{\rho=1}^{\infty}
   (\psi_\mu,\varphi_\rho)(\varphi_\rho,\psi_\nu)
  =(\psi_\mu,\psi_\nu)=\delta_{\mu\nu}.
 \tag{U\(_1\)}\label{eq:U1}
\end{align}
Moreover,
\begin{align}
  \varphi_\mu
  &=\sum_{\nu=1}^{\infty}(\varphi_\mu,\psi_\nu)\psi_\nu
   =\sum_{\nu=1}^{\infty}u_{\mu\nu}\psi_\nu,\notag\\
  \psi_\mu
  &=\sum_{\nu=1}^{\infty}(\psi_\mu,\varphi_\nu)\varphi_\nu
   =\sum_{\nu=1}^{\infty}\overline{u_{\nu\mu}}\varphi_\nu.
  \tag{U\(_2\)}\label{eq:U2}
\end{align}

The usual transformation formulas also hold for
\(a_{\mu\nu},b_{\mu\nu}\):
\begin{align*}
 a_{\mu\nu}
 &=(\varphi_\mu,R\varphi_\nu)\\
 &=\sum_{\rho=1}^{\infty}
   (\varphi_\mu,\psi_\rho)(\psi_\rho,R\varphi_\nu)\\
 &=\sum_{\rho=1}^{\infty}
   (\varphi_\mu,\psi_\rho)(R\psi_\rho,\varphi_\nu)\\
 &=\sum_{\rho=1}^{\infty}
   (\varphi_\mu,\psi_\rho)
   \left[\sum_{\sigma=1}^{\infty}
      (R\psi_\rho,\psi_\sigma)(\psi_\sigma,\varphi_\nu)\right]\\
 &=\sum_{\rho=1}^{\infty}
   \left(\sum_{\sigma=1}^{\infty}
     b_{\rho\sigma}u_{\mu\rho}\overline{u_{\nu\sigma}}\right).
\end{align*}
Interchanging \(\mu,\nu\), and \(\rho,\sigma\), and taking complex
conjugates, using
\(a_{\mu\nu}=\overline{a_{\nu\mu}}\) and
\(b_{\rho\sigma}=\overline{b_{\sigma\rho}}\), gives
\begin{equation}
 a_{\mu\nu}
 =\sum_{\rho=1}^{\infty}
   \left(\sum_{\sigma=1}^{\infty}
     b_{\rho\sigma}u_{\mu\rho}\overline{u_{\nu\sigma}}\right)
 =\sum_{\sigma=1}^{\infty}
   \left(\sum_{\rho=1}^{\infty}
     b_{\rho\sigma}u_{\mu\rho}\overline{u_{\nu\sigma}}\right).
 \tag{U\(_3\)}\label{eq:U3}
\end{equation}
Similarly, the inverse formulas are\translatornote{In the second
iterated sum of \eqref{eq:U4}, the source prints \(b_{\rho\sigma}\).
The inverse transformation requires \(a_{\rho\sigma}\), as used in
both sums here.}
\begin{equation}
 b_{\mu\nu}
 =\sum_{\rho=1}^{\infty}
   \left(\sum_{\sigma=1}^{\infty}
     a_{\rho\sigma}\overline{u_{\rho\mu}}u_{\sigma\nu}\right)
 =\sum_{\sigma=1}^{\infty}
   \left(\sum_{\rho=1}^{\infty}
     a_{\rho\sigma}\overline{u_{\rho\mu}}u_{\sigma\nu}\right).
 \tag{U\(_4\)}\label{eq:U4}
\end{equation}

Thus far the customary theory of unitary transformations of Hermitian
matrices remains valid, but no farther.  Given an H.\ square-summable
matrix \(A=\{a_{\mu\nu}\}\), a complete normalized orthogonal system
\(\varphi_1,\varphi_2,\ldots\), and an arbitrary unitary transformation
matrix \(\{u_{\mu\nu}\}\) satisfying \eqref{eq:U1}, an attempt to apply
\eqref{eq:U2}--\eqref{eq:U4} need not initially produce any convergent
series at all.  Even if all the convergences in \eqref{eq:U3} and
\eqref{eq:U4} occur, the closed linear Hermitian operator of
\(A,\{\varphi_1,\varphi_2,\ldots\}\) may change under the
transformation.  We shall not pursue these relations here; they will
be discussed fully elsewhere (cf.\ \hyperref[fn:source-22]{Note~22}).

\subsection*{4.}
\label{appIII:4}
\addcontentsline{toc}{subsection}{Appendix III.4}

It remains to mention how to determine the closed linear manifolds
\(\mathfrak E,\mathfrak F\) associated with the closed linear Hermitian
operator \(\widetilde R\) of
\(A,\{\varphi_1,\varphi_2,\ldots\}\), where \(R\) is its elementary
Hermitian operator; cf.\ \hyperref[thm:22]{Theorem 22}.
\(\mathfrak E\), respectively \(\mathfrak F\), is the set of all
\((\widetilde R+i1)f\), respectively \((\widetilde R-i1)f\).
Closedness of \(\widetilde R\), together with the properties of
\(\widetilde R\pm i1\) established in \hyperref[thm:22]{Theorem 22}, readily shows that
this is the set of limit points of \((R\pm i1)f\).\footnote{\label{fn:76}If a
sequence \((\widetilde R\pm i1)f\) converges, then so does
\(U^{\pm1}(\widetilde R\pm i1)f=(\widetilde R\mp i1)f\); hence \(f\) and
\(\widetilde Rf\) converge.  The converse is immediate: if \(f\) and
\(\widetilde Rf\) converge, then so does
\((\widetilde R\pm i1)f\).}
This, in turn, is the linear manifold spanned by the
\((R\pm i1)f\), hence by the \((R\pm i1)\varphi_\mu\).
Accordingly, \(\mathfrak E\), respectively \(\mathfrak F\), is the
closed linear manifold spanned by
\[
  (R\pm i1)\varphi_\mu
  =\sum_{\nu=1}^{\infty}a_{\mu\nu}\varphi_\nu
   \pm i\varphi_\mu.
\]

\section*{Appendix IV. The Operator \texorpdfstring{\(\overline R\)}{R-bar}}
\label{app:IV}
\addcontentsline{toc}{section}{Appendix IV. The Operator \texorpdfstring{\(\overline R\)}{R-bar}}

\subsection*{1.}
\label{appIV:1}
\addcontentsline{toc}{subsection}{Appendix IV.1}

The operator \(\overline R\) and its Cayley transform \(\overline U\)
were defined in \hyperref[chap:X]{Chapter X},
\hyperref[thm:37]{Theorem 37}, as follows.  Let
\(\varphi_0,\varphi_1,\varphi_2,\ldots\) be a complete normalized
orthogonal system and put
\[
  \overline U\left(\sum_{\nu=0}^{\infty}x_\nu\varphi_\nu\right)
  =\sum_{\nu=0}^{\infty}x_\nu\varphi_{\nu+1},
  \qquad
  \sum_{\nu=1}^{\infty}|x_\nu|^2<\infty.
\]
Thus \(\overline Rf\) is defined for all
\begin{align*}
 f=\varphi-\overline U\varphi
 &=x_0\varphi_0+(x_1-x_0)\varphi_1
   +(x_2-x_1)\varphi_2+\cdots,
\end{align*}
and then
\[
  \overline Rf=i(\varphi+\overline U\varphi)
  =ix_0\varphi_0+i(x_0+x_1)\varphi_1
   +i(x_1+x_2)\varphi_2+\cdots.
\]
Equivalently,
\[
  \overline R\left(\sum_{\nu=0}^{\infty}y_\nu\varphi_\nu\right)
\]
is defined when
\[
  |y_0|^2+|y_0+y_1|^2+|y_0+y_1+y_2|^2+\cdots
\]
converges; for necessarily
\[
  x_0=y_0,\quad x_1=y_0+y_1,\quad
  x_2=y_0+y_1+y_2,\ldots.
\]
Its value is then
\[
  iy_0\varphi_0+(2iy_0+iy_1)\varphi_1
  +(2iy_0+2iy_1+iy_2)\varphi_2+\cdots.
\]
In this way we obtain a kind of matrix for \(\overline R\):
\[
  a_{\mu\nu}=
  \begin{cases}
    0,&\mu>\nu,\\
    i,&\mu=\nu,\\
    2i,&\mu<\nu.
  \end{cases}
\]
It is not, of course, a matrix in the sense of
\hyperref[app:III]{Appendix III}: it is neither Hermitian nor
square-summable.  This is unsurprising, since all
\(\overline R\varphi_\mu\), \(\mu=0,1,2,\ldots\), are undefined.

To give a genuine matrix for \(\overline R\), we must first find a
complete normalized orthogonal system
\(\psi_1,\psi_2,\ldots\) for which every
\(\overline R\psi_\mu\) is defined.  Put, for \(m,k=0,1,2,\ldots\),
\begin{align*}
 \omega_{m,k}
 &=\varphi_{m2^{k+1}}+\varphi_{m2^{k+1}+1}
   +\cdots+\varphi_{(2m+1)2^k-1}\\
 &\quad-\varphi_{(2m+1)2^k}
   -\varphi_{(2m+1)2^k+1}
   -\cdots-\varphi_{(m+1)2^{k+1}-1}.
\end{align*}
Our criterion shows that
\(\overline R(\sum_{\nu=0}^{\infty}x_\nu\varphi_\nu)\) is certainly
defined if only finitely many \(x_\nu\) are nonzero and
\(\sum_{\nu=0}^{\infty}x_\nu=0\); this is the case here.

The equality
\[
  (\omega_{m,k},\omega_{n,l})=0
\]
certainly holds when the two expressions have no
\(\varphi_\nu\)-term in common.  It also holds when all the
\(\varphi_\nu\) occurring in one occur in the other, all with
coefficient \(+1\), or all with coefficient \(-1\).  Except when
\(m=n,k=l\), one of these alternatives always occurs; on the other
hand,
\[
  |\omega_{m,k}|^2=2^{k+1}.
\]
Consequently
\[
  \psi_{m,k}=2^{-(k+1)/2}\omega_{m,k}
\]
form a normalized orthogonal system, and all
\(\overline R\psi_{m,k}\) are defined.

This system is complete.  Indeed, let
\[
  f=\sum_{\nu=0}^{\infty}x_\nu\varphi_\nu
\]
be orthogonal to all \(\psi_{m,k}\), equivalently to all
\(\omega_{m,k}\).  This says
\begin{align*}
 &x_{m2^{k+1}}+x_{m2^{k+1}+1}+\cdots
  +x_{(2m+1)2^k-1}\\
 &\qquad
 -x_{(2m+1)2^k}-x_{(2m+1)2^k+1}-\cdots
 -x_{(m+1)2^{k+1}-1}=0.
\end{align*}
For \(k=0\), this yields
\[
  x_0-x_1=0,\quad x_2-x_3=0,\quad
  x_4-x_5=0,\quad x_6-x_7=0,\ldots;
\]
that is,
\[
  x_0=x_1,\quad x_2=x_3,\quad
  x_4=x_5,\quad x_6=x_7,\ldots.
\]
For \(k=1\), it gives
\[
  x_0+x_1-x_2-x_3=0,\quad
  x_4+x_5-x_6-x_7=0,\ldots,
\]
so
\[
  2x_0=2x_2,\quad2x_4=2x_6,\ldots,
\]
and therefore
\[
  x_0=x_1=x_2=x_3,\qquad
  x_4=x_5=x_6=x_7,\ldots.
\]
For \(k=2\), it gives
\[
  x_0+x_1+x_2+x_3-x_4-x_5-x_6-x_7=0,
\]
and hence \(4x_0=4x_4\), and so forth.  Thus, altogether,
\[
  x_0=x_1=x_2=\cdots.
\]
But \(\sum_{\nu=0}^{\infty}|x_\nu|^2\) can be finite only if all these
numbers vanish.  Hence \(f=0\), proving completeness.

\subsection*{2.}
\label{appIV:2}
\addcontentsline{toc}{subsection}{Appendix IV.2}

Define the matrix \(\{b_{m,k;n,l}\}\) by
\[
  b_{m,k;n,l}
  =(\psi_{m,k},\overline R\psi_{n,l})
  =2^{-(k+l)/2-1}(\omega_{m,k},\overline R\omega_{n,l}).
\]
It is H.\ square-summable, and its elementary Hermitian operator \(S\)
is the restriction of \(\overline R\) to the \(\psi_{m,k}\).
Consequently \(\overline R\) is an extension of \(S\), and, being
closed and linear, an extension of \(\overline S\).  If \(\overline S\)
is maximal, then necessarily
\[
  \overline R=\overline S;
\]
that is, \(\overline R\) is the closed linear Hermitian operator of
\(\{b_{m,k;n,l}\},\{\psi_{m,k}\}\).

We prove maximality by showing that the manifold \(\mathfrak E\) of
\(\overline S\) is all of \(\mathfrak H\).  By
\hyperref[appIII:4]{Appendix III.4}, the vectors
\(\overline R\psi_{m,k}+i\psi_{m,k}\), or equivalently
\(\overline R\omega_{m,k}+i\omega_{m,k}\), must span
\(\mathfrak H\) as a closed linear manifold.  Thus it is enough to
show that if
\[
  f=\sum_{\nu=0}^{\infty}x_\nu\varphi_\nu
\]
is orthogonal to all \(\overline R\omega_{m,k}+i\omega_{m,k}\), then
\(f=0\).

The formulas in the preceding section give
\begin{align*}
 \overline R\omega_{m,k}
 &=i\varphi_{m2^{k+1}}
   +3i\varphi_{m2^{k+1}+1}
   +\cdots+(2^{k+1}-1)i\varphi_{(2m+1)2^k-1}\\
 &\quad +(2^{k+1}-1)i\varphi_{(2m+1)2^k}
   +(2^{k+1}-3)i\varphi_{(2m+1)2^k+1}
   +\cdots\\
 &\quad +3i\varphi_{(m+1)2^{k+1}-2}
   +i\varphi_{(m+1)2^{k+1}-1},
\end{align*}
and therefore
\begin{align*}
 \overline R\omega_{m,k}+i\omega_{m,k}
 &=2i\varphi_{m2^{k+1}}
   +4i\varphi_{m2^{k+1}+1}
   +\cdots+2^{k+1}i\varphi_{(2m+1)2^k-1}\\
 &\quad +(2^{k+1}-2)i\varphi_{(2m+1)2^k}
   +(2^{k+1}-4)i\varphi_{(2m+1)2^k+1}
   +\cdots\\
 &\quad +2i\varphi_{(m+1)2^{k+1}-2}.
\end{align*}
Orthogonality of \(f\) to these vectors means
\[
  2i\overline{x}_{m2^{k+1}}
  +4i\overline{x}_{m2^{k+1}+1}
  +\cdots
  +2^{k+1}i\overline{x}_{(2m+1)2^k-1}
  +\cdots
  +2i\overline{x}_{(m+1)2^{k+1}-2}=0.
\]
Thus
\[
  x_0=0,\quad x_2=0,\quad x_4=0,\ldots;
\]
\[
  x_0+2x_1+x_2=0,\quad
  x_4+2x_5+x_6=0,\ldots;
\]
\[
  x_0+2x_1+3x_2+4x_3+3x_4+2x_5+x_6=0,\ldots,
\]
and so forth.  It follows successively that
\[
  x_0=x_1=x_2=\cdots=0,
\]
and hence \(f=0\).

It remains to calculate
\[
  b_{m,k;n,l}
  =2^{-(k+l)/2-1}(\omega_{m,k},\overline R\omega_{n,l}),
\]
which is immediate because both vectors have been expressed in terms
of the \(\varphi_\nu\).  If \(m\neq n\) and \(k=l\), there are no common
\(\varphi_\nu\)-terms, so the result is zero; for \(m=n,k=l\), direct
calculation again gives zero.  Let \(k>l\).  If \(n\) lies outside
\[
  m2^{k-l},\ldots,(m+1)2^{k-l}-1,
\]
there are again no common \(\varphi_\nu\)-terms, and the result is
zero.  If \(n\) lies in this interval, then
\[
  n=m2^{k-l}+\rho
  \quad\text{or}\quad
  n=(m+1)2^{k-l}-\rho-1,
  \qquad
  \rho=0,1,\ldots,2^{k-l-1}-1.
\]
In the two cases, respectively,
\begin{align*}
 (\omega_{m,k},\overline R\omega_{n,l})
 &=
 \sum_{\alpha=\rho2^{l+1}}^{(2\rho+1)2^l-1}
      \bigl[-i(2\alpha+1)\bigr]
 -\sum_{\alpha=(2\rho+1)2^l}^{(\rho+1)2^{l+1}-1}
      \bigl[-i(2\alpha+1)\bigr]\\
 &=-i\,2^l(-2\cdot2^l)=i\,2^{2l+1},
\end{align*}
or
\begin{align*}
 (\omega_{m,k},\overline R\omega_{n,l})
 &=
 \sum_{\alpha=(2\rho+1)2^l}^{(\rho+1)2^{l+1}-1}
      \bigl[-i(2\alpha+1)\bigr]
 -\sum_{\alpha=\rho2^{l+1}}^{(2\rho+1)2^l-1}
      \bigl[-i(2\alpha+1)\bigr]\\
 &=-i\,2^l(2\cdot2^l)=-i\,2^{2l+1}.
\end{align*}
For \(k<l\), Hermitian symmetry and interchange of \(m,n\), and
\(k,l\), give the corresponding values
\[
  0,\qquad -i\,2^{2k+1},\qquad i\,2^{2k+1}.
\]
After the normalizing factor in \(b_{m,k;n,l}\), the values
\[
  0,\quad\pm i\,2^{2k+1},\quad\pm i\,2^{2l+1}
\]
become
\[
  0,\quad
  \pm i\,2^{(3k-l)/2},\quad
  \pm i\,2^{(3l-k)/2},
\]
respectively.

We reduce the double index \(m,k=0,1,2,\ldots\) to the single index
\(p=1,2,\ldots\) by
\[
  p=(2m+1)2^k.
\]
The result may then be stated as follows.

\begin{namedstatement}{Theorem}{\(1^{**}\)}
\label{thm:appIV-1}
In the complete normalized orthogonal system
\(\psi_1,\psi_2,\ldots\) described above, there is an H.\
square-summable matrix \(B=\{b_{pq}\}\) such that \(\overline R\) is
the closed linear Hermitian operator of
\(B,\{\psi_1,\psi_2,\ldots\}\).  Its entries
\(b_{pq}\), \(p,q=1,2,\ldots\), are defined as follows.

Partition the sequence of positive integers into cells.  The interval
\[
  m2^{k+1}+1,\ldots,(m+1)2^{k+1}
\]
is a cell; the intervals
\[
  m2^{k+1}+1,\ldots,(2m+1)2^k
\]
and
\[
  (2m+1)2^k+1,\ldots,(m+1)2^{k+1}
\]
are its first and second halves, respectively;
\((2m+1)2^k\) is its center and \(k\) its degree.  Every positive
integer \(p\) is the center of exactly one cell, called the cell of
\(p\).

If neither of the cells of \(p,q\) is a proper part of the other, then
\[
  b_{pq}=0.
\]
If one is a proper part of the other, and \(k,l\) are the degrees of
the smaller and larger cells, respectively, then
\[
  b_{pq}=\pm i\,2^{(3k-l)/2}.
\]
The sign is determined by the following table:
\[
\begin{array}{c|cc}
 &\text{it lies in the first half}
 &\text{it lies in the second half}\\ \hline
\text{\(p\)'s cell is smaller}&-&+\\
\text{\(q\)'s cell is smaller}&+&-
\end{array}
\]
Thus \(B\) is \(i\) times a skew-symmetric matrix.\translatornote{In the
reproduced source diagram, the German labels mean: \emph{Zelle},
``cell''; \emph{Erste Hälfte}, ``first half''; \emph{Mitte}, ``center'';
\emph{Zweite Hälfte}, ``second half''; and \emph{Grad}, ``degree.''
Thus \emph{Grad 0}, \ldots, \emph{Grad 3} denote degrees \(0\) through
\(3\).}
\end{namedstatement}

\begin{center}
  \includegraphics[width=\textwidth]{fig1.png}
\end{center}

\subsection*{3.}
\label{appIV:3}
\addcontentsline{toc}{subsection}{Appendix IV.3}

We obtain a better view of \(\overline R\) by realizing it in function
spaces.  Let \(\Omega\) be the interval \(0,1\), and form its function
space: the set of all Lebesgue-measurable functions \(f(x)\) on
\(0<x<1\) for which
\[
  \int_0^1|f(x)|^2\,dx<\infty.
\]
By \hyperref[app:I]{Appendix I}, this is a Hilbert space isomorphic to
\(\mathfrak H\); denote it by \(\mathfrak H^*\).

The functions
\[
  e^{2n\pi ix},\qquad n=0,\pm1,\pm2,\ldots,
\]
form a complete normalized orthogonal system in \(\mathfrak H^*\).
Let \(\mathfrak I\) be the closed linear manifold spanned by
\[
  e^{2n\pi ix},\qquad n=0,1,2,\ldots.
\]
It is infinite-dimensional and therefore itself a Hilbert space; in
\(\mathfrak I\) we shall realize \(\overline R\).

The operator
\[
  Uf(x)=e^{2\pi ix}f(x)
\]
is isometric, indeed unitary, on \(\mathfrak H^*\), and maps
\(\mathfrak I\) into a part of itself.  Put
\[
  \varphi_\nu(x)=e^{2\nu\pi ix},\qquad \nu=0,1,2,\ldots.
\]
Then \(\varphi_0,\varphi_1,\varphi_2,\ldots\) is a complete normalized
orthogonal system in \(\mathfrak I\), and
\[
  U\varphi_\nu=\varphi_{\nu+1}.
\]
Thus \(U\) agrees on \(\mathfrak I\) with \(\overline U\), and
\(\overline R\) is its Cayley transform.  Hence
\(\overline Rf(x)\), for \(f\in\mathfrak I\), is defined when
\[
  f(x)=\varphi(x)-U\varphi(x)
      =(1-e^{2\pi ix})\varphi(x),
  \qquad \varphi\in\mathfrak I,
\]
and then
\[
  \overline Rf(x)
  =i\bigl(\varphi(x)+U\varphi(x)\bigr)
  =i(1+e^{2\pi ix})\varphi(x).
\]
Equivalently, it is defined when both
\[
  f(x),\qquad \frac{f(x)}{1-e^{2\pi ix}}
\]
belong to \(\mathfrak I\), and in that case
\[
  \overline Rf(x)
  =i\,\frac{1+e^{2\pi ix}}{1-e^{2\pi ix}}\,f(x)
  =-\cot(\pi x)\,f(x).
\]

Let \(f(x)\in\mathfrak I\).  For
\(f(x)/(1-e^{2\pi ix})\) to belong to \(\mathfrak I\), its absolute
square must be integrable and it must be orthogonal to every
\[
  e^{2n\pi ix},\qquad n=-1,-2,\ldots.
\]
This implies
\[
  \int_0^1\cot^2(\pi x)|f(x)|^2\,dx<\infty,
\]
since \(\overline Rf(x)=-\cot(\pi x)f(x)\).  We now show that this
condition conversely implies
\[
  \frac{f(x)}{1-e^{2\pi ix}}\in\mathfrak I
  \qquad (f\in\mathfrak I).
\]
Indeed,
\[
 \left|\frac1{1-e^{2\pi ix}}\right|^2
 =\left|\frac{\cot(\pi x)+i}{2i}\right|^2
 \leq\frac12\cot^2(\pi x)+\frac12.
\]
Thus
\[
  \int_0^1
  \left|\frac{f(x)}{1-e^{2\pi ix}}\right|^2dx<\infty,
\]
and the functions
\[
  \left|\frac{f(x)}{1-\rho e^{2\pi ix}}\right|^2,
  \qquad 0<\rho<1,
\]
are absolutely uniformly integrable.  Consequently, as
\(0<\rho<1\), \(\rho\to1\),
\[
 \int_0^1
 \frac{f(x)}{1-\rho e^{2\pi ix}}e^{-2n\pi ix}\,dx
 \longrightarrow
 \int_0^1
 \frac{f(x)}{1-e^{2\pi ix}}e^{-2n\pi ix}\,dx.
\]
We must show that the right-hand side vanishes for
\(n=-1,-2,\ldots\); it is enough to prove this for the left-hand side,
for every \(0<\rho<1\).  Now
\[
  \frac1{1-\rho e^{2\pi ix}}
  =\sum_{\nu=0}^{\infty}\rho^\nu e^{2\nu\pi ix},
\]
and, for fixed \(\rho\), the series converges uniformly in \(x\).
Therefore
\begin{align*}
 \sum_{\nu=0}^{N}\rho^\nu
   \int_0^1 f(x)e^{-2(n-\nu)\pi ix}\,dx
 &=
 \int_0^1f(x)
    \left(\sum_{\nu=0}^{N}\rho^\nu e^{2\nu\pi ix}\right)
    e^{-2n\pi ix}\,dx\\
 &\longrightarrow
 \int_0^1\frac{f(x)}{1-\rho e^{2\pi ix}}
    e^{-2n\pi ix}\,dx
\end{align*}
as \(N\to\infty\).  Since \(n=-1,-2,\ldots\) and
\(\nu=0,1,2,\ldots\), one has \(n-\nu=-1,-2,\ldots\).  Because
\(f\in\mathfrak I\), the left-hand side is always zero; hence the
right-hand side vanishes as well.

We can now prove the following.

\begin{namedstatement}{Theorem}{\(2^{**}\)}
\label{thm:appIV-2}
Let \(\mathfrak H^*\) be the Hilbert function space just described, and
let \(\mathfrak I_k^+\), respectively \(\mathfrak I_k^-\), be the
closed linear manifolds spanned by all
\[
  e^{2n\pi ix}\quad\text{with }n\geq k,
  \qquad\text{respectively with }n\leq k,
\]
where \(n,k=0,\pm1,\pm2,\ldots\).  The
\(\mathfrak I_k^\pm\) are infinite-dimensional, and hence are again
Hilbert spaces.

Define \(R\) on \(\mathfrak H^*\) as follows: if
\(f(x)\) and \(-\cot(\pi x)f(x)\) belong to \(\mathfrak H^*\), that is,
if
\[
  \int_0^1|f(x)|^2\,dx<\infty,\qquad
  \int_0^1\cot^2(\pi x)|f(x)|^2\,dx<\infty,
\]
then \(Rf(x)\) is defined and
\[
  Rf(x)=-\cot(\pi x)f(x).
\]
If \(f(x)\) lies in one of the \(\mathfrak I_k^\pm\) and \(Rf(x)\) is
defined, then \(Rf(x)\) lies in the same
\(\mathfrak I_k^\pm\).  Thus \(R\) may be restricted on each
\(\mathfrak I_k^\pm\) to a Hermitian operator there.\footnote{\label{fn:77}This is
by no means the reducibility of \(R\) by
\(\mathfrak I_k^\pm\): whenever \(Rf\) is defined,
\(RP_{\mathfrak I_k^\pm}f\) need not also be defined.}
On every \(\mathfrak I_k^+\) this gives \(\overline R\), and on every
\(\mathfrak I_k^-\) it gives \(-\overline R\).
\end{namedstatement}

\begin{proof}
For \(\mathfrak I_0^+=\mathfrak I\), this has already been proved.
The unitary map
\[
  f(x)\longmapsto e^{2k\pi ix}f(x)
\]
takes
\((\mathfrak H^*,R,\mathfrak I_0^+)\) to
\((\mathfrak H^*,R,\mathfrak I_k^+)\).  The unitary map
\[
  f(x)\longmapsto f(1-x)
\]
takes \((\mathfrak H^*,R,\mathfrak I_{-k}^+)\) to
\((\mathfrak H^*,-R,\mathfrak I_k^-)\).  Hence the assertions hold for
all \(\mathfrak I_k^\pm\).
\end{proof}

It is also easy to show that this \(R\) is a hypermaximal Hermitian
operator on \(\mathfrak H^*\).

\subsection*{4.}
\label{appIV:4}
\addcontentsline{toc}{subsection}{Appendix IV.4}

To every
\[
  f(x)=\sum_{n=0}^{\infty}a_ne^{2n\pi ix},
  \qquad
  \sum_{n=0}^{\infty}|a_n|^2<\infty,
\]
in \(\mathfrak I_0^+\),\footnote{\label{fn:78}The series
\(\sum_{n=0}^{\infty}a_ne^{2n\pi ix}\) converges, in the metric used
throughout this paper, ``in the mean'':
\[
  \int_0^1\left|
  f(x)-\sum_{n=0}^{N}a_ne^{2n\pi ix}
  \right|^2dx\longrightarrow0
  \qquad(N\to\infty),
\]
but not necessarily pointwise.}
we may associate the function
\[
  F(z)=\sum_{n=0}^{\infty}a_nz^n
\]
analytic in the unit disk; it is analytic there because
\(\sum|a_n|^2<\infty\), and hence \(a_n\to0\).  The circle Parseval
identity gives\translatornote{The source prints \(r^n\) and omits the
normalizing factor \(1/(2\pi i)\) in the contour integral.  Both have
been corrected to the standard Parseval identity.}
\[
 \sum_{n=0}^{\infty}|a_n|^2
 =\lim_{r\to1}\sum_{n=0}^{\infty}r^{2n}|a_n|^2
 =\lim_{r\to1}\frac{1}{2\pi i}
   \int_{|z|=r}|F(z)|^2\,\frac{dz}{z}.
\]
Thus the functions \(F(z)\) analytic in the unit disk for which
\[
  \frac{1}{2\pi i}\int_{|z|=r}|F(z)|^2\,\frac{dz}{z}
\]
remains bounded as \(r\to1\) form a Hilbert space.  Since
\[
  -\cot(\pi x)
  =i\,\frac{e^{2\pi ix}+1}{e^{2\pi ix}-1},
\]
\(\overline R\) is represented there by
\[
  \overline RF(z)=i\,\frac{z+1}{z-1}F(z),
\]
provided the result again belongs to the stated class.

Another noteworthy realization of \(\overline R\) is obtained by
considering functions \(f(x)\) defined on \(0<x<\infty\) with
\[
  \int_0^\infty|f(x)|^2\,dx<\infty.
\]
It then corresponds to the differential operator
\[
  i\,\frac{d}{dx}
\]
with the boundary condition \(f(0)=0\).  We shall not pursue this
realization here.\footnote{\label{fn:79}It will be discussed in the work mentioned
in \hyperref[fn:source-22]{Note~22}.}

It remains to verify the properties of \(\overline R\) used in
\hyperref[thm:40]{Theorem 40}.  First,
\(\overline R\) is not everywhere defined and does not assume every
value.  Realize it, by \hyperref[thm:appIV-2]{Theorem \(2^{**}\)}, in
\(\mathfrak I_0^+\).  Since
\[
  \int_0^1\cot^2(\pi x)\,dx
  =\int_0^1\tan^2(\pi x)\,dx
  =\infty,
\]
\(\overline R\) is undefined for \(f(x)=1\), and never assumes that
value.

Second, \(\overline R\) is bounded neither below nor above.  Realize it
as in \hyperref[thm:appIV-1]{Theorem \(1^{**}\)} and put
\[
  f=\psi_{2^\nu}\pm i\psi_{2^{\nu+1}}.
\]
Then
\[
  (f,\overline Rf)=\mp2^{\nu+\frac12},
  \qquad |f|^2=2.
\]
Thus, for any \(C\), values with
\[
  (f,\overline Rf)>C|f|^2
  \qquad\text{or}\qquad
  (f,\overline Rf)<C|f|^2
\]
are obtained by choosing \(2^{\nu-\frac12}>|C|\).


\end{document}

